% main.tex — Master compilation file
% Paper: "Tokenized Housing and Lifecycle Portfolio Choice:
%         A Decoupling of Location from Housing Exposure"
%
% Target journal: Review of Financial Studies (RFS); backup RAPS / REE.
% Status: skeleton sections complete (s1–s6); [PLACEHOLDER] cells await server1 results.
%
% Compile: pdflatex main.tex  (or lualatex / xelatex)
%          bibtex main
%          pdflatex main.tex  (twice for cross-references)
%
% Or with latexmk:
%   latexmk -pdf main.tex

\documentclass[12pt]{article}

% ─── Packages ───────────────────────────────────────────────────────────────
\usepackage[margin=1.2in]{geometry}
\usepackage{amsmath,amssymb,amsthm}
\usepackage{natbib}
\usepackage{hyperref}
\usepackage{booktabs}
\usepackage{graphicx}
\usepackage{xcolor}
\usepackage{setspace}
\usepackage{microtype}
\usepackage{caption}
\usepackage{subcaption}
\usepackage{appendix}
\usepackage{textcomp}   % \textsubscript

% ─── Layout ─────────────────────────────────────────────────────────────────
\onehalfspacing
\setlength{\parskip}{0pt}
\setlength{\parindent}{1.5em}

% ─── Theorem environments ────────────────────────────────────────────────────
\newtheorem{proposition}{Proposition}
\newtheorem{lemma}{Lemma}
\newtheorem{corollary}{Corollary}
\theoremstyle{definition}
\newtheorem{assumption}{Assumption}
\newtheorem{definition}{Definition}
\theoremstyle{remark}
\newtheorem{remark}{Remark}

% ─── Placeholder command ─────────────────────────────────────────────────────
% Usage: \ph{description}  — red placeholder for numerical results pending server1
\newcommand{\ph}[1]{\textcolor{red}{\textsc{[#1]}}}

% ─── Custom commands ─────────────────────────────────────────────────────────
\newcommand{\cev}{\mathrm{CEV}}
\newcommand{\xprev}{x_{\mathrm{prev}}}
\newcommand{\xnew}{x_{\mathrm{new}}}
\newcommand{\tausell}{\tau_{\mathrm{sell}}}
\newcommand{\taubuy}{\tau_{\mathrm{buy}}}
\newcommand{\tautoken}{\tau_{\mathrm{token}}}
\newcommand{\preloc}{p_{\mathrm{reloc}}}
\newcommand{\rhoAB}{\rho_{AB}}
\newcommand{\deltaown}{\delta_{\mathrm{own}}}
\newcommand{\ellA}{\ell_A}
\newcommand{\ellB}{\ell_B}
\newcommand{\EsubzeroTwoL}{\text{E0}}
\newcommand{\EonesubTwoL}{\text{E1}_{\text{2L}}}
\newcommand{\EtwosubTwoL}{\text{E2}_{\text{2L}}}

% ─── Bibliography style ───────────────────────────────────────────────────────
% RFS uses author-year citations.
\bibliographystyle{plainnat}

% ─────────────────────────────────────────────────────────────────────────────
\begin{document}
% ─────────────────────────────────────────────────────────────────────────────

\title{Tokenized Housing and Lifecycle Portfolio Choice:\\
       A Decoupling of Location from Housing Exposure%
       \thanks{We thank seminar participants for helpful comments.
               All errors are our own.
               Code and replication files at \url{https://github.com/[repo]}.}}

\author{[Authors]}

\date{May 2026\\\textit{Preliminary and incomplete — do not cite or circulate.}}

\maketitle

% ─── Abstract ────────────────────────────────────────────────────────────────
\begin{abstract}
We quantify the lifetime welfare value of residential housing tokenization
using a calibrated two-location lifecycle portfolio model with stochastic
relocation shocks.
Households choose among three regimes: rent-only, traditional binary
homeownership (with 8.5\% round-trip transaction cost at each relocation),
and continuous fractional token ownership portable across locations.
The model tracks prior-period token holdings as explicit state variables,
which creates a forward-looking motive to pre-accumulate tokens of a future
destination before relocating.
Our headline result is $\cev(\EtwosubTwoL\ \text{vs}\ \EonesubTwoL)
= \ph{X.XX}\%$ at baseline calibration
($p_{\text{reloc}} = 6\%$, $\taubuy = 2.5\%$, $\rhoAB = 0.50$).
The total gain decomposes into an avoided-transaction-cost channel (\ph{X}\%)
from escaping the forced sell-and-buy round-trip, and a pre-buying hedge
channel (\ph{X}\%) that is uniquely activated by token portability: by
pre-accumulating tokens of the destination location, households reduce the
marginal buying cost on future relocation.
Falsification tests confirm the mechanism: the hedge channel collapses when
the relocation probability is zero and attenuates as cross-location
correlations approach one.
This constitutes a structurally novel welfare channel outside of
\citet{Liu2021}, \citet{YaoZhang2005}, and \citet{KMW2018}.
\end{abstract}

\medskip
\noindent\textbf{Keywords:} tokenized housing; lifecycle portfolio choice;
relocation; transaction costs; household finance; blockchain.\\[2pt]
\noindent\textbf{JEL codes:} G11; G51; R21; R31.

\clearpage

% ─────────────────────────────────────────────────────────────────────────────
% s1_intro.tex — Introduction
% Paper: "Tokenized Housing and Lifecycle Portfolio Choice:
%         A Decoupling of Location from Housing Exposure"
%
% Source: paper/outline_v4.md §1 (sections 1.1–1.4)
%         docs/welfare_decomp_v4.md §5–6 (contribution framing)
%         docs/calibration_v3.md (empirical anchors)
%
% Status: skeleton complete; [PLACEHOLDER] marks require server1 numerical results.
%         Fill in before submission: CEV headline, channel magnitudes, H1 evidence.
%
% Compilation: % s1_intro.tex — Introduction
% Paper: "Tokenized Housing and Lifecycle Portfolio Choice:
%         A Decoupling of Location from Housing Exposure"
%
% Source: paper/outline_v4.md §1 (sections 1.1–1.4)
%         docs/welfare_decomp_v4.md §5–6 (contribution framing)
%         docs/calibration_v3.md (empirical anchors)
%
% Status: skeleton complete; [PLACEHOLDER] marks require server1 numerical results.
%         Fill in before submission: CEV headline, channel magnitudes, H1 evidence.
%
% Compilation: % s1_intro.tex — Introduction
% Paper: "Tokenized Housing and Lifecycle Portfolio Choice:
%         A Decoupling of Location from Housing Exposure"
%
% Source: paper/outline_v4.md §1 (sections 1.1–1.4)
%         docs/welfare_decomp_v4.md §5–6 (contribution framing)
%         docs/calibration_v3.md (empirical anchors)
%
% Status: skeleton complete; [PLACEHOLDER] marks require server1 numerical results.
%         Fill in before submission: CEV headline, channel magnitudes, H1 evidence.
%
% Compilation: \input{sections/s1_intro} from main.tex
%              Requires amsmath, hyperref, natbib (or biblatex).

% ─────────────────────────────────────────────────────────────────────────────
\section{Introduction}
\label{sec:intro}
% ─────────────────────────────────────────────────────────────────────────────

American households relocate frequently.
The Panel Study of Income Dynamics (PSID) records inter-MSA mobility rates of
5--7\% per year among working-age adults, implying that more than 70\% of
households move at least once between ages~25 and~65.
Each relocation under traditional homeownership forces a \emph{sell-and-buy
round-trip}: the household sells the current residence, incurring a broker
commission and closing costs of roughly 6\% of home value, then purchases at
the new location, adding origination fees, title insurance, and inspection
costs of a further 2--3\%.
The aggregate round-trip transaction cost is approximately 8--10\% of home
value per relocation event \citep{NAR2023, CFPB2022}.

These frictions impose a compounding welfare loss over the lifecycle.
A household that relocates four times between ages~25 and~65~---~a plausible
modal count under PSID mobility rates---~faces cumulative round-trip costs
of 32--40\% of a home's value, equivalent to multiple years of mortgage
payments forgone.
Beyond the direct cost, each forced sale severs the household's
\emph{location-specific housing-market exposure}: upon arriving at
location~$B$, the former location-$A$ owner no longer participates in
location~$A$'s price appreciation, even if labor, family, or human-capital
ties make future return migration likely.
Diversified real estate investment trusts (REITs) aggregate at the
portfolio level and provide no location-specific hedge.

Fractional residential housing tokens---blockchain-registered claims to
location-specific property cash flows, exemplified by RealT-class
platforms---offer a structural remedy.
A tokenized holding at location~$A$ can be \emph{retained across
relocation} to location~$B$ at a transfer cost of approximately 1\%,
versus the 8.5\% round-trip for direct ownership.
More importantly, a forward-looking household at location~$A$ who
anticipates a possible future relocation to~$B$ can \emph{pre-accumulate}
fractional location-$B$ tokens incrementally, paying the 2.5\% buying cost
in small installments rather than as a lump-sum at the moment of arrival.
This pre-accumulation motive is the structurally novel mechanism that
distinguishes token portability from both direct ownership and diversified REITs:
the household \emph{decouples its geographic location from its housing-market exposure}.

This paper provides the first lifecycle welfare quantification of
location-decoupling.
We build a two-location lifecycle portfolio model with stochastic relocation
shocks, location-correlated house-price processes, and three ownership regimes:
rent-only (E0), traditional binary homeownership that forces a full sell-and-buy
at every relocation (E1\_2L), and continuous fractional token ownership that
is portable across locations and priced correctly via a 6D state space that
tracks prior-period token holdings (E2\_2L).
The 6D state---which includes last period's token positions $(x_{A,\text{prev}},
x_{B,\text{prev}})$ as explicit state variables---enables the model to price
the per-period buying cost correctly and to generate the forward-looking
pre-accumulation motive that is the paper's central mechanism.

Our headline result is:
\begin{equation}
    \operatorname{CEV}(E2\_2L \text{ vs } E1\_2L) = \textsc{[x.xx\%]},
    \label{eq:headline_cev}
\end{equation}
at baseline calibration ($p_{\text{reloc}} = 0.06$, $\tau_{\text{buy}} = 2.5\%$,
$\rho_{AB} = 0.50$).
[\textsc{Placeholder: fill in once server1 runs confirm.}]
This gain decomposes into three channels:
\begin{enumerate}
    \item \emph{Avoided-transaction-cost channel}: \textsc{[x.x\%]}.
        At each relocation event, the E1\_2L household pays
        $\tau_{\text{sell}} + \tau_{\text{buy}} \approx 8.5\%$;
        the E2\_2L household retains tokens at $\tau_{\text{token}} \approx 1\%$.
        This channel is unique to token portability and has no analog in
        \citet{Liu2021}, who abstracts from mobility.

    \item \emph{Maintained-hedge channel}: \textsc{[x.x\%]}.
        A household at location~$B$ who holds $x_{A,\text{prev}} > 0$ continues
        to participate in location~$A$'s price appreciation.
        The mechanism activates through the pre-accumulation motive enabled by
        the 6D state: the household optimally builds $x_B > 0$ at location~$A$
        to reduce future buying costs on relocation (the ``pre-buy hedge'').
        [\textsc{Placeholder: confirm sign via $\bar{x}_B > 0$ at $\ell = A$
        in server1 output (H1 test).}]

    \item \emph{Continuous-ownership channel}: \textsc{[x.x\%]}.
        Continuous fractional ownership of the occupied residence provides
        rent saving of $\delta_{\text{own}} = \rho - m = 4\%$ per unit held.
        This channel is common to \citet{Liu2021} and replicates his
        MHS-relaxation result within our model.
\end{enumerate}
Channels~(1) and~(2) are the paper's structural contribution beyond
\citet{Liu2021}: token portability uniquely provides (1) because direct
owners face forced sales, and REITs provide no location-specific claims
for channel~(2).
The cross-term between channels is empirically negligible (\textsc{[x.xx\%]})
as in our v3 channel-decomposition benchmark, confirming additive separability.

\paragraph{Falsification structure.}
We pre-register three mechanism tests that must pass for the token-portability
claim to be credible.
(r)~Setting $p_{\text{reloc}} = 0$ shuts off the mobility channel;
    channels~(1) and~(2) must collapse to zero.
(m)~Setting $\rho_{AB} \to 1$ eliminates diversification benefit;
    the maintained-hedge channel must attenuate monotonically.
(q)~Setting $\tau_{\text{buy}} = 0$ eliminates the pre-accumulation motive;
    the pre-buy hedge channel must vanish.
[\textsc{Placeholder: report $p$-values or CEV changes from server1 sweep results.}]

\paragraph{Related literature.}
Our paper connects four strands.

\medskip\noindent
\textbf{Lifecycle housing-portfolio choice.}
\citet{YaoZhang2005} and \citet{Cocco2005} establish the canonical
binary-tenure lifecycle model; we extend by adding a second location,
stochastic relocation, and token portability.
\citet{KraftMunk2011} and \citet{KMW2018} introduce habit and
adjustment-cost extensions; we add geographic mobility.
\citet{Liu2021} introduces the Mortgage-with-Housing-Supply (MHS)
relaxation: continuous ownership fraction within a single location.
Our continuous-$x$ channel (E2\_2L at a single location) replicates Liu's
mechanism; our paper's novel contribution lies in channels~(1) and~(2)
above, which are outside Liu's single-location framework.

\medskip\noindent
\textbf{Housing as hedge against mobility risk.}
\citet{SinaiSouleles2005} show that homeownership hedges against rent-price
risk for households who expect to remain in a location.
\citet{Davidoff2006} documents income-housing correlation as a source of
labor-market risk for owner-occupiers.
\citet{BFN2014} (Bagliano, Fugazza, and Nicodano) show that housing and
equity exposure interact through labor income in a lifecycle model.
Our contribution is distinct: we model the \emph{loss} of location-specific
hedge at relocation and quantify the welfare value of maintaining it via tokens.

\medskip\noindent
\textbf{Transaction costs in housing markets.}
\citet{FlavinYamashita2002} document that transaction costs constrain
portfolio rebalancing in a housing lifecycle context.
\citet{Han2013} quantifies the lock-in effect of transaction costs on
tenure choice.
\citet{PiazzesiSchneider2016} embed transaction costs in an equilibrium
housing-search model.
Our model treats transaction costs as the \emph{source} of the E1\_2L
welfare loss; token portability is valued precisely because it reduces these costs.

\medskip\noindent
\textbf{Tokenization and digital assets.}
\citet{CongLiWang2021} provide the foundational theory of tokenomics in a
general equilibrium setting.
\citet{Swinkels2023} surveys the institutional structure of residential
real estate tokenization platforms.
Our contribution is a lifecycle welfare quantification \emph{within} a
standard Cocco-class model, grounded in the PSID-/NAR-calibrated parameter
space, rather than a platform mechanism design.

\paragraph{Partial equilibrium caveat.}
All results are partial-equilibrium: house prices are exogenous.
We interpret $\operatorname{CEV}(E2\_2L \text{ vs } E1\_2L)$ as the
household-level welfare value of \emph{access} to tokenized instruments,
not as an aggregate welfare estimate.
A general-equilibrium extension with endogenous prices and adoption externalities
is left for future work.

\paragraph{Road map.}
Section~\ref{sec:model} presents the two-location lifecycle model.
Section~\ref{sec:calibration} describes the calibration.
Section~\ref{sec:results} reports baseline welfare results and the channel
decomposition.
Section~\ref{sec:discussion} discusses the contribution relative to
\citet{Liu2021} and the robustness of the falsification structure.
Section~\ref{sec:conclusion} concludes.
All numerical details are in the appendices.
 from main.tex
%              Requires amsmath, hyperref, natbib (or biblatex).

% ─────────────────────────────────────────────────────────────────────────────
\section{Introduction}
\label{sec:intro}
% ─────────────────────────────────────────────────────────────────────────────

American households relocate frequently.
The Panel Study of Income Dynamics (PSID) records inter-MSA mobility rates of
5--7\% per year among working-age adults, implying that more than 70\% of
households move at least once between ages~25 and~65.
Each relocation under traditional homeownership forces a \emph{sell-and-buy
round-trip}: the household sells the current residence, incurring a broker
commission and closing costs of roughly 6\% of home value, then purchases at
the new location, adding origination fees, title insurance, and inspection
costs of a further 2--3\%.
The aggregate round-trip transaction cost is approximately 8--10\% of home
value per relocation event \citep{NAR2023, CFPB2022}.

These frictions impose a compounding welfare loss over the lifecycle.
A household that relocates four times between ages~25 and~65~---~a plausible
modal count under PSID mobility rates---~faces cumulative round-trip costs
of 32--40\% of a home's value, equivalent to multiple years of mortgage
payments forgone.
Beyond the direct cost, each forced sale severs the household's
\emph{location-specific housing-market exposure}: upon arriving at
location~$B$, the former location-$A$ owner no longer participates in
location~$A$'s price appreciation, even if labor, family, or human-capital
ties make future return migration likely.
Diversified real estate investment trusts (REITs) aggregate at the
portfolio level and provide no location-specific hedge.

Fractional residential housing tokens---blockchain-registered claims to
location-specific property cash flows, exemplified by RealT-class
platforms---offer a structural remedy.
A tokenized holding at location~$A$ can be \emph{retained across
relocation} to location~$B$ at a transfer cost of approximately 1\%,
versus the 8.5\% round-trip for direct ownership.
More importantly, a forward-looking household at location~$A$ who
anticipates a possible future relocation to~$B$ can \emph{pre-accumulate}
fractional location-$B$ tokens incrementally, paying the 2.5\% buying cost
in small installments rather than as a lump-sum at the moment of arrival.
This pre-accumulation motive is the structurally novel mechanism that
distinguishes token portability from both direct ownership and diversified REITs:
the household \emph{decouples its geographic location from its housing-market exposure}.

This paper provides the first lifecycle welfare quantification of
location-decoupling.
We build a two-location lifecycle portfolio model with stochastic relocation
shocks, location-correlated house-price processes, and three ownership regimes:
rent-only (E0), traditional binary homeownership that forces a full sell-and-buy
at every relocation (E1\_2L), and continuous fractional token ownership that
is portable across locations and priced correctly via a 6D state space that
tracks prior-period token holdings (E2\_2L).
The 6D state---which includes last period's token positions $(x_{A,\text{prev}},
x_{B,\text{prev}})$ as explicit state variables---enables the model to price
the per-period buying cost correctly and to generate the forward-looking
pre-accumulation motive that is the paper's central mechanism.

Our headline result is:
\begin{equation}
    \operatorname{CEV}(E2\_2L \text{ vs } E1\_2L) = \textsc{[x.xx\%]},
    \label{eq:headline_cev}
\end{equation}
at baseline calibration ($p_{\text{reloc}} = 0.06$, $\tau_{\text{buy}} = 2.5\%$,
$\rho_{AB} = 0.50$).
[\textsc{Placeholder: fill in once server1 runs confirm.}]
This gain decomposes into three channels:
\begin{enumerate}
    \item \emph{Avoided-transaction-cost channel}: \textsc{[x.x\%]}.
        At each relocation event, the E1\_2L household pays
        $\tau_{\text{sell}} + \tau_{\text{buy}} \approx 8.5\%$;
        the E2\_2L household retains tokens at $\tau_{\text{token}} \approx 1\%$.
        This channel is unique to token portability and has no analog in
        \citet{Liu2021}, who abstracts from mobility.

    \item \emph{Maintained-hedge channel}: \textsc{[x.x\%]}.
        A household at location~$B$ who holds $x_{A,\text{prev}} > 0$ continues
        to participate in location~$A$'s price appreciation.
        The mechanism activates through the pre-accumulation motive enabled by
        the 6D state: the household optimally builds $x_B > 0$ at location~$A$
        to reduce future buying costs on relocation (the ``pre-buy hedge'').
        [\textsc{Placeholder: confirm sign via $\bar{x}_B > 0$ at $\ell = A$
        in server1 output (H1 test).}]

    \item \emph{Continuous-ownership channel}: \textsc{[x.x\%]}.
        Continuous fractional ownership of the occupied residence provides
        rent saving of $\delta_{\text{own}} = \rho - m = 4\%$ per unit held.
        This channel is common to \citet{Liu2021} and replicates his
        MHS-relaxation result within our model.
\end{enumerate}
Channels~(1) and~(2) are the paper's structural contribution beyond
\citet{Liu2021}: token portability uniquely provides (1) because direct
owners face forced sales, and REITs provide no location-specific claims
for channel~(2).
The cross-term between channels is empirically negligible (\textsc{[x.xx\%]})
as in our v3 channel-decomposition benchmark, confirming additive separability.

\paragraph{Falsification structure.}
We pre-register three mechanism tests that must pass for the token-portability
claim to be credible.
(r)~Setting $p_{\text{reloc}} = 0$ shuts off the mobility channel;
    channels~(1) and~(2) must collapse to zero.
(m)~Setting $\rho_{AB} \to 1$ eliminates diversification benefit;
    the maintained-hedge channel must attenuate monotonically.
(q)~Setting $\tau_{\text{buy}} = 0$ eliminates the pre-accumulation motive;
    the pre-buy hedge channel must vanish.
[\textsc{Placeholder: report $p$-values or CEV changes from server1 sweep results.}]

\paragraph{Related literature.}
Our paper connects four strands.

\medskip\noindent
\textbf{Lifecycle housing-portfolio choice.}
\citet{YaoZhang2005} and \citet{Cocco2005} establish the canonical
binary-tenure lifecycle model; we extend by adding a second location,
stochastic relocation, and token portability.
\citet{KraftMunk2011} and \citet{KMW2018} introduce habit and
adjustment-cost extensions; we add geographic mobility.
\citet{Liu2021} introduces the Mortgage-with-Housing-Supply (MHS)
relaxation: continuous ownership fraction within a single location.
Our continuous-$x$ channel (E2\_2L at a single location) replicates Liu's
mechanism; our paper's novel contribution lies in channels~(1) and~(2)
above, which are outside Liu's single-location framework.

\medskip\noindent
\textbf{Housing as hedge against mobility risk.}
\citet{SinaiSouleles2005} show that homeownership hedges against rent-price
risk for households who expect to remain in a location.
\citet{Davidoff2006} documents income-housing correlation as a source of
labor-market risk for owner-occupiers.
\citet{BFN2014} (Bagliano, Fugazza, and Nicodano) show that housing and
equity exposure interact through labor income in a lifecycle model.
Our contribution is distinct: we model the \emph{loss} of location-specific
hedge at relocation and quantify the welfare value of maintaining it via tokens.

\medskip\noindent
\textbf{Transaction costs in housing markets.}
\citet{FlavinYamashita2002} document that transaction costs constrain
portfolio rebalancing in a housing lifecycle context.
\citet{Han2013} quantifies the lock-in effect of transaction costs on
tenure choice.
\citet{PiazzesiSchneider2016} embed transaction costs in an equilibrium
housing-search model.
Our model treats transaction costs as the \emph{source} of the E1\_2L
welfare loss; token portability is valued precisely because it reduces these costs.

\medskip\noindent
\textbf{Tokenization and digital assets.}
\citet{CongLiWang2021} provide the foundational theory of tokenomics in a
general equilibrium setting.
\citet{Swinkels2023} surveys the institutional structure of residential
real estate tokenization platforms.
Our contribution is a lifecycle welfare quantification \emph{within} a
standard Cocco-class model, grounded in the PSID-/NAR-calibrated parameter
space, rather than a platform mechanism design.

\paragraph{Partial equilibrium caveat.}
All results are partial-equilibrium: house prices are exogenous.
We interpret $\operatorname{CEV}(E2\_2L \text{ vs } E1\_2L)$ as the
household-level welfare value of \emph{access} to tokenized instruments,
not as an aggregate welfare estimate.
A general-equilibrium extension with endogenous prices and adoption externalities
is left for future work.

\paragraph{Road map.}
Section~\ref{sec:model} presents the two-location lifecycle model.
Section~\ref{sec:calibration} describes the calibration.
Section~\ref{sec:results} reports baseline welfare results and the channel
decomposition.
Section~\ref{sec:discussion} discusses the contribution relative to
\citet{Liu2021} and the robustness of the falsification structure.
Section~\ref{sec:conclusion} concludes.
All numerical details are in the appendices.
 from main.tex
%              Requires amsmath, hyperref, natbib (or biblatex).

% ─────────────────────────────────────────────────────────────────────────────
\section{Introduction}
\label{sec:intro}
% ─────────────────────────────────────────────────────────────────────────────

American households relocate frequently.
The Panel Study of Income Dynamics (PSID) records inter-MSA mobility rates of
5--7\% per year among working-age adults, implying that more than 70\% of
households move at least once between ages~25 and~65.
Each relocation under traditional homeownership forces a \emph{sell-and-buy
round-trip}: the household sells the current residence, incurring a broker
commission and closing costs of roughly 6\% of home value, then purchases at
the new location, adding origination fees, title insurance, and inspection
costs of a further 2--3\%.
The aggregate round-trip transaction cost is approximately 8--10\% of home
value per relocation event \citep{NAR2023, CFPB2022}.

These frictions impose a compounding welfare loss over the lifecycle.
A household that relocates four times between ages~25 and~65~---~a plausible
modal count under PSID mobility rates---~faces cumulative round-trip costs
of 32--40\% of a home's value, equivalent to multiple years of mortgage
payments forgone.
Beyond the direct cost, each forced sale severs the household's
\emph{location-specific housing-market exposure}: upon arriving at
location~$B$, the former location-$A$ owner no longer participates in
location~$A$'s price appreciation, even if labor, family, or human-capital
ties make future return migration likely.
Diversified real estate investment trusts (REITs) aggregate at the
portfolio level and provide no location-specific hedge.

Fractional residential housing tokens---blockchain-registered claims to
location-specific property cash flows, exemplified by RealT-class
platforms---offer a structural remedy.
A tokenized holding at location~$A$ can be \emph{retained across
relocation} to location~$B$ at a transfer cost of approximately 1\%,
versus the 8.5\% round-trip for direct ownership.
More importantly, a forward-looking household at location~$A$ who
anticipates a possible future relocation to~$B$ can \emph{pre-accumulate}
fractional location-$B$ tokens incrementally, paying the 2.5\% buying cost
in small installments rather than as a lump-sum at the moment of arrival.
This pre-accumulation motive is the structurally novel mechanism that
distinguishes token portability from both direct ownership and diversified REITs:
the household \emph{decouples its geographic location from its housing-market exposure}.

This paper provides the first lifecycle welfare quantification of
location-decoupling.
We build a two-location lifecycle portfolio model with stochastic relocation
shocks, location-correlated house-price processes, and three ownership regimes:
rent-only (E0), traditional binary homeownership that forces a full sell-and-buy
at every relocation (E1\_2L), and continuous fractional token ownership that
is portable across locations and priced correctly via a 6D state space that
tracks prior-period token holdings (E2\_2L).
The 6D state---which includes last period's token positions $(x_{A,\text{prev}},
x_{B,\text{prev}})$ as explicit state variables---enables the model to price
the per-period buying cost correctly and to generate the forward-looking
pre-accumulation motive that is the paper's central mechanism.

Our headline result is:
\begin{equation}
    \operatorname{CEV}(E2\_2L \text{ vs } E1\_2L) = \textsc{[x.xx\%]},
    \label{eq:headline_cev}
\end{equation}
at baseline calibration ($p_{\text{reloc}} = 0.06$, $\tau_{\text{buy}} = 2.5\%$,
$\rho_{AB} = 0.50$).
[\textsc{Placeholder: fill in once server1 runs confirm.}]
This gain decomposes into three channels:
\begin{enumerate}
    \item \emph{Avoided-transaction-cost channel}: \textsc{[x.x\%]}.
        At each relocation event, the E1\_2L household pays
        $\tau_{\text{sell}} + \tau_{\text{buy}} \approx 8.5\%$;
        the E2\_2L household retains tokens at $\tau_{\text{token}} \approx 1\%$.
        This channel is unique to token portability and has no analog in
        \citet{Liu2021}, who abstracts from mobility.

    \item \emph{Maintained-hedge channel}: \textsc{[x.x\%]}.
        A household at location~$B$ who holds $x_{A,\text{prev}} > 0$ continues
        to participate in location~$A$'s price appreciation.
        The mechanism activates through the pre-accumulation motive enabled by
        the 6D state: the household optimally builds $x_B > 0$ at location~$A$
        to reduce future buying costs on relocation (the ``pre-buy hedge'').
        [\textsc{Placeholder: confirm sign via $\bar{x}_B > 0$ at $\ell = A$
        in server1 output (H1 test).}]

    \item \emph{Continuous-ownership channel}: \textsc{[x.x\%]}.
        Continuous fractional ownership of the occupied residence provides
        rent saving of $\delta_{\text{own}} = \rho - m = 4\%$ per unit held.
        This channel is common to \citet{Liu2021} and replicates his
        MHS-relaxation result within our model.
\end{enumerate}
Channels~(1) and~(2) are the paper's structural contribution beyond
\citet{Liu2021}: token portability uniquely provides (1) because direct
owners face forced sales, and REITs provide no location-specific claims
for channel~(2).
The cross-term between channels is empirically negligible (\textsc{[x.xx\%]})
as in our v3 channel-decomposition benchmark, confirming additive separability.

\paragraph{Falsification structure.}
We pre-register three mechanism tests that must pass for the token-portability
claim to be credible.
(r)~Setting $p_{\text{reloc}} = 0$ shuts off the mobility channel;
    channels~(1) and~(2) must collapse to zero.
(m)~Setting $\rho_{AB} \to 1$ eliminates diversification benefit;
    the maintained-hedge channel must attenuate monotonically.
(q)~Setting $\tau_{\text{buy}} = 0$ eliminates the pre-accumulation motive;
    the pre-buy hedge channel must vanish.
[\textsc{Placeholder: report $p$-values or CEV changes from server1 sweep results.}]

\paragraph{Related literature.}
Our paper connects four strands.

\medskip\noindent
\textbf{Lifecycle housing-portfolio choice.}
\citet{YaoZhang2005} and \citet{Cocco2005} establish the canonical
binary-tenure lifecycle model; we extend by adding a second location,
stochastic relocation, and token portability.
\citet{KraftMunk2011} and \citet{KMW2018} introduce habit and
adjustment-cost extensions; we add geographic mobility.
\citet{Liu2021} introduces the Mortgage-with-Housing-Supply (MHS)
relaxation: continuous ownership fraction within a single location.
Our continuous-$x$ channel (E2\_2L at a single location) replicates Liu's
mechanism; our paper's novel contribution lies in channels~(1) and~(2)
above, which are outside Liu's single-location framework.

\medskip\noindent
\textbf{Housing as hedge against mobility risk.}
\citet{SinaiSouleles2005} show that homeownership hedges against rent-price
risk for households who expect to remain in a location.
\citet{Davidoff2006} documents income-housing correlation as a source of
labor-market risk for owner-occupiers.
\citet{BFN2014} (Bagliano, Fugazza, and Nicodano) show that housing and
equity exposure interact through labor income in a lifecycle model.
Our contribution is distinct: we model the \emph{loss} of location-specific
hedge at relocation and quantify the welfare value of maintaining it via tokens.

\medskip\noindent
\textbf{Transaction costs in housing markets.}
\citet{FlavinYamashita2002} document that transaction costs constrain
portfolio rebalancing in a housing lifecycle context.
\citet{Han2013} quantifies the lock-in effect of transaction costs on
tenure choice.
\citet{PiazzesiSchneider2016} embed transaction costs in an equilibrium
housing-search model.
Our model treats transaction costs as the \emph{source} of the E1\_2L
welfare loss; token portability is valued precisely because it reduces these costs.

\medskip\noindent
\textbf{Tokenization and digital assets.}
\citet{CongLiWang2021} provide the foundational theory of tokenomics in a
general equilibrium setting.
\citet{Swinkels2023} surveys the institutional structure of residential
real estate tokenization platforms.
Our contribution is a lifecycle welfare quantification \emph{within} a
standard Cocco-class model, grounded in the PSID-/NAR-calibrated parameter
space, rather than a platform mechanism design.

\paragraph{Partial equilibrium caveat.}
All results are partial-equilibrium: house prices are exogenous.
We interpret $\operatorname{CEV}(E2\_2L \text{ vs } E1\_2L)$ as the
household-level welfare value of \emph{access} to tokenized instruments,
not as an aggregate welfare estimate.
A general-equilibrium extension with endogenous prices and adoption externalities
is left for future work.

\paragraph{Road map.}
Section~\ref{sec:model} presents the two-location lifecycle model.
Section~\ref{sec:calibration} describes the calibration.
Section~\ref{sec:results} reports baseline welfare results and the channel
decomposition.
Section~\ref{sec:discussion} discusses the contribution relative to
\citet{Liu2021} and the robustness of the falsification structure.
Section~\ref{sec:conclusion} concludes.
All numerical details are in the appendices.

% ─────────────────────────────────────────────────────────────────────────────

\clearpage

% ─────────────────────────────────────────────────────────────────────────────
% s2_model.tex — Model section
% Paper: "Tokenized Housing and Lifecycle Portfolio Choice:
%         A Decoupling of Location from Housing Exposure"
%
% Source: docs/methods_v3.md (v3/v4 spec, cloud agent fire 7)
%         docs/calibration_v3.md (parameter table)
%         paper/outline_v4.md (section structure)
%
% Status: complete draft; numerical placeholders ({\sc [X]}) to be filled
%         once server1 baseline results are available.
%
% Compilation: included from main.tex via % s2_model.tex — Model section
% Paper: "Tokenized Housing and Lifecycle Portfolio Choice:
%         A Decoupling of Location from Housing Exposure"
%
% Source: docs/methods_v3.md (v3/v4 spec, cloud agent fire 7)
%         docs/calibration_v3.md (parameter table)
%         paper/outline_v4.md (section structure)
%
% Status: complete draft; numerical placeholders ({\sc [X]}) to be filled
%         once server1 baseline results are available.
%
% Compilation: included from main.tex via % s2_model.tex — Model section
% Paper: "Tokenized Housing and Lifecycle Portfolio Choice:
%         A Decoupling of Location from Housing Exposure"
%
% Source: docs/methods_v3.md (v3/v4 spec, cloud agent fire 7)
%         docs/calibration_v3.md (parameter table)
%         paper/outline_v4.md (section structure)
%
% Status: complete draft; numerical placeholders ({\sc [X]}) to be filled
%         once server1 baseline results are available.
%
% Compilation: included from main.tex via \input{sections/s2_model}
%              Requires amsmath, booktabs, multirow, hyperref.

% ─────────────────────────────────────────────────────────────────────────────
\section{Model}
\label{sec:model}
% ─────────────────────────────────────────────────────────────────────────────

\subsection{Economic Environment}
\label{sec:model:environment}

The economy is populated by a unit mass of households who live for a
finite horizon of $T$ periods, corresponding to ages $25$ through $80$
(retirement at age $65$).
There are two locations, $A$ and $B$, each associated with its own housing
market.
At every date $t$, each household occupies exactly one location
$\ell_t \in \{A, B\}$.
Households have Epstein-Zin preferences that reduce to CRRA over
non-housing consumption in the baseline:
\begin{equation}
    U_t = \mathbb{E}_t\!\left[
        \sum_{s=t}^{T} \beta^{s-t}\,\pi_s\,
        u\!\left(\hat{c}_s\right)
    \right],
    \qquad
    u(c) = \frac{c^{1-\gamma}}{1-\gamma},
    \label{eq:utility}
\end{equation}
where $\beta$ is the time-preference factor, $\pi_s$ is the deterministic
survival probability at age $s$, and $\hat{c}_t$ is non-housing consumption
normalized by the house-price index.%
\footnote{We abstract from housing services in utility and focus on the
financial welfare gain from token portability. An extension with housing
services in utility is discussed in Appendix~\ref{app:robustness}.}

All quantities are normalized by the beginning-of-period house-price index
$P_t$, so the model unit is one house.

\subsection{State Space}
\label{sec:model:state}

The household's state vector in the baseline (\emph{v4 specification}) is
six-dimensional:
\begin{equation}
    \mathbf{s}_t \;=\;
    \bigl(t,\;w_t,\;z_t,\;\ell_t,\;x_{A,t-1},\;x_{B,t-1}\bigr),
    \label{eq:state}
\end{equation}
where:
\begin{itemize}[noitemsep]
    \item $w_t \in [w_{\min}, w_{\max}]$: beginning-of-period normalized
          financial wealth;
    \item $z_t \in [z_{\min}, z_{\max}]$: persistent component of labor income;
    \item $\ell_t \in \{A, B\}$: current location (discrete);
    \item $x_{A,t-1},\, x_{B,t-1} \geq 0$: token holdings at locations $A$ and
          $B$ carried forward from period $t-1$ (the \emph{lagged-position state}).
\end{itemize}
The lagged-position state $(x_{A,t-1}, x_{B,t-1})$ is the key novelty of the
v4 specification.  Tracking prior holdings creates an \emph{intertemporal
incentive}: a household at $\ell = A$ can pre-accumulate $x_{B,t-1} > 0$
at low incremental cost, reducing the buying cost it would face upon
relocation to $B$ in a future period.
Without this state variable, per-period buying costs are approximated as
zero (ignoring the pre-accumulation margin), which was the v3 Option~3
approximation.

At $t = 1$, all households enter with $x_{A,0} = x_{B,0} = 0$ (no prior
token positions).

\subsection{Stochastic Relocation Shock}
\label{sec:model:relocation}

Each period, the household receives a Bernoulli relocation shock:
\begin{equation}
    \mathbf{1}_{\text{reloc},t} \;\sim\; \operatorname{Bernoulli}\!\bigl(p_t^{\text{reloc}}\bigr),
    \qquad
    p_t^{\text{reloc}} = \begin{cases}
        p_{\text{work}} & \text{if}\; t \leq \tau_{\text{ret}}, \\
        p_{\text{ret}}  & \text{if}\; t  >  \tau_{\text{ret}},
    \end{cases}
    \label{eq:reloc_prob}
\end{equation}
where $\tau_{\text{ret}}$ is the retirement age.  In the baseline,
$p_{\text{work}} = 0.06$ (6\,\% per year for working-age households) and
$p_{\text{ret}} = 0.02$.  The working-age rate is anchored to
inter-MSA annual mobility rates in the PSID; see
Table~\ref{tab:calibration} and \citet{calibrationV3} for details.

When relocation occurs ($\mathbf{1}_{\text{reloc},t} = 1$), the household's
location toggles: $\ell_{t+1} = \ell_t' \neq \ell_t$.  The implications
for housing positions differ by regime (Section~\ref{sec:model:regimes}).

\subsection{Regime Taxonomy}
\label{sec:model:regimes}

We compare three regimes that differ in the admissible set of housing
positions:

\begin{table}[ht]
    \centering
    \caption{Regime Taxonomy}
    \label{tab:regimes}
    \renewcommand{\arraystretch}{1.20}
    \begin{tabular}{llll}
        \toprule
        Regime & Description & $x_{A,t}$ admissible & $x_{B,t}$ admissible \\
        \midrule
        E0     & Rent-only         & $0$ & $0$ \\
        E1\_2L & Binary ownership  & $\{0,1\}$ if $\ell=A$; $0$ if $\ell=B$
                                   & $\{0,1\}$ if $\ell=B$; $0$ if $\ell=A$ \\
        E2\_2L & Continuous tokens & $[0, \bar{x}]$ & $[0, \bar{x}]$ \\
        \bottomrule
    \end{tabular}
    \smallskip
    \begin{minipage}{\linewidth}
        \footnotesize
        Notes: $\bar{x}$ is the wealth-constrained maximum (Section~\ref{sec:model:budget}).
        In E1\_2L, the household can own \emph{only} at its current location; the
        non-occupied-location holding is zero by admissibility (reflecting the
        physical-residence requirement of traditional homeownership).
    \end{minipage}
\end{table}

\noindent\textbf{E0 (rent-only).}  The household never owns housing and pays
the full rent-to-price ratio $\rho$ each period.  This regime serves as the
no-ownership benchmark.

\noindent\textbf{E1\_2L (binary homeownership).}  The household owns either
zero or one unit of the property at its \emph{current} location.  It cannot
hold the non-occupied-location property.  This regime represents traditional
homeownership: a household at $A$ cannot hold a financial claim on a
location-$B$ property.  On relocation, the entire owned unit is sold
(at a discount, see Section~\ref{sec:model:txcost}), and the household
arrives at the new location with no prior position ($x_{t-1} = 0$ at
both locations).

\noindent\textbf{E2\_2L (fractional tokens, portable).}  The household can
hold any non-negative fractional amount at \emph{both} locations simultaneously.
Tokens are portable across relocations: moving from $A$ to $B$ does not
trigger a forced sale of location-$A$ tokens.  Transaction costs are charged
on the \emph{change} in holdings each period (see Section~\ref{sec:model:txcost}),
not as a lump-sum on relocation events.

The headline welfare object is $\operatorname{CEV}(\text{E2\_2L vs E1\_2L})$
--- the proportional consumption supplement that makes E1\_2L households
indifferent to E2\_2L access.

\subsection{Housing-Cost Rule}
\label{sec:model:kappa}

Net housing cost per period (as a fraction of house value) depends on the
household's regime and current holdings:
\begin{equation}
    \kappa(x_{A,t}, x_{B,t}, \ell_t) = \begin{cases}
        \rho                           & \text{(E0)} \\[4pt]
        \rho\,\mathbf{1}[x_{\ell,t} < 1]
        + m\,\mathbf{1}[x_{\ell,t} \geq 1]  & \text{(E1\_2L)} \\[4pt]
        \rho - x_{\ell,t}\,(\rho - m) & \text{(E2\_2L)}
    \end{cases}
    \label{eq:kappa}
\end{equation}
where $x_{\ell,t} \equiv x_{A,t}\,\mathbf{1}[\ell_t = A] + x_{B,t}\,\mathbf{1}[\ell_t = B]$
is the token holding at the occupied location, $\rho$ is the rent-to-price
ratio, and $m$ is the maintenance cost.

The parameter $\delta_{\text{own}} \equiv \rho - m$ is the per-unit annual
welfare gain from owning the occupied residence.  In the baseline calibration,
$\delta_{\text{own}} = 0.04$ (4\,\% per year).

\paragraph{The fixed-kappa convention.}
A key identifying restriction is that in E2\_2L, \emph{only} the
occupied-location token $x_{\ell,t}$ reduces net rent.  The non-occupied
token $x_{\ell',t}$ is a purely financial asset — it earns capital gains
but does not reduce rent (the household cannot simultaneously occupy and
rent out the same unit).  This convention, confirmed by falsification
tests in \citet{researchLogMay01}, ensures that cross-location holdings
are motivated exclusively by the financial return and the pre-accumulation
motive, not by a mechanical rent-saving subsidy.

\subsection{Period Budget Constraint}
\label{sec:model:budget}

Each period, the household allocates normalized wealth $w_t$ across
non-housing consumption $\hat{c}_t$, bond savings $b_t$, equity $s_t$,
and housing token purchases $(x_{A,t}, x_{B,t})$:
\begin{equation}
    \hat{c}_t + \kappa(x_{A,t}, x_{B,t}, \ell_t)
    + b_t + s_t + x_{A,t} + x_{B,t} + \tau_t^{\text{tx}} = w_t,
    \label{eq:budget}
\end{equation}
where $\tau_t^{\text{tx}}$ is the per-period transaction cost
(Section~\ref{sec:model:txcost}).  We require $\hat{c}_t > 0$,
$s_t \geq 0$, and $b_t \geq -\ell_{\max} \cdot x_{\ell,t}$ (LTV
mortgage constraint; $\ell_{\max} = 0$ in the baseline).

\subsection{Transaction Costs}
\label{sec:model:txcost}

\paragraph{E1\_2L forced-relocation sell.}
When an E1\_2L household at location $\ell$ relocates to $\ell' \neq \ell$,
it is forced to sell its housing unit.  The sell proceeds are discounted by
a transaction cost $\tau_{\text{sell}} \approx 6\,\%$ (NAR commission):
the sell \emph{factor} applied to location-$\ell$ holdings in the wealth
transition is $(1 - \tau_{\text{sell}})$.
Additionally, on arrival at the new location, the household pays a buying
cost $\tau_{\text{buy}} \approx 2.5\,\%$ of the purchase price when
acquiring a new unit.

\paragraph{E2\_2L per-period transaction cost.}
In the v4 specification, $x_{A,t-1}$ and $x_{B,t-1}$ are tracked as state
variables.  Each period, the household pays a transaction cost based on the
\emph{change} in holdings:
\begin{align}
    \Delta_A &= x_{A,t} - x_{A,t-1}, \qquad
    \Delta_B = x_{B,t} - x_{B,t-1}, \notag \\
    \tau_t^{\text{tx}} &= \tau_{\text{buy}} \,
        \bigl[\max(\Delta_A, 0) + \max(\Delta_B, 0)\bigr]
        + \tau_{\text{token}} \,
        \bigl[\max(-\Delta_A, 0) + \max(-\Delta_B, 0)\bigr],
    \label{eq:txcost}
\end{align}
where $\tau_{\text{buy}} = 0.025$ (per-unit buying cost) and
$\tau_{\text{token}} = 0.01$ (per-unit token-sell cost).
Tokens retained unchanged across periods incur no cost ($\Delta = 0$).

\paragraph{The pre-accumulation motive.}
Equation~\eqref{eq:txcost} creates an \emph{intertemporal incentive} absent
from earlier specifications.  A household at $\ell = A$ today faces an
expected future buying cost of $\tau_{\text{buy}} \cdot x_{B,t}^{*}$ upon
relocation to $B$, where $x_{B,t}^{*}$ is the desired holding at $B$.
By pre-accumulating $x_{B,t-1} > 0$ \emph{before} relocation, the household
reduces this future cost by $\tau_{\text{buy}} \cdot x_{B,t-1}$ per unit
pre-held.  The per-period expected saving is
$p_{\text{work}} \cdot \tau_{\text{buy}} \approx 0.15\,\%$ per unit of
$x_{B,t-1}$ held at $\ell = A$.

Under E1\_2L, cross-location pre-accumulation is inadmissible
(Table~\ref{tab:regimes}), so E1\_2L households cannot exploit this
margin.  This asymmetry is the structural source of the welfare gap.

\subsection{Asset Returns}
\label{sec:model:returns}

\paragraph{Housing returns.}
We decompose the single-location log housing return into a shared aggregate
factor and a location-specific idiosyncratic component:
\begin{align}
    \iota_{A,t+1} &= \sqrt{2}\,\sigma_\iota\,\xi_{A,t+1}, \notag \\
    \iota_{B,t+1} &= \rho_{AB}\,\iota_{A,t+1}
                     + \sqrt{1 - \rho_{AB}^2}\,\sqrt{2}\,\sigma_\iota\,\xi_{B,t+1},
                     \label{eq:chol} \\
    R_{A,t+1} &= \exp\!\bigl(\mu_h + \sqrt{2}\,\sigma_{\text{div}}\,\eta_{\text{div},t+1}
                               + \iota_{A,t+1}\bigr), \notag \\
    R_{B,t+1} &= \exp\!\bigl(\mu_h + \sqrt{2}\,\sigma_{\text{div}}\,\eta_{\text{div},t+1}
                               + \iota_{B,t+1}\bigr), \notag
\end{align}
where $\eta_{\text{div},t+1}$ is the \emph{shared} aggregate housing factor
(same realization for $A$ and $B$), $\xi_{A,t+1}$ and $\xi_{B,t+1}$ are
independent standard normal shocks, and $\sigma_{\text{div}}^2 + \sigma_\iota^2 =
\sigma_h^2$ (variance decomposition).
The Cholesky factorization in~\eqref{eq:chol} implies
$\operatorname{Corr}(\iota_{A,t+1}, \iota_{B,t+1}) = \rho_{AB}$.

The baseline $\rho_{AB} = 0.50$ is the midpoint of the Case-Shiller MSA-pair
idiosyncratic correlation range ($0.30$--$0.70$ for inter-regional US moves;
see \citealt{calibrationV3}).  Sensitivity analysis over
$\rho_{AB} \in \{0, 0.25, 0.50, 0.75, 0.95\}$ is reported in
Section~\ref{sec:results:sensitivity}.

\paragraph{Equity return.}
$R_{S,t+1} = \exp(\mu_s + \sqrt{2}\,\sigma_s\,\eta_{s,t+1})$,
where $\mu_s = \ln(R_f + \pi_s) - \frac{1}{2}\sigma_s^2$ and $\pi_s$
is the equity premium.

\paragraph{House-price normalization.}
The house-price factor $\mathit{hp}_{t+1} = \exp(g_h + \sqrt{2}\,\sigma_\xi\,\xi_{h,t+1})$
normalizes next-period wealth; all financial quantities are expressed in
units of current house prices.

\subsection{Income Process}
\label{sec:model:income}

Labor income follows the \citet{CoccoGomesMaenhout2005} specification.
The deterministic age profile is:
\begin{equation}
    f(\text{age}) = -2.17042 + 0.16818\left(\tfrac{\text{age}}{10}\right)
                   - 0.03230\left(\tfrac{\text{age}}{10}\right)^{\!2}
                   + 0.00200\left(\tfrac{\text{age}}{10}\right)^{\!3}.
    \label{eq:income_profile}
\end{equation}
The normalized persistent income component evolves as:
\begin{align}
    z_{t+1} &= z_t\,\exp\!\bigl(f(\text{age}_{t+1}) - f(\text{age}_t)
               + u_{t+1}\bigr)\big/\mathit{hp}_{t+1},
               & u_{t+1} &\sim \mathcal{N}(0, \sigma_u^2), \notag \\
    y_{t+1} &= z_{t+1}\,\exp(\varepsilon_{t+1}),
               & \varepsilon_{t+1} &\sim \mathcal{N}(0, \sigma_\varepsilon^2).
               \label{eq:income}
\end{align}
At retirement, $z_{t+1} = \lambda_{\text{ret}} z_t / \mathit{hp}_{t+1}$
(replacement rate $\lambda_{\text{ret}} = 0.65$); post-retirement income
is constant at $y_{t+1} = z_t / \mathit{hp}_{t+1}$.

\subsection{Wealth Transition}
\label{sec:model:wealth}

Beginning-of-next-period wealth is:
\begin{equation}
    w_{t+1} = \frac{b_t\,r_b + s_t\,R_{S,t+1}
                    + x_{A,t}\,R_{A,t+1}\,\phi_{A,t+1}
                    + x_{B,t}\,R_{B,t+1}\,\phi_{B,t+1}}
                   {\mathit{hp}_{t+1}}
              + y_{t+1},
    \label{eq:wealth_transition}
\end{equation}
where $r_b = R_f$ for $b_t \geq 0$ and $r_b = R_f + r_{\text{prem}}$
for $b_t < 0$, and $\phi_{\ell,t+1}$ is the sell factor:
\begin{equation}
    \phi_{\ell,t+1} = \begin{cases}
        1 - \tau_{\text{sell}} & \text{if E1\_2L and relocating from } \ell, \\
        1 & \text{otherwise.}
    \end{cases}
    \label{eq:sell_factor}
\end{equation}
Tokens in E2\_2L are portable: $\phi = 1$ always.  The forced-sell discount
$(1 - \tau_{\text{sell}})$ under E1\_2L is the first component of the
transaction-cost channel in the welfare decomposition
(Section~\ref{sec:model:welfare}).

\subsection{Bellman Equation}
\label{sec:model:bellman}

For regime $R \in \{\text{E0, E1\_2L, E2\_2L}\}$ and v4 state
$\mathbf{s}_t = (t, w_t, z_t, \ell_t, x_{A,t-1}, x_{B,t-1})$:
\begin{equation}
    V_t^R\!\bigl(w, z, \ell, x_{Ap}, x_{Bp}\bigr)
    = \max_{\hat{c},\, b,\, s,\, x_A,\, x_B \,\in\, \mathcal{A}_R}
    \left[
        u(\hat{c}) + \beta\,\pi_t\,
        \mathcal{E}_t^R\!\bigl(w, z, \ell, x_A, x_B\bigr)
    \right],
    \label{eq:bellman}
\end{equation}
where $\mathcal{A}_R$ is the regime-specific admissible set
(Table~\ref{tab:regimes} and budget constraint~\eqref{eq:budget}),
and the continuation value is:
\begin{multline}
    \mathcal{E}_t^R\!\bigl(w, z, \ell, x_A, x_B\bigr) \\
    = \sum_q \omega_q\,\mathit{hp}_q^{1-\gamma}
    \Bigl[
        (1 - p_t^{\text{reloc}})\,
        V_{t+1}^R\!\bigl(w_q^{\text{stay}},  z_q', \ell,  x_A^{\ast}, x_B^{\ast}\bigr)
        + p_t^{\text{reloc}}\,
        V_{t+1}^R\!\bigl(w_q^{\text{reloc}}, z_q', \ell', x_A^{\dagger}, x_B^{\dagger}\bigr)
    \Bigr],
    \label{eq:ev}
\end{multline}
where $(\omega_q, \mathit{hp}_q, w_q^{\text{stay}}, w_q^{\text{reloc}}, z_q')$
are Gauss-Hermite quadrature weights and realizations (7D; $3^7 = 2187$
integration points), and the next-period $x$-state updates are:
\begin{align}
    (x_A^{\ast}, x_B^{\ast})   &= (x_A, x_B),
        \quad &\text{(stay; E0, E2\_2L, E1\_2L)} \label{eq:xprev_stay} \\
    (x_A^{\dagger}, x_B^{\dagger}) &= (x_A, x_B),
        \quad &\text{(reloc; E0, E2\_2L — portable)} \label{eq:xprev_reloc_e2} \\
    (x_A^{\dagger}, x_B^{\dagger}) &= (0, 0),
        \quad &\text{(reloc; E1\_2L — forced sale, fresh start)}
        \label{eq:xprev_reloc_e1}
\end{align}
Equation~\eqref{eq:xprev_reloc_e1} captures the key asymmetry: under
E1\_2L, relocation resets the lagged-position state to zero (the household
must buy fresh at the new location, paying $\tau_{\text{buy}}$ on the full
amount), while under E2\_2L, equation~\eqref{eq:xprev_reloc_e2} shows that
tokens carry forward unchanged.

The terminal condition is $V_{T+1}^R(w, \cdot) = u(w)$
(bequest-free, all wealth consumed at death).

\subsection{Welfare Measure and Decomposition}
\label{sec:model:welfare}

\paragraph{Certainty-equivalent consumption (CEV).}
The welfare comparison between regimes $R_a$ and $R_b$ is measured by the
proportional consumption supplement $\lambda$ such that:
\begin{equation}
    V_1^{R_b}\!\bigl((1 + \lambda)\,w^{\ast}\bigr) = V_1^{R_a}\!\bigl(w^{\ast}\bigr),
    \label{eq:cev_def}
\end{equation}
evaluated at the representative initial state $w^{\ast}$ (midpoint of the
wealth grid, $z = z_{\text{mid}}$, $\ell = A$, $x_{Ap} = x_{Bp} = 0$).
Under CRRA:
\begin{equation}
    1 + \lambda = \left(\frac{V_1^{R_a}}{V_1^{R_b}}\right)^{\!\!1/(1-\gamma)}.
    \label{eq:cev_crra}
\end{equation}

\paragraph{Channel decomposition.}
The headline CEV $\operatorname{CEV}(\text{E2\_2L vs E1\_2L})$ decomposes into
two channels via the counterfactual regime E1\_2L$_{\text{NOTX}}$ (E1\_2L
with $\tau_{\text{sell}} = 0$):
\begin{align}
    \underbrace{\operatorname{CEV}(\text{E2\_2L vs E1\_2L})}_{\text{total}}
    &=
    \underbrace{\operatorname{CEV}(\text{E1\_2L}_{\text{NOTX}} \text{ vs E1\_2L})}_{\text{avoided-tx channel}}
    +
    \underbrace{\operatorname{CEV}(\text{E2\_2L vs E1\_2L}_{\text{NOTX}})}_{\text{maintained-hedge channel}}
    + \xi_{\text{cross}},
    \label{eq:decomp}
\end{align}
where $\xi_{\text{cross}}$ is the cross-term, which was empirically near zero
($< 0.1\%$) in the v3 specification and is expected to remain small in v4.
We report $\xi_{\text{cross}}$ explicitly to verify separability.

Table~\ref{tab:cev_decomp} reports the decomposition with numerical results
to be filled after server1 baseline runs are complete.

\begin{table}[ht]
    \centering
    \caption{CEV Decomposition: Tokenization Welfare Channels}
    \label{tab:cev_decomp}
    \renewcommand{\arraystretch}{1.20}
    \begin{tabular}{lcc}
        \toprule
        Channel & CEV (\%) & Share (\%) \\
        \midrule
        Total tokenization gain:
            $\operatorname{CEV}(\text{E2\_2L vs E1\_2L})$
            & {\sc [x.xx]} & 100 \\
        \quad Avoided-tx channel:
            $\operatorname{CEV}(\text{E1\_2L}_{\text{NOTX}} \text{ vs E1\_2L})$
            & {\sc [x.xx]} & {\sc [xx]} \\
        \quad Maintained-hedge channel:
            $\operatorname{CEV}(\text{E2\_2L vs E1\_2L}_{\text{NOTX}})$
            & {\sc [x.xx]} & {\sc [xx]} \\
        \quad Cross-term $\xi_{\text{cross}}$
            & {\sc [x.xx]} & {\sc [xx]} \\
        \midrule
        Option 1 incremental gain:
            $\operatorname{CEV}(\text{E2\_2L v4 vs E2\_2L v3})$
            & {\sc [x.xx]} & --- \\
        \bottomrule
    \end{tabular}
    \smallskip
    \begin{minipage}{\linewidth}
        \footnotesize
        Notes: CEVs evaluated at representative initial state
        $(w^{\ast}, z_{\text{mid}}, \ell = A, x_{Ap} = x_{Bp} = 0)$
        using the v4 solver (6D state; $\text{N\_W}=15$,
        $\text{N\_Z}=5$, $\text{N\_X\_PREV}=3$, $\text{GH\_NODES}=3$).
        {\sc [x.xx]}: placeholder to be filled from server1 runs.
        Baseline calibration: Table~\ref{tab:calibration}.
    \end{minipage}
\end{table}

\paragraph{Hypothesis tests.}
We pre-register three empirical hypotheses (H1--H3) that must hold for
the option-1 mechanism to support the RFS-level contribution claim:
\begin{description}[noitemsep]
    \item[H1] $\overline{x}_{B,1}^{\text{E2}} > 0$ at $\ell = A$
              (cross-location pre-accumulation activates).
    \item[H2] $\operatorname{CEV}(\text{E2\_2L v4 vs E1\_2L v4}) > 4.255\,\%$
              (beats the Option~3 approximation baseline).
    \item[H3] $\operatorname{CEV}(\text{E2\_2L v4 vs E2\_2L v3}) \in [0.5, 1.5]\,\%$
              (incremental hedge value from proper state extension).
\end{description}
If all three hold: the mechanism is RFS-credible and Phase~2
(calibration, sensitivity, manuscript) proceeds.  If any fails, the
paper targets Real Estate Economics or Journal of Housing Economics
with the Option~3 result ($+4.26\,\%$).

\subsection{Numerical Solution}
\label{sec:model:numerical}

We solve the Bellman equation~\eqref{eq:bellman} by backward induction
(value function iteration) using the v4 Julia solver
(\texttt{src/vfi\_solver\_v4.jl}, 954 LOC).
The value function is approximated on a discrete 6D grid; continuation
values are computed by 7D Gauss-Hermite quadrature ($3^7 = 2187$ points).

\paragraph{Interpolation.}
The continuation-value lookup interpolates in $(w, z)$ bilinearly for
a given $(\ell, x_{Ap}, x_{Bp})$ grid cell.  The lagged-position state
$(x_{Ap}, x_{Bp})$ uses a coarse grid with $\text{N\_X\_PREV} = 3$
points ($\{0, 1, 2\}$ in the baseline), sufficient for the first-cut
test of H1.  Sensitivity to grid refinement is reported in
Appendix~\ref{app:grid}.

\paragraph{Grid sizes and computation.}
Baseline coarse grid: $\text{N\_W}=15$, $\text{N\_Z}=5$,
$\text{N\_X\_PREV}=3$, $\text{ASSET\_GRID}=7$, $\text{X\_GRID}=4$.
State-space size per period: $15 \times 5 \times 2 \times 3 \times 3 = 1350$
points (versus 294 in v3 with $\text{N\_W}=21$, $\text{N\_Z}=7$,
$\text{N\_ell}=2$).
Estimated wall time per regime: 2--4 hours single-threaded on server1.

%              Requires amsmath, booktabs, multirow, hyperref.

% ─────────────────────────────────────────────────────────────────────────────
\section{Model}
\label{sec:model}
% ─────────────────────────────────────────────────────────────────────────────

\subsection{Economic Environment}
\label{sec:model:environment}

The economy is populated by a unit mass of households who live for a
finite horizon of $T$ periods, corresponding to ages $25$ through $80$
(retirement at age $65$).
There are two locations, $A$ and $B$, each associated with its own housing
market.
At every date $t$, each household occupies exactly one location
$\ell_t \in \{A, B\}$.
Households have Epstein-Zin preferences that reduce to CRRA over
non-housing consumption in the baseline:
\begin{equation}
    U_t = \mathbb{E}_t\!\left[
        \sum_{s=t}^{T} \beta^{s-t}\,\pi_s\,
        u\!\left(\hat{c}_s\right)
    \right],
    \qquad
    u(c) = \frac{c^{1-\gamma}}{1-\gamma},
    \label{eq:utility}
\end{equation}
where $\beta$ is the time-preference factor, $\pi_s$ is the deterministic
survival probability at age $s$, and $\hat{c}_t$ is non-housing consumption
normalized by the house-price index.%
\footnote{We abstract from housing services in utility and focus on the
financial welfare gain from token portability. An extension with housing
services in utility is discussed in Appendix~\ref{app:robustness}.}

All quantities are normalized by the beginning-of-period house-price index
$P_t$, so the model unit is one house.

\subsection{State Space}
\label{sec:model:state}

The household's state vector in the baseline (\emph{v4 specification}) is
six-dimensional:
\begin{equation}
    \mathbf{s}_t \;=\;
    \bigl(t,\;w_t,\;z_t,\;\ell_t,\;x_{A,t-1},\;x_{B,t-1}\bigr),
    \label{eq:state}
\end{equation}
where:
\begin{itemize}[noitemsep]
    \item $w_t \in [w_{\min}, w_{\max}]$: beginning-of-period normalized
          financial wealth;
    \item $z_t \in [z_{\min}, z_{\max}]$: persistent component of labor income;
    \item $\ell_t \in \{A, B\}$: current location (discrete);
    \item $x_{A,t-1},\, x_{B,t-1} \geq 0$: token holdings at locations $A$ and
          $B$ carried forward from period $t-1$ (the \emph{lagged-position state}).
\end{itemize}
The lagged-position state $(x_{A,t-1}, x_{B,t-1})$ is the key novelty of the
v4 specification.  Tracking prior holdings creates an \emph{intertemporal
incentive}: a household at $\ell = A$ can pre-accumulate $x_{B,t-1} > 0$
at low incremental cost, reducing the buying cost it would face upon
relocation to $B$ in a future period.
Without this state variable, per-period buying costs are approximated as
zero (ignoring the pre-accumulation margin), which was the v3 Option~3
approximation.

At $t = 1$, all households enter with $x_{A,0} = x_{B,0} = 0$ (no prior
token positions).

\subsection{Stochastic Relocation Shock}
\label{sec:model:relocation}

Each period, the household receives a Bernoulli relocation shock:
\begin{equation}
    \mathbf{1}_{\text{reloc},t} \;\sim\; \operatorname{Bernoulli}\!\bigl(p_t^{\text{reloc}}\bigr),
    \qquad
    p_t^{\text{reloc}} = \begin{cases}
        p_{\text{work}} & \text{if}\; t \leq \tau_{\text{ret}}, \\
        p_{\text{ret}}  & \text{if}\; t  >  \tau_{\text{ret}},
    \end{cases}
    \label{eq:reloc_prob}
\end{equation}
where $\tau_{\text{ret}}$ is the retirement age.  In the baseline,
$p_{\text{work}} = 0.06$ (6\,\% per year for working-age households) and
$p_{\text{ret}} = 0.02$.  The working-age rate is anchored to
inter-MSA annual mobility rates in the PSID; see
Table~\ref{tab:calibration} and \citet{calibrationV3} for details.

When relocation occurs ($\mathbf{1}_{\text{reloc},t} = 1$), the household's
location toggles: $\ell_{t+1} = \ell_t' \neq \ell_t$.  The implications
for housing positions differ by regime (Section~\ref{sec:model:regimes}).

\subsection{Regime Taxonomy}
\label{sec:model:regimes}

We compare three regimes that differ in the admissible set of housing
positions:

\begin{table}[ht]
    \centering
    \caption{Regime Taxonomy}
    \label{tab:regimes}
    \renewcommand{\arraystretch}{1.20}
    \begin{tabular}{llll}
        \toprule
        Regime & Description & $x_{A,t}$ admissible & $x_{B,t}$ admissible \\
        \midrule
        E0     & Rent-only         & $0$ & $0$ \\
        E1\_2L & Binary ownership  & $\{0,1\}$ if $\ell=A$; $0$ if $\ell=B$
                                   & $\{0,1\}$ if $\ell=B$; $0$ if $\ell=A$ \\
        E2\_2L & Continuous tokens & $[0, \bar{x}]$ & $[0, \bar{x}]$ \\
        \bottomrule
    \end{tabular}
    \smallskip
    \begin{minipage}{\linewidth}
        \footnotesize
        Notes: $\bar{x}$ is the wealth-constrained maximum (Section~\ref{sec:model:budget}).
        In E1\_2L, the household can own \emph{only} at its current location; the
        non-occupied-location holding is zero by admissibility (reflecting the
        physical-residence requirement of traditional homeownership).
    \end{minipage}
\end{table}

\noindent\textbf{E0 (rent-only).}  The household never owns housing and pays
the full rent-to-price ratio $\rho$ each period.  This regime serves as the
no-ownership benchmark.

\noindent\textbf{E1\_2L (binary homeownership).}  The household owns either
zero or one unit of the property at its \emph{current} location.  It cannot
hold the non-occupied-location property.  This regime represents traditional
homeownership: a household at $A$ cannot hold a financial claim on a
location-$B$ property.  On relocation, the entire owned unit is sold
(at a discount, see Section~\ref{sec:model:txcost}), and the household
arrives at the new location with no prior position ($x_{t-1} = 0$ at
both locations).

\noindent\textbf{E2\_2L (fractional tokens, portable).}  The household can
hold any non-negative fractional amount at \emph{both} locations simultaneously.
Tokens are portable across relocations: moving from $A$ to $B$ does not
trigger a forced sale of location-$A$ tokens.  Transaction costs are charged
on the \emph{change} in holdings each period (see Section~\ref{sec:model:txcost}),
not as a lump-sum on relocation events.

The headline welfare object is $\operatorname{CEV}(\text{E2\_2L vs E1\_2L})$
--- the proportional consumption supplement that makes E1\_2L households
indifferent to E2\_2L access.

\subsection{Housing-Cost Rule}
\label{sec:model:kappa}

Net housing cost per period (as a fraction of house value) depends on the
household's regime and current holdings:
\begin{equation}
    \kappa(x_{A,t}, x_{B,t}, \ell_t) = \begin{cases}
        \rho                           & \text{(E0)} \\[4pt]
        \rho\,\mathbf{1}[x_{\ell,t} < 1]
        + m\,\mathbf{1}[x_{\ell,t} \geq 1]  & \text{(E1\_2L)} \\[4pt]
        \rho - x_{\ell,t}\,(\rho - m) & \text{(E2\_2L)}
    \end{cases}
    \label{eq:kappa}
\end{equation}
where $x_{\ell,t} \equiv x_{A,t}\,\mathbf{1}[\ell_t = A] + x_{B,t}\,\mathbf{1}[\ell_t = B]$
is the token holding at the occupied location, $\rho$ is the rent-to-price
ratio, and $m$ is the maintenance cost.

The parameter $\delta_{\text{own}} \equiv \rho - m$ is the per-unit annual
welfare gain from owning the occupied residence.  In the baseline calibration,
$\delta_{\text{own}} = 0.04$ (4\,\% per year).

\paragraph{The fixed-kappa convention.}
A key identifying restriction is that in E2\_2L, \emph{only} the
occupied-location token $x_{\ell,t}$ reduces net rent.  The non-occupied
token $x_{\ell',t}$ is a purely financial asset — it earns capital gains
but does not reduce rent (the household cannot simultaneously occupy and
rent out the same unit).  This convention, confirmed by falsification
tests in \citet{researchLogMay01}, ensures that cross-location holdings
are motivated exclusively by the financial return and the pre-accumulation
motive, not by a mechanical rent-saving subsidy.

\subsection{Period Budget Constraint}
\label{sec:model:budget}

Each period, the household allocates normalized wealth $w_t$ across
non-housing consumption $\hat{c}_t$, bond savings $b_t$, equity $s_t$,
and housing token purchases $(x_{A,t}, x_{B,t})$:
\begin{equation}
    \hat{c}_t + \kappa(x_{A,t}, x_{B,t}, \ell_t)
    + b_t + s_t + x_{A,t} + x_{B,t} + \tau_t^{\text{tx}} = w_t,
    \label{eq:budget}
\end{equation}
where $\tau_t^{\text{tx}}$ is the per-period transaction cost
(Section~\ref{sec:model:txcost}).  We require $\hat{c}_t > 0$,
$s_t \geq 0$, and $b_t \geq -\ell_{\max} \cdot x_{\ell,t}$ (LTV
mortgage constraint; $\ell_{\max} = 0$ in the baseline).

\subsection{Transaction Costs}
\label{sec:model:txcost}

\paragraph{E1\_2L forced-relocation sell.}
When an E1\_2L household at location $\ell$ relocates to $\ell' \neq \ell$,
it is forced to sell its housing unit.  The sell proceeds are discounted by
a transaction cost $\tau_{\text{sell}} \approx 6\,\%$ (NAR commission):
the sell \emph{factor} applied to location-$\ell$ holdings in the wealth
transition is $(1 - \tau_{\text{sell}})$.
Additionally, on arrival at the new location, the household pays a buying
cost $\tau_{\text{buy}} \approx 2.5\,\%$ of the purchase price when
acquiring a new unit.

\paragraph{E2\_2L per-period transaction cost.}
In the v4 specification, $x_{A,t-1}$ and $x_{B,t-1}$ are tracked as state
variables.  Each period, the household pays a transaction cost based on the
\emph{change} in holdings:
\begin{align}
    \Delta_A &= x_{A,t} - x_{A,t-1}, \qquad
    \Delta_B = x_{B,t} - x_{B,t-1}, \notag \\
    \tau_t^{\text{tx}} &= \tau_{\text{buy}} \,
        \bigl[\max(\Delta_A, 0) + \max(\Delta_B, 0)\bigr]
        + \tau_{\text{token}} \,
        \bigl[\max(-\Delta_A, 0) + \max(-\Delta_B, 0)\bigr],
    \label{eq:txcost}
\end{align}
where $\tau_{\text{buy}} = 0.025$ (per-unit buying cost) and
$\tau_{\text{token}} = 0.01$ (per-unit token-sell cost).
Tokens retained unchanged across periods incur no cost ($\Delta = 0$).

\paragraph{The pre-accumulation motive.}
Equation~\eqref{eq:txcost} creates an \emph{intertemporal incentive} absent
from earlier specifications.  A household at $\ell = A$ today faces an
expected future buying cost of $\tau_{\text{buy}} \cdot x_{B,t}^{*}$ upon
relocation to $B$, where $x_{B,t}^{*}$ is the desired holding at $B$.
By pre-accumulating $x_{B,t-1} > 0$ \emph{before} relocation, the household
reduces this future cost by $\tau_{\text{buy}} \cdot x_{B,t-1}$ per unit
pre-held.  The per-period expected saving is
$p_{\text{work}} \cdot \tau_{\text{buy}} \approx 0.15\,\%$ per unit of
$x_{B,t-1}$ held at $\ell = A$.

Under E1\_2L, cross-location pre-accumulation is inadmissible
(Table~\ref{tab:regimes}), so E1\_2L households cannot exploit this
margin.  This asymmetry is the structural source of the welfare gap.

\subsection{Asset Returns}
\label{sec:model:returns}

\paragraph{Housing returns.}
We decompose the single-location log housing return into a shared aggregate
factor and a location-specific idiosyncratic component:
\begin{align}
    \iota_{A,t+1} &= \sqrt{2}\,\sigma_\iota\,\xi_{A,t+1}, \notag \\
    \iota_{B,t+1} &= \rho_{AB}\,\iota_{A,t+1}
                     + \sqrt{1 - \rho_{AB}^2}\,\sqrt{2}\,\sigma_\iota\,\xi_{B,t+1},
                     \label{eq:chol} \\
    R_{A,t+1} &= \exp\!\bigl(\mu_h + \sqrt{2}\,\sigma_{\text{div}}\,\eta_{\text{div},t+1}
                               + \iota_{A,t+1}\bigr), \notag \\
    R_{B,t+1} &= \exp\!\bigl(\mu_h + \sqrt{2}\,\sigma_{\text{div}}\,\eta_{\text{div},t+1}
                               + \iota_{B,t+1}\bigr), \notag
\end{align}
where $\eta_{\text{div},t+1}$ is the \emph{shared} aggregate housing factor
(same realization for $A$ and $B$), $\xi_{A,t+1}$ and $\xi_{B,t+1}$ are
independent standard normal shocks, and $\sigma_{\text{div}}^2 + \sigma_\iota^2 =
\sigma_h^2$ (variance decomposition).
The Cholesky factorization in~\eqref{eq:chol} implies
$\operatorname{Corr}(\iota_{A,t+1}, \iota_{B,t+1}) = \rho_{AB}$.

The baseline $\rho_{AB} = 0.50$ is the midpoint of the Case-Shiller MSA-pair
idiosyncratic correlation range ($0.30$--$0.70$ for inter-regional US moves;
see \citealt{calibrationV3}).  Sensitivity analysis over
$\rho_{AB} \in \{0, 0.25, 0.50, 0.75, 0.95\}$ is reported in
Section~\ref{sec:results:sensitivity}.

\paragraph{Equity return.}
$R_{S,t+1} = \exp(\mu_s + \sqrt{2}\,\sigma_s\,\eta_{s,t+1})$,
where $\mu_s = \ln(R_f + \pi_s) - \frac{1}{2}\sigma_s^2$ and $\pi_s$
is the equity premium.

\paragraph{House-price normalization.}
The house-price factor $\mathit{hp}_{t+1} = \exp(g_h + \sqrt{2}\,\sigma_\xi\,\xi_{h,t+1})$
normalizes next-period wealth; all financial quantities are expressed in
units of current house prices.

\subsection{Income Process}
\label{sec:model:income}

Labor income follows the \citet{CoccoGomesMaenhout2005} specification.
The deterministic age profile is:
\begin{equation}
    f(\text{age}) = -2.17042 + 0.16818\left(\tfrac{\text{age}}{10}\right)
                   - 0.03230\left(\tfrac{\text{age}}{10}\right)^{\!2}
                   + 0.00200\left(\tfrac{\text{age}}{10}\right)^{\!3}.
    \label{eq:income_profile}
\end{equation}
The normalized persistent income component evolves as:
\begin{align}
    z_{t+1} &= z_t\,\exp\!\bigl(f(\text{age}_{t+1}) - f(\text{age}_t)
               + u_{t+1}\bigr)\big/\mathit{hp}_{t+1},
               & u_{t+1} &\sim \mathcal{N}(0, \sigma_u^2), \notag \\
    y_{t+1} &= z_{t+1}\,\exp(\varepsilon_{t+1}),
               & \varepsilon_{t+1} &\sim \mathcal{N}(0, \sigma_\varepsilon^2).
               \label{eq:income}
\end{align}
At retirement, $z_{t+1} = \lambda_{\text{ret}} z_t / \mathit{hp}_{t+1}$
(replacement rate $\lambda_{\text{ret}} = 0.65$); post-retirement income
is constant at $y_{t+1} = z_t / \mathit{hp}_{t+1}$.

\subsection{Wealth Transition}
\label{sec:model:wealth}

Beginning-of-next-period wealth is:
\begin{equation}
    w_{t+1} = \frac{b_t\,r_b + s_t\,R_{S,t+1}
                    + x_{A,t}\,R_{A,t+1}\,\phi_{A,t+1}
                    + x_{B,t}\,R_{B,t+1}\,\phi_{B,t+1}}
                   {\mathit{hp}_{t+1}}
              + y_{t+1},
    \label{eq:wealth_transition}
\end{equation}
where $r_b = R_f$ for $b_t \geq 0$ and $r_b = R_f + r_{\text{prem}}$
for $b_t < 0$, and $\phi_{\ell,t+1}$ is the sell factor:
\begin{equation}
    \phi_{\ell,t+1} = \begin{cases}
        1 - \tau_{\text{sell}} & \text{if E1\_2L and relocating from } \ell, \\
        1 & \text{otherwise.}
    \end{cases}
    \label{eq:sell_factor}
\end{equation}
Tokens in E2\_2L are portable: $\phi = 1$ always.  The forced-sell discount
$(1 - \tau_{\text{sell}})$ under E1\_2L is the first component of the
transaction-cost channel in the welfare decomposition
(Section~\ref{sec:model:welfare}).

\subsection{Bellman Equation}
\label{sec:model:bellman}

For regime $R \in \{\text{E0, E1\_2L, E2\_2L}\}$ and v4 state
$\mathbf{s}_t = (t, w_t, z_t, \ell_t, x_{A,t-1}, x_{B,t-1})$:
\begin{equation}
    V_t^R\!\bigl(w, z, \ell, x_{Ap}, x_{Bp}\bigr)
    = \max_{\hat{c},\, b,\, s,\, x_A,\, x_B \,\in\, \mathcal{A}_R}
    \left[
        u(\hat{c}) + \beta\,\pi_t\,
        \mathcal{E}_t^R\!\bigl(w, z, \ell, x_A, x_B\bigr)
    \right],
    \label{eq:bellman}
\end{equation}
where $\mathcal{A}_R$ is the regime-specific admissible set
(Table~\ref{tab:regimes} and budget constraint~\eqref{eq:budget}),
and the continuation value is:
\begin{multline}
    \mathcal{E}_t^R\!\bigl(w, z, \ell, x_A, x_B\bigr) \\
    = \sum_q \omega_q\,\mathit{hp}_q^{1-\gamma}
    \Bigl[
        (1 - p_t^{\text{reloc}})\,
        V_{t+1}^R\!\bigl(w_q^{\text{stay}},  z_q', \ell,  x_A^{\ast}, x_B^{\ast}\bigr)
        + p_t^{\text{reloc}}\,
        V_{t+1}^R\!\bigl(w_q^{\text{reloc}}, z_q', \ell', x_A^{\dagger}, x_B^{\dagger}\bigr)
    \Bigr],
    \label{eq:ev}
\end{multline}
where $(\omega_q, \mathit{hp}_q, w_q^{\text{stay}}, w_q^{\text{reloc}}, z_q')$
are Gauss-Hermite quadrature weights and realizations (7D; $3^7 = 2187$
integration points), and the next-period $x$-state updates are:
\begin{align}
    (x_A^{\ast}, x_B^{\ast})   &= (x_A, x_B),
        \quad &\text{(stay; E0, E2\_2L, E1\_2L)} \label{eq:xprev_stay} \\
    (x_A^{\dagger}, x_B^{\dagger}) &= (x_A, x_B),
        \quad &\text{(reloc; E0, E2\_2L — portable)} \label{eq:xprev_reloc_e2} \\
    (x_A^{\dagger}, x_B^{\dagger}) &= (0, 0),
        \quad &\text{(reloc; E1\_2L — forced sale, fresh start)}
        \label{eq:xprev_reloc_e1}
\end{align}
Equation~\eqref{eq:xprev_reloc_e1} captures the key asymmetry: under
E1\_2L, relocation resets the lagged-position state to zero (the household
must buy fresh at the new location, paying $\tau_{\text{buy}}$ on the full
amount), while under E2\_2L, equation~\eqref{eq:xprev_reloc_e2} shows that
tokens carry forward unchanged.

The terminal condition is $V_{T+1}^R(w, \cdot) = u(w)$
(bequest-free, all wealth consumed at death).

\subsection{Welfare Measure and Decomposition}
\label{sec:model:welfare}

\paragraph{Certainty-equivalent consumption (CEV).}
The welfare comparison between regimes $R_a$ and $R_b$ is measured by the
proportional consumption supplement $\lambda$ such that:
\begin{equation}
    V_1^{R_b}\!\bigl((1 + \lambda)\,w^{\ast}\bigr) = V_1^{R_a}\!\bigl(w^{\ast}\bigr),
    \label{eq:cev_def}
\end{equation}
evaluated at the representative initial state $w^{\ast}$ (midpoint of the
wealth grid, $z = z_{\text{mid}}$, $\ell = A$, $x_{Ap} = x_{Bp} = 0$).
Under CRRA:
\begin{equation}
    1 + \lambda = \left(\frac{V_1^{R_a}}{V_1^{R_b}}\right)^{\!\!1/(1-\gamma)}.
    \label{eq:cev_crra}
\end{equation}

\paragraph{Channel decomposition.}
The headline CEV $\operatorname{CEV}(\text{E2\_2L vs E1\_2L})$ decomposes into
two channels via the counterfactual regime E1\_2L$_{\text{NOTX}}$ (E1\_2L
with $\tau_{\text{sell}} = 0$):
\begin{align}
    \underbrace{\operatorname{CEV}(\text{E2\_2L vs E1\_2L})}_{\text{total}}
    &=
    \underbrace{\operatorname{CEV}(\text{E1\_2L}_{\text{NOTX}} \text{ vs E1\_2L})}_{\text{avoided-tx channel}}
    +
    \underbrace{\operatorname{CEV}(\text{E2\_2L vs E1\_2L}_{\text{NOTX}})}_{\text{maintained-hedge channel}}
    + \xi_{\text{cross}},
    \label{eq:decomp}
\end{align}
where $\xi_{\text{cross}}$ is the cross-term, which was empirically near zero
($< 0.1\%$) in the v3 specification and is expected to remain small in v4.
We report $\xi_{\text{cross}}$ explicitly to verify separability.

Table~\ref{tab:cev_decomp} reports the decomposition with numerical results
to be filled after server1 baseline runs are complete.

\begin{table}[ht]
    \centering
    \caption{CEV Decomposition: Tokenization Welfare Channels}
    \label{tab:cev_decomp}
    \renewcommand{\arraystretch}{1.20}
    \begin{tabular}{lcc}
        \toprule
        Channel & CEV (\%) & Share (\%) \\
        \midrule
        Total tokenization gain:
            $\operatorname{CEV}(\text{E2\_2L vs E1\_2L})$
            & {\sc [x.xx]} & 100 \\
        \quad Avoided-tx channel:
            $\operatorname{CEV}(\text{E1\_2L}_{\text{NOTX}} \text{ vs E1\_2L})$
            & {\sc [x.xx]} & {\sc [xx]} \\
        \quad Maintained-hedge channel:
            $\operatorname{CEV}(\text{E2\_2L vs E1\_2L}_{\text{NOTX}})$
            & {\sc [x.xx]} & {\sc [xx]} \\
        \quad Cross-term $\xi_{\text{cross}}$
            & {\sc [x.xx]} & {\sc [xx]} \\
        \midrule
        Option 1 incremental gain:
            $\operatorname{CEV}(\text{E2\_2L v4 vs E2\_2L v3})$
            & {\sc [x.xx]} & --- \\
        \bottomrule
    \end{tabular}
    \smallskip
    \begin{minipage}{\linewidth}
        \footnotesize
        Notes: CEVs evaluated at representative initial state
        $(w^{\ast}, z_{\text{mid}}, \ell = A, x_{Ap} = x_{Bp} = 0)$
        using the v4 solver (6D state; $\text{N\_W}=15$,
        $\text{N\_Z}=5$, $\text{N\_X\_PREV}=3$, $\text{GH\_NODES}=3$).
        {\sc [x.xx]}: placeholder to be filled from server1 runs.
        Baseline calibration: Table~\ref{tab:calibration}.
    \end{minipage}
\end{table}

\paragraph{Hypothesis tests.}
We pre-register three empirical hypotheses (H1--H3) that must hold for
the option-1 mechanism to support the RFS-level contribution claim:
\begin{description}[noitemsep]
    \item[H1] $\overline{x}_{B,1}^{\text{E2}} > 0$ at $\ell = A$
              (cross-location pre-accumulation activates).
    \item[H2] $\operatorname{CEV}(\text{E2\_2L v4 vs E1\_2L v4}) > 4.255\,\%$
              (beats the Option~3 approximation baseline).
    \item[H3] $\operatorname{CEV}(\text{E2\_2L v4 vs E2\_2L v3}) \in [0.5, 1.5]\,\%$
              (incremental hedge value from proper state extension).
\end{description}
If all three hold: the mechanism is RFS-credible and Phase~2
(calibration, sensitivity, manuscript) proceeds.  If any fails, the
paper targets Real Estate Economics or Journal of Housing Economics
with the Option~3 result ($+4.26\,\%$).

\subsection{Numerical Solution}
\label{sec:model:numerical}

We solve the Bellman equation~\eqref{eq:bellman} by backward induction
(value function iteration) using the v4 Julia solver
(\texttt{src/vfi\_solver\_v4.jl}, 954 LOC).
The value function is approximated on a discrete 6D grid; continuation
values are computed by 7D Gauss-Hermite quadrature ($3^7 = 2187$ points).

\paragraph{Interpolation.}
The continuation-value lookup interpolates in $(w, z)$ bilinearly for
a given $(\ell, x_{Ap}, x_{Bp})$ grid cell.  The lagged-position state
$(x_{Ap}, x_{Bp})$ uses a coarse grid with $\text{N\_X\_PREV} = 3$
points ($\{0, 1, 2\}$ in the baseline), sufficient for the first-cut
test of H1.  Sensitivity to grid refinement is reported in
Appendix~\ref{app:grid}.

\paragraph{Grid sizes and computation.}
Baseline coarse grid: $\text{N\_W}=15$, $\text{N\_Z}=5$,
$\text{N\_X\_PREV}=3$, $\text{ASSET\_GRID}=7$, $\text{X\_GRID}=4$.
State-space size per period: $15 \times 5 \times 2 \times 3 \times 3 = 1350$
points (versus 294 in v3 with $\text{N\_W}=21$, $\text{N\_Z}=7$,
$\text{N\_ell}=2$).
Estimated wall time per regime: 2--4 hours single-threaded on server1.

%              Requires amsmath, booktabs, multirow, hyperref.

% ─────────────────────────────────────────────────────────────────────────────
\section{Model}
\label{sec:model}
% ─────────────────────────────────────────────────────────────────────────────

\subsection{Economic Environment}
\label{sec:model:environment}

The economy is populated by a unit mass of households who live for a
finite horizon of $T$ periods, corresponding to ages $25$ through $80$
(retirement at age $65$).
There are two locations, $A$ and $B$, each associated with its own housing
market.
At every date $t$, each household occupies exactly one location
$\ell_t \in \{A, B\}$.
Households have Epstein-Zin preferences that reduce to CRRA over
non-housing consumption in the baseline:
\begin{equation}
    U_t = \mathbb{E}_t\!\left[
        \sum_{s=t}^{T} \beta^{s-t}\,\pi_s\,
        u\!\left(\hat{c}_s\right)
    \right],
    \qquad
    u(c) = \frac{c^{1-\gamma}}{1-\gamma},
    \label{eq:utility}
\end{equation}
where $\beta$ is the time-preference factor, $\pi_s$ is the deterministic
survival probability at age $s$, and $\hat{c}_t$ is non-housing consumption
normalized by the house-price index.%
\footnote{We abstract from housing services in utility and focus on the
financial welfare gain from token portability. An extension with housing
services in utility is discussed in Appendix~\ref{app:robustness}.}

All quantities are normalized by the beginning-of-period house-price index
$P_t$, so the model unit is one house.

\subsection{State Space}
\label{sec:model:state}

The household's state vector in the baseline (\emph{v4 specification}) is
six-dimensional:
\begin{equation}
    \mathbf{s}_t \;=\;
    \bigl(t,\;w_t,\;z_t,\;\ell_t,\;x_{A,t-1},\;x_{B,t-1}\bigr),
    \label{eq:state}
\end{equation}
where:
\begin{itemize}[noitemsep]
    \item $w_t \in [w_{\min}, w_{\max}]$: beginning-of-period normalized
          financial wealth;
    \item $z_t \in [z_{\min}, z_{\max}]$: persistent component of labor income;
    \item $\ell_t \in \{A, B\}$: current location (discrete);
    \item $x_{A,t-1},\, x_{B,t-1} \geq 0$: token holdings at locations $A$ and
          $B$ carried forward from period $t-1$ (the \emph{lagged-position state}).
\end{itemize}
The lagged-position state $(x_{A,t-1}, x_{B,t-1})$ is the key novelty of the
v4 specification.  Tracking prior holdings creates an \emph{intertemporal
incentive}: a household at $\ell = A$ can pre-accumulate $x_{B,t-1} > 0$
at low incremental cost, reducing the buying cost it would face upon
relocation to $B$ in a future period.
Without this state variable, per-period buying costs are approximated as
zero (ignoring the pre-accumulation margin), which was the v3 Option~3
approximation.

At $t = 1$, all households enter with $x_{A,0} = x_{B,0} = 0$ (no prior
token positions).

\subsection{Stochastic Relocation Shock}
\label{sec:model:relocation}

Each period, the household receives a Bernoulli relocation shock:
\begin{equation}
    \mathbf{1}_{\text{reloc},t} \;\sim\; \operatorname{Bernoulli}\!\bigl(p_t^{\text{reloc}}\bigr),
    \qquad
    p_t^{\text{reloc}} = \begin{cases}
        p_{\text{work}} & \text{if}\; t \leq \tau_{\text{ret}}, \\
        p_{\text{ret}}  & \text{if}\; t  >  \tau_{\text{ret}},
    \end{cases}
    \label{eq:reloc_prob}
\end{equation}
where $\tau_{\text{ret}}$ is the retirement age.  In the baseline,
$p_{\text{work}} = 0.06$ (6\,\% per year for working-age households) and
$p_{\text{ret}} = 0.02$.  The working-age rate is anchored to
inter-MSA annual mobility rates in the PSID; see
Table~\ref{tab:calibration} and \citet{calibrationV3} for details.

When relocation occurs ($\mathbf{1}_{\text{reloc},t} = 1$), the household's
location toggles: $\ell_{t+1} = \ell_t' \neq \ell_t$.  The implications
for housing positions differ by regime (Section~\ref{sec:model:regimes}).

\subsection{Regime Taxonomy}
\label{sec:model:regimes}

We compare three regimes that differ in the admissible set of housing
positions:

\begin{table}[ht]
    \centering
    \caption{Regime Taxonomy}
    \label{tab:regimes}
    \renewcommand{\arraystretch}{1.20}
    \begin{tabular}{llll}
        \toprule
        Regime & Description & $x_{A,t}$ admissible & $x_{B,t}$ admissible \\
        \midrule
        E0     & Rent-only         & $0$ & $0$ \\
        E1\_2L & Binary ownership  & $\{0,1\}$ if $\ell=A$; $0$ if $\ell=B$
                                   & $\{0,1\}$ if $\ell=B$; $0$ if $\ell=A$ \\
        E2\_2L & Continuous tokens & $[0, \bar{x}]$ & $[0, \bar{x}]$ \\
        \bottomrule
    \end{tabular}
    \smallskip
    \begin{minipage}{\linewidth}
        \footnotesize
        Notes: $\bar{x}$ is the wealth-constrained maximum (Section~\ref{sec:model:budget}).
        In E1\_2L, the household can own \emph{only} at its current location; the
        non-occupied-location holding is zero by admissibility (reflecting the
        physical-residence requirement of traditional homeownership).
    \end{minipage}
\end{table}

\noindent\textbf{E0 (rent-only).}  The household never owns housing and pays
the full rent-to-price ratio $\rho$ each period.  This regime serves as the
no-ownership benchmark.

\noindent\textbf{E1\_2L (binary homeownership).}  The household owns either
zero or one unit of the property at its \emph{current} location.  It cannot
hold the non-occupied-location property.  This regime represents traditional
homeownership: a household at $A$ cannot hold a financial claim on a
location-$B$ property.  On relocation, the entire owned unit is sold
(at a discount, see Section~\ref{sec:model:txcost}), and the household
arrives at the new location with no prior position ($x_{t-1} = 0$ at
both locations).

\noindent\textbf{E2\_2L (fractional tokens, portable).}  The household can
hold any non-negative fractional amount at \emph{both} locations simultaneously.
Tokens are portable across relocations: moving from $A$ to $B$ does not
trigger a forced sale of location-$A$ tokens.  Transaction costs are charged
on the \emph{change} in holdings each period (see Section~\ref{sec:model:txcost}),
not as a lump-sum on relocation events.

The headline welfare object is $\operatorname{CEV}(\text{E2\_2L vs E1\_2L})$
--- the proportional consumption supplement that makes E1\_2L households
indifferent to E2\_2L access.

\subsection{Housing-Cost Rule}
\label{sec:model:kappa}

Net housing cost per period (as a fraction of house value) depends on the
household's regime and current holdings:
\begin{equation}
    \kappa(x_{A,t}, x_{B,t}, \ell_t) = \begin{cases}
        \rho                           & \text{(E0)} \\[4pt]
        \rho\,\mathbf{1}[x_{\ell,t} < 1]
        + m\,\mathbf{1}[x_{\ell,t} \geq 1]  & \text{(E1\_2L)} \\[4pt]
        \rho - x_{\ell,t}\,(\rho - m) & \text{(E2\_2L)}
    \end{cases}
    \label{eq:kappa}
\end{equation}
where $x_{\ell,t} \equiv x_{A,t}\,\mathbf{1}[\ell_t = A] + x_{B,t}\,\mathbf{1}[\ell_t = B]$
is the token holding at the occupied location, $\rho$ is the rent-to-price
ratio, and $m$ is the maintenance cost.

The parameter $\delta_{\text{own}} \equiv \rho - m$ is the per-unit annual
welfare gain from owning the occupied residence.  In the baseline calibration,
$\delta_{\text{own}} = 0.04$ (4\,\% per year).

\paragraph{The fixed-kappa convention.}
A key identifying restriction is that in E2\_2L, \emph{only} the
occupied-location token $x_{\ell,t}$ reduces net rent.  The non-occupied
token $x_{\ell',t}$ is a purely financial asset — it earns capital gains
but does not reduce rent (the household cannot simultaneously occupy and
rent out the same unit).  This convention, confirmed by falsification
tests in \citet{researchLogMay01}, ensures that cross-location holdings
are motivated exclusively by the financial return and the pre-accumulation
motive, not by a mechanical rent-saving subsidy.

\subsection{Period Budget Constraint}
\label{sec:model:budget}

Each period, the household allocates normalized wealth $w_t$ across
non-housing consumption $\hat{c}_t$, bond savings $b_t$, equity $s_t$,
and housing token purchases $(x_{A,t}, x_{B,t})$:
\begin{equation}
    \hat{c}_t + \kappa(x_{A,t}, x_{B,t}, \ell_t)
    + b_t + s_t + x_{A,t} + x_{B,t} + \tau_t^{\text{tx}} = w_t,
    \label{eq:budget}
\end{equation}
where $\tau_t^{\text{tx}}$ is the per-period transaction cost
(Section~\ref{sec:model:txcost}).  We require $\hat{c}_t > 0$,
$s_t \geq 0$, and $b_t \geq -\ell_{\max} \cdot x_{\ell,t}$ (LTV
mortgage constraint; $\ell_{\max} = 0$ in the baseline).

\subsection{Transaction Costs}
\label{sec:model:txcost}

\paragraph{E1\_2L forced-relocation sell.}
When an E1\_2L household at location $\ell$ relocates to $\ell' \neq \ell$,
it is forced to sell its housing unit.  The sell proceeds are discounted by
a transaction cost $\tau_{\text{sell}} \approx 6\,\%$ (NAR commission):
the sell \emph{factor} applied to location-$\ell$ holdings in the wealth
transition is $(1 - \tau_{\text{sell}})$.
Additionally, on arrival at the new location, the household pays a buying
cost $\tau_{\text{buy}} \approx 2.5\,\%$ of the purchase price when
acquiring a new unit.

\paragraph{E2\_2L per-period transaction cost.}
In the v4 specification, $x_{A,t-1}$ and $x_{B,t-1}$ are tracked as state
variables.  Each period, the household pays a transaction cost based on the
\emph{change} in holdings:
\begin{align}
    \Delta_A &= x_{A,t} - x_{A,t-1}, \qquad
    \Delta_B = x_{B,t} - x_{B,t-1}, \notag \\
    \tau_t^{\text{tx}} &= \tau_{\text{buy}} \,
        \bigl[\max(\Delta_A, 0) + \max(\Delta_B, 0)\bigr]
        + \tau_{\text{token}} \,
        \bigl[\max(-\Delta_A, 0) + \max(-\Delta_B, 0)\bigr],
    \label{eq:txcost}
\end{align}
where $\tau_{\text{buy}} = 0.025$ (per-unit buying cost) and
$\tau_{\text{token}} = 0.01$ (per-unit token-sell cost).
Tokens retained unchanged across periods incur no cost ($\Delta = 0$).

\paragraph{The pre-accumulation motive.}
Equation~\eqref{eq:txcost} creates an \emph{intertemporal incentive} absent
from earlier specifications.  A household at $\ell = A$ today faces an
expected future buying cost of $\tau_{\text{buy}} \cdot x_{B,t}^{*}$ upon
relocation to $B$, where $x_{B,t}^{*}$ is the desired holding at $B$.
By pre-accumulating $x_{B,t-1} > 0$ \emph{before} relocation, the household
reduces this future cost by $\tau_{\text{buy}} \cdot x_{B,t-1}$ per unit
pre-held.  The per-period expected saving is
$p_{\text{work}} \cdot \tau_{\text{buy}} \approx 0.15\,\%$ per unit of
$x_{B,t-1}$ held at $\ell = A$.

Under E1\_2L, cross-location pre-accumulation is inadmissible
(Table~\ref{tab:regimes}), so E1\_2L households cannot exploit this
margin.  This asymmetry is the structural source of the welfare gap.

\subsection{Asset Returns}
\label{sec:model:returns}

\paragraph{Housing returns.}
We decompose the single-location log housing return into a shared aggregate
factor and a location-specific idiosyncratic component:
\begin{align}
    \iota_{A,t+1} &= \sqrt{2}\,\sigma_\iota\,\xi_{A,t+1}, \notag \\
    \iota_{B,t+1} &= \rho_{AB}\,\iota_{A,t+1}
                     + \sqrt{1 - \rho_{AB}^2}\,\sqrt{2}\,\sigma_\iota\,\xi_{B,t+1},
                     \label{eq:chol} \\
    R_{A,t+1} &= \exp\!\bigl(\mu_h + \sqrt{2}\,\sigma_{\text{div}}\,\eta_{\text{div},t+1}
                               + \iota_{A,t+1}\bigr), \notag \\
    R_{B,t+1} &= \exp\!\bigl(\mu_h + \sqrt{2}\,\sigma_{\text{div}}\,\eta_{\text{div},t+1}
                               + \iota_{B,t+1}\bigr), \notag
\end{align}
where $\eta_{\text{div},t+1}$ is the \emph{shared} aggregate housing factor
(same realization for $A$ and $B$), $\xi_{A,t+1}$ and $\xi_{B,t+1}$ are
independent standard normal shocks, and $\sigma_{\text{div}}^2 + \sigma_\iota^2 =
\sigma_h^2$ (variance decomposition).
The Cholesky factorization in~\eqref{eq:chol} implies
$\operatorname{Corr}(\iota_{A,t+1}, \iota_{B,t+1}) = \rho_{AB}$.

The baseline $\rho_{AB} = 0.50$ is the midpoint of the Case-Shiller MSA-pair
idiosyncratic correlation range ($0.30$--$0.70$ for inter-regional US moves;
see \citealt{calibrationV3}).  Sensitivity analysis over
$\rho_{AB} \in \{0, 0.25, 0.50, 0.75, 0.95\}$ is reported in
Section~\ref{sec:results:sensitivity}.

\paragraph{Equity return.}
$R_{S,t+1} = \exp(\mu_s + \sqrt{2}\,\sigma_s\,\eta_{s,t+1})$,
where $\mu_s = \ln(R_f + \pi_s) - \frac{1}{2}\sigma_s^2$ and $\pi_s$
is the equity premium.

\paragraph{House-price normalization.}
The house-price factor $\mathit{hp}_{t+1} = \exp(g_h + \sqrt{2}\,\sigma_\xi\,\xi_{h,t+1})$
normalizes next-period wealth; all financial quantities are expressed in
units of current house prices.

\subsection{Income Process}
\label{sec:model:income}

Labor income follows the \citet{CoccoGomesMaenhout2005} specification.
The deterministic age profile is:
\begin{equation}
    f(\text{age}) = -2.17042 + 0.16818\left(\tfrac{\text{age}}{10}\right)
                   - 0.03230\left(\tfrac{\text{age}}{10}\right)^{\!2}
                   + 0.00200\left(\tfrac{\text{age}}{10}\right)^{\!3}.
    \label{eq:income_profile}
\end{equation}
The normalized persistent income component evolves as:
\begin{align}
    z_{t+1} &= z_t\,\exp\!\bigl(f(\text{age}_{t+1}) - f(\text{age}_t)
               + u_{t+1}\bigr)\big/\mathit{hp}_{t+1},
               & u_{t+1} &\sim \mathcal{N}(0, \sigma_u^2), \notag \\
    y_{t+1} &= z_{t+1}\,\exp(\varepsilon_{t+1}),
               & \varepsilon_{t+1} &\sim \mathcal{N}(0, \sigma_\varepsilon^2).
               \label{eq:income}
\end{align}
At retirement, $z_{t+1} = \lambda_{\text{ret}} z_t / \mathit{hp}_{t+1}$
(replacement rate $\lambda_{\text{ret}} = 0.65$); post-retirement income
is constant at $y_{t+1} = z_t / \mathit{hp}_{t+1}$.

\subsection{Wealth Transition}
\label{sec:model:wealth}

Beginning-of-next-period wealth is:
\begin{equation}
    w_{t+1} = \frac{b_t\,r_b + s_t\,R_{S,t+1}
                    + x_{A,t}\,R_{A,t+1}\,\phi_{A,t+1}
                    + x_{B,t}\,R_{B,t+1}\,\phi_{B,t+1}}
                   {\mathit{hp}_{t+1}}
              + y_{t+1},
    \label{eq:wealth_transition}
\end{equation}
where $r_b = R_f$ for $b_t \geq 0$ and $r_b = R_f + r_{\text{prem}}$
for $b_t < 0$, and $\phi_{\ell,t+1}$ is the sell factor:
\begin{equation}
    \phi_{\ell,t+1} = \begin{cases}
        1 - \tau_{\text{sell}} & \text{if E1\_2L and relocating from } \ell, \\
        1 & \text{otherwise.}
    \end{cases}
    \label{eq:sell_factor}
\end{equation}
Tokens in E2\_2L are portable: $\phi = 1$ always.  The forced-sell discount
$(1 - \tau_{\text{sell}})$ under E1\_2L is the first component of the
transaction-cost channel in the welfare decomposition
(Section~\ref{sec:model:welfare}).

\subsection{Bellman Equation}
\label{sec:model:bellman}

For regime $R \in \{\text{E0, E1\_2L, E2\_2L}\}$ and v4 state
$\mathbf{s}_t = (t, w_t, z_t, \ell_t, x_{A,t-1}, x_{B,t-1})$:
\begin{equation}
    V_t^R\!\bigl(w, z, \ell, x_{Ap}, x_{Bp}\bigr)
    = \max_{\hat{c},\, b,\, s,\, x_A,\, x_B \,\in\, \mathcal{A}_R}
    \left[
        u(\hat{c}) + \beta\,\pi_t\,
        \mathcal{E}_t^R\!\bigl(w, z, \ell, x_A, x_B\bigr)
    \right],
    \label{eq:bellman}
\end{equation}
where $\mathcal{A}_R$ is the regime-specific admissible set
(Table~\ref{tab:regimes} and budget constraint~\eqref{eq:budget}),
and the continuation value is:
\begin{multline}
    \mathcal{E}_t^R\!\bigl(w, z, \ell, x_A, x_B\bigr) \\
    = \sum_q \omega_q\,\mathit{hp}_q^{1-\gamma}
    \Bigl[
        (1 - p_t^{\text{reloc}})\,
        V_{t+1}^R\!\bigl(w_q^{\text{stay}},  z_q', \ell,  x_A^{\ast}, x_B^{\ast}\bigr)
        + p_t^{\text{reloc}}\,
        V_{t+1}^R\!\bigl(w_q^{\text{reloc}}, z_q', \ell', x_A^{\dagger}, x_B^{\dagger}\bigr)
    \Bigr],
    \label{eq:ev}
\end{multline}
where $(\omega_q, \mathit{hp}_q, w_q^{\text{stay}}, w_q^{\text{reloc}}, z_q')$
are Gauss-Hermite quadrature weights and realizations (7D; $3^7 = 2187$
integration points), and the next-period $x$-state updates are:
\begin{align}
    (x_A^{\ast}, x_B^{\ast})   &= (x_A, x_B),
        \quad &\text{(stay; E0, E2\_2L, E1\_2L)} \label{eq:xprev_stay} \\
    (x_A^{\dagger}, x_B^{\dagger}) &= (x_A, x_B),
        \quad &\text{(reloc; E0, E2\_2L — portable)} \label{eq:xprev_reloc_e2} \\
    (x_A^{\dagger}, x_B^{\dagger}) &= (0, 0),
        \quad &\text{(reloc; E1\_2L — forced sale, fresh start)}
        \label{eq:xprev_reloc_e1}
\end{align}
Equation~\eqref{eq:xprev_reloc_e1} captures the key asymmetry: under
E1\_2L, relocation resets the lagged-position state to zero (the household
must buy fresh at the new location, paying $\tau_{\text{buy}}$ on the full
amount), while under E2\_2L, equation~\eqref{eq:xprev_reloc_e2} shows that
tokens carry forward unchanged.

The terminal condition is $V_{T+1}^R(w, \cdot) = u(w)$
(bequest-free, all wealth consumed at death).

\subsection{Welfare Measure and Decomposition}
\label{sec:model:welfare}

\paragraph{Certainty-equivalent consumption (CEV).}
The welfare comparison between regimes $R_a$ and $R_b$ is measured by the
proportional consumption supplement $\lambda$ such that:
\begin{equation}
    V_1^{R_b}\!\bigl((1 + \lambda)\,w^{\ast}\bigr) = V_1^{R_a}\!\bigl(w^{\ast}\bigr),
    \label{eq:cev_def}
\end{equation}
evaluated at the representative initial state $w^{\ast}$ (midpoint of the
wealth grid, $z = z_{\text{mid}}$, $\ell = A$, $x_{Ap} = x_{Bp} = 0$).
Under CRRA:
\begin{equation}
    1 + \lambda = \left(\frac{V_1^{R_a}}{V_1^{R_b}}\right)^{\!\!1/(1-\gamma)}.
    \label{eq:cev_crra}
\end{equation}

\paragraph{Channel decomposition.}
The headline CEV $\operatorname{CEV}(\text{E2\_2L vs E1\_2L})$ decomposes into
two channels via the counterfactual regime E1\_2L$_{\text{NOTX}}$ (E1\_2L
with $\tau_{\text{sell}} = 0$):
\begin{align}
    \underbrace{\operatorname{CEV}(\text{E2\_2L vs E1\_2L})}_{\text{total}}
    &=
    \underbrace{\operatorname{CEV}(\text{E1\_2L}_{\text{NOTX}} \text{ vs E1\_2L})}_{\text{avoided-tx channel}}
    +
    \underbrace{\operatorname{CEV}(\text{E2\_2L vs E1\_2L}_{\text{NOTX}})}_{\text{maintained-hedge channel}}
    + \xi_{\text{cross}},
    \label{eq:decomp}
\end{align}
where $\xi_{\text{cross}}$ is the cross-term, which was empirically near zero
($< 0.1\%$) in the v3 specification and is expected to remain small in v4.
We report $\xi_{\text{cross}}$ explicitly to verify separability.

Table~\ref{tab:cev_decomp} reports the decomposition with numerical results
to be filled after server1 baseline runs are complete.

\begin{table}[ht]
    \centering
    \caption{CEV Decomposition: Tokenization Welfare Channels}
    \label{tab:cev_decomp}
    \renewcommand{\arraystretch}{1.20}
    \begin{tabular}{lcc}
        \toprule
        Channel & CEV (\%) & Share (\%) \\
        \midrule
        Total tokenization gain:
            $\operatorname{CEV}(\text{E2\_2L vs E1\_2L})$
            & {\sc [x.xx]} & 100 \\
        \quad Avoided-tx channel:
            $\operatorname{CEV}(\text{E1\_2L}_{\text{NOTX}} \text{ vs E1\_2L})$
            & {\sc [x.xx]} & {\sc [xx]} \\
        \quad Maintained-hedge channel:
            $\operatorname{CEV}(\text{E2\_2L vs E1\_2L}_{\text{NOTX}})$
            & {\sc [x.xx]} & {\sc [xx]} \\
        \quad Cross-term $\xi_{\text{cross}}$
            & {\sc [x.xx]} & {\sc [xx]} \\
        \midrule
        Option 1 incremental gain:
            $\operatorname{CEV}(\text{E2\_2L v4 vs E2\_2L v3})$
            & {\sc [x.xx]} & --- \\
        \bottomrule
    \end{tabular}
    \smallskip
    \begin{minipage}{\linewidth}
        \footnotesize
        Notes: CEVs evaluated at representative initial state
        $(w^{\ast}, z_{\text{mid}}, \ell = A, x_{Ap} = x_{Bp} = 0)$
        using the v4 solver (6D state; $\text{N\_W}=15$,
        $\text{N\_Z}=5$, $\text{N\_X\_PREV}=3$, $\text{GH\_NODES}=3$).
        {\sc [x.xx]}: placeholder to be filled from server1 runs.
        Baseline calibration: Table~\ref{tab:calibration}.
    \end{minipage}
\end{table}

\paragraph{Hypothesis tests.}
We pre-register three empirical hypotheses (H1--H3) that must hold for
the option-1 mechanism to support the RFS-level contribution claim:
\begin{description}[noitemsep]
    \item[H1] $\overline{x}_{B,1}^{\text{E2}} > 0$ at $\ell = A$
              (cross-location pre-accumulation activates).
    \item[H2] $\operatorname{CEV}(\text{E2\_2L v4 vs E1\_2L v4}) > 4.255\,\%$
              (beats the Option~3 approximation baseline).
    \item[H3] $\operatorname{CEV}(\text{E2\_2L v4 vs E2\_2L v3}) \in [0.5, 1.5]\,\%$
              (incremental hedge value from proper state extension).
\end{description}
If all three hold: the mechanism is RFS-credible and Phase~2
(calibration, sensitivity, manuscript) proceeds.  If any fails, the
paper targets Real Estate Economics or Journal of Housing Economics
with the Option~3 result ($+4.26\,\%$).

\subsection{Numerical Solution}
\label{sec:model:numerical}

We solve the Bellman equation~\eqref{eq:bellman} by backward induction
(value function iteration) using the v4 Julia solver
(\texttt{src/vfi\_solver\_v4.jl}, 954 LOC).
The value function is approximated on a discrete 6D grid; continuation
values are computed by 7D Gauss-Hermite quadrature ($3^7 = 2187$ points).

\paragraph{Interpolation.}
The continuation-value lookup interpolates in $(w, z)$ bilinearly for
a given $(\ell, x_{Ap}, x_{Bp})$ grid cell.  The lagged-position state
$(x_{Ap}, x_{Bp})$ uses a coarse grid with $\text{N\_X\_PREV} = 3$
points ($\{0, 1, 2\}$ in the baseline), sufficient for the first-cut
test of H1.  Sensitivity to grid refinement is reported in
Appendix~\ref{app:grid}.

\paragraph{Grid sizes and computation.}
Baseline coarse grid: $\text{N\_W}=15$, $\text{N\_Z}=5$,
$\text{N\_X\_PREV}=3$, $\text{ASSET\_GRID}=7$, $\text{X\_GRID}=4$.
State-space size per period: $15 \times 5 \times 2 \times 3 \times 3 = 1350$
points (versus 294 in v3 with $\text{N\_W}=21$, $\text{N\_Z}=7$,
$\text{N\_ell}=2$).
Estimated wall time per regime: 2--4 hours single-threaded on server1.

% ─────────────────────────────────────────────────────────────────────────────

\clearpage

% ─────────────────────────────────────────────────────────────────────────────
% s3_calibration.tex — Calibration section
% Paper: "Tokenized Housing and Lifecycle Portfolio Choice:
%         A Decoupling of Location from Housing Exposure"
%
% Source: docs/calibration_v3.md (cloud agent fire 5)
%         docs/welfare_decomp_v4.md (channel decomposition plan)
%         paper/outline_v4.md (section structure)
%
% Status: complete draft; one parameter (rho_AB sensitivity) awaits
%         server1 baseline results before final anchor confirmation (H2' gate).
%
% Compilation: included from main.tex via % s3_calibration.tex — Calibration section
% Paper: "Tokenized Housing and Lifecycle Portfolio Choice:
%         A Decoupling of Location from Housing Exposure"
%
% Source: docs/calibration_v3.md (cloud agent fire 5)
%         docs/welfare_decomp_v4.md (channel decomposition plan)
%         paper/outline_v4.md (section structure)
%
% Status: complete draft; one parameter (rho_AB sensitivity) awaits
%         server1 baseline results before final anchor confirmation (H2' gate).
%
% Compilation: included from main.tex via % s3_calibration.tex — Calibration section
% Paper: "Tokenized Housing and Lifecycle Portfolio Choice:
%         A Decoupling of Location from Housing Exposure"
%
% Source: docs/calibration_v3.md (cloud agent fire 5)
%         docs/welfare_decomp_v4.md (channel decomposition plan)
%         paper/outline_v4.md (section structure)
%
% Status: complete draft; one parameter (rho_AB sensitivity) awaits
%         server1 baseline results before final anchor confirmation (H2' gate).
%
% Compilation: included from main.tex via \input{sections/s3_calibration}
%              Requires amsmath, booktabs, multirow, siunitx (for tabular).

% ─────────────────────────────────────────────────────────────────────────────
\section{Calibration}
\label{sec:calibration}
% ─────────────────────────────────────────────────────────────────────────────

Table~\ref{tab:parameters} summarises the baseline parameter values.
The calibration follows the household-finance lifecycle tradition
(Cocco, Gomes, and Maenhout, 2005; Yao and Zhang, 2005), extended
with three v4-specific objects: the relocation probability $p_{\text{work}}$,
the round-trip transaction-cost pair $(\tau_{\text{sell}}, \tau_{\text{buy}})$,
and the cross-location idiosyncratic correlation $\rho_{AB}$.
We discuss each group in turn.

% ── Table 1: Parameter summary ───────────────────────────────────────────────
\begin{table}[t]
\centering
\caption{Baseline Parameters}
\label{tab:parameters}
\small
\begin{tabular}{llcl}
\toprule
Parameter & Symbol & Value & Source \\
\midrule
\multicolumn{4}{l}{\textit{Preferences and lifecycle}} \\
Risk aversion         & $\gamma$         & 5.0  & CGM (2005) \\
Time discount         & $\beta$          & 0.96 & CGM (2005) \\
Entry age / retire / terminal & --- & 25 / 65 / 80 & Standard \\
\addlinespace
\multicolumn{4}{l}{\textit{Financial assets}} \\
Risk-free gross return     & $R_f$           & 1.02  & CGM (2005) \\
Equity risk premium        & $\mu_S - r_f$   & 4\%   & CGM (2005) \\
Stock return volatility    & $\sigma_S$      & 15.7\% & CGM (2005) \\
\addlinespace
\multicolumn{4}{l}{\textit{Income process (CGM 2005)}} \\
Permanent shock variance   & $\sigma_u^2$    & 0.0106 & CGM Table~I \\
Transitory shock variance  & $\sigma_\varepsilon^2$ & 0.0738 & CGM Table~I \\
Retirement replacement     & $\lambda$       & 0.65  & CGM (2005) \\
\addlinespace
\multicolumn{4}{l}{\textit{Housing return (Cocco 2005; Yao-Zhang 2005)}} \\
Expected log price growth  & $g_h$           & 1.6\%  & Cocco (2005) \\
Total housing return std   & $\sigma_h$      & 11.5\% & Cocco (2005) \\
Aggregate factor std       & $\sigma_{\text{div}}$ & 10.0\% & Case-Shiller / NAREIT \\
Idiosyncratic std (derived) & $\sigma_\iota$ & 5.7\%  & $\sqrt{\sigma_h^2 - \sigma_{\text{div}}^2}$ \\
Rent-to-price ratio        & $\varrho$       & 5.0\%  & Yao and Zhang (2005) \\
Maintenance-to-price ratio & $m$             & 1.0\%  & Cocco (2005) \\
Rent-saving wedge (derived) & $\delta_{\text{own}} = \varrho - m$ & 4.0\% & --- \\
\addlinespace
\multicolumn{4}{l}{\textit{Mobility (PSID-anchored)}} \\
Relocation prob, working age & $p_{\text{work}}$ & 6.0\% & PSID inter-MSA \\
Relocation prob, retired     & $p_{\text{ret}}$  & 2.0\% & PSID retired \\
\addlinespace
\multicolumn{4}{l}{\textit{Transaction costs}} \\
Traditional sell cost  & $\tau_{\text{sell}}$ & 6.0\% & NAR (2023) \\
Traditional buy cost   & $\tau_{\text{buy}}$  & 2.5\% & CFPB / ALTA \\
Token transfer cost    & $\tau_{\text{token}}$& 1.0\% & RealT / Lofty surveys \\
\addlinespace
\multicolumn{4}{l}{\textit{Cross-location correlation}} \\
Idiosyncratic corr., MSA pair & $\rho_{AB}$ & 0.50 & Case-Shiller (2000--2023) \\
\addlinespace
\multicolumn{4}{l}{\textit{Mortgage (baseline off)}} \\
LTV constraint & $\overline{\text{ltv}}$ & 0.0 & Baseline no-mortgage \\
\bottomrule
\end{tabular}
\end{table}

\subsection{Preferences and Income Process}
\label{sec:calibration:prefs}

Preferences and the income process are taken directly from
Cocco, Gomes, and Maenhout (2005, CGM hereafter).
Risk aversion $\gamma = 5$ and the discount factor $\beta = 0.96$ are
their baseline values.
The labor income age profile follows the CGM cubic polynomial fit
to PSID data:
\begin{equation}
    f(\text{age}) \;=\;
    -2.170 + 0.168\,(\text{age}/10)
    - 0.032\,(\text{age}/10)^2
    + 0.002\,(\text{age}/10)^3.
    \label{eq:income_profile}
\end{equation}
Permanent and transitory income shocks have variances $\sigma_u^2 = 0.0106$
and $\sigma_\varepsilon^2 = 0.0738$ (CGM Table~I).
The retirement replacement rate $\lambda = 0.65$ matches the median
Social Security replacement rate at median lifetime earnings.

Financial asset parameters follow the same source:
the equity premium is $4\%$, stock return volatility $15.7\%$,
and the risk-free rate $2\%$ real.

\subsection{Housing Return Decomposition}
\label{sec:calibration:housing}

Single-location housing return volatility $\sigma_h = 11.5\%$ is from
Cocco (2005).
We decompose the single-location log return into a shared aggregate factor
$\eta_{\text{div}}$ and an idiosyncratic component $\iota_\ell$:
\begin{equation}
    \log R_\ell \;=\; \mu_h + \eta_{\text{div}} + \iota_\ell,
    \qquad
    \eta_{\text{div}} \sim \mathcal{N}(0,\sigma_{\text{div}}^2),
    \quad
    \iota_\ell \sim \mathcal{N}(0,\sigma_\iota^2),
    \label{eq:return_decomp}
\end{equation}
with $\sigma_h^2 = \sigma_{\text{div}}^2 + \sigma_\iota^2$.
We set $\sigma_{\text{div}} = 10\%$ so that the aggregate factor accounts
for approximately $75\%$ of total single-location housing variance
($\sigma_{\text{div}}^2 / \sigma_h^2 \approx 0.756$), consistent with
the common factor dominating MSA-level variance in Case-Shiller data
and with Piazzesi and Schneider~(2016).
The idiosyncratic component $\sigma_\iota = \sqrt{0.115^2 - 0.10^2} \approx 5.7\%$
is calibrated residually.

The idiosyncratic components at the two locations follow a bivariate normal
with correlation $\rho_{AB}$ (Section~\ref{sec:calibration:rhoAB}).
The Cholesky factorisation is:
\begin{align}
    \iota_A &= \sigma_\iota\,\xi_A, \label{eq:chol_A} \\
    \iota_B &= \rho_{AB}\,\iota_A + \sqrt{1 - \rho_{AB}^2}\,\sigma_\iota\,\xi_B,
    \label{eq:chol_B}
\end{align}
with $\xi_A, \xi_B \sim \mathcal{N}(0,1)$ independent.

The rent-to-price ratio $\varrho = 5\%$ follows Yao and Zhang~(2005).
Maintenance costs $m = 1\%$ follow Cocco~(2005).
The rent-saving wedge $\delta_{\text{own}} = \varrho - m = 4\%$ governs
the per-period benefit of occupying a fractional token position
($\kappa$ rule, equation~(\ref{eq:kappa}) in Section~\ref{sec:model:kappa}).

\subsection{Mobility Rate}
\label{sec:calibration:mobility}

We anchor the relocation probability to PSID inter-MSA mobility data.
Total annual relocation rates in the PSID average $7$--$12\%$ for
working-age households but include short-distance, within-MSA moves that
do not change the relevant housing market.
Long-distance, inter-MSA moves account for approximately $40$--$50\%$
of total PSID mobility, yielding an inter-MSA rate of $5$--$7\%$ annually
for prime-age workers.
We set $p_{\text{work}} = 6\%$ (PSID midpoint) for working ages
$25$--$64$ and $p_{\text{ret}} = 2\%$ for ages $65$--$80$.

This choice is consistent with Bagliano, Fugazza, and Nicodano~(2014, RFS),
who calibrate career-driven mobility shocks at $5$--$8\%$ annually, and
broadly in the range of Yao and Zhang~(2005) ($4\%$) and Saks and
Wozniak~(2011, JLE) ($2$--$4\%$ inter-state, implying $4$--$8\%$ inter-MSA).

The sensitivity grid (Table~\ref{tab:sensitivity_grid} below)
spans $p_{\text{work}} \in \{0, 0.02, 0.06, 0.12\}$.
At $p_{\text{work}} = 0$, the cross-location hedge channel must vanish by
construction (no relocation, no cost-saving motive); this is a pre-registered
falsification test (Referee Round~4, test~r).

\subsection{Transaction Costs}
\label{sec:calibration:txcost}

\paragraph{Traditional housing ($E1\_2L$).}
We set the seller commission and closing costs at $\tau_{\text{sell}} = 6\%$,
consistent with the NAR~(2023) aggregate for a typical single-family
home sale ($5$--$6\%$ agent commission plus $0.5$--$1\%$ closing costs).
The buyer cost is $\tau_{\text{buy}} = 2.5\%$,
covering origination fees ($0.75\%$), title insurance and escrow
($0.75\%$), appraisal and inspection ($0.25\%$), and prepaid items
($0.25\%$), consistent with CFPB and ALTA surveys.
The implied round-trip cost for an $E1\_2L$ relocation is
$\tau_{\text{sell}} + \tau_{\text{buy}} = 8.5\%$,
within the $5$--$10\%$ range used in the structural housing literature
(Diaz and Luengo-Prado, 2008; Head, Lloyd-Ellis, and Sun, 2014).

\paragraph{Tokenized housing ($E2\_2L$).}
Platform surveys of RealT, Lofty, and DigiShares~(2024) report trading
fees in the range $0.1$--$0.5\%$ plus blockchain gas and settlement
of $0.05$--$0.3\%$ on Layer-2 networks.
We set $\tau_{\text{token}} = 1\%$ to include a regulatory-compliance
margin of $0.5\%$ (reflecting SEC Regulation~D and AML overhead).
The structural distinction from $E1\_2L$ is that $\tau_{\text{token}} \ll \tau_{\text{sell}} + \tau_{\text{buy}}$:
tokens are portable across relocations, so the $8.5\%$ round-trip is avoided
for every unit of cross-location position retained.

\paragraph{Expected hedge premium.}
The per-period value of pre-positioning one unit of $x_B$ while at
location $A$ is
\begin{equation}
    \text{hedge premium per unit}
    \;=\;
    p_{\text{work}} \times \tau_{\text{buy}}
    \;\approx\;
    0.06 \times 0.025
    \;=\; 0.15\%\ \text{per period}.
    \label{eq:hedge_premium}
\end{equation}
This is the incremental saving on buying costs that accrues when the
household has pre-positioned $x_B > 0$ prior to relocation.
Accumulated over a working-life horizon, equation~(\ref{eq:hedge_premium})
implies a lifetime CEV contribution of $1$--$2\%$ from the hedge
channel alone, with the continuous-$x$ rent-saving channel adding a further
$3$--$4\%$ (see Section~\ref{sec:results:decomp}).

\subsection{Cross-Location House-Price Correlation}
\label{sec:calibration:rhoAB}

The parameter $\rho_{AB}$ governs the idiosyncratic correlation between
location $A$ and location $B$ returns.
Using the decomposition~(\ref{eq:return_decomp}),
the observed return correlation between two MSAs is:
\begin{equation}
    \operatorname{Corr}(R_A, R_B)
    \;\approx\;
    \frac{\sigma_{\text{div}}^2}{\sigma_h^2}
    \;+\;
    \frac{\sigma_\iota^2}{\sigma_h^2}\,\rho_{AB}
    \;\approx\; 0.756 + 0.244\,\rho_{AB}.
    \label{eq:rhoAB_mapping}
\end{equation}
We set $\rho_{AB} = 0.50$, which maps to an observed return correlation of
$\approx 0.88$, consistent with pairs of geographically proximate
US metros in the Case-Shiller 20-city panel (2000--2023).
This is the midpoint of the $0.30$--$0.70$ range that covers inter-regional
moves (different metro areas): at one extreme, Chicago--Phoenix
($\hat{\rho}_{AB} \approx 0.35$) and at the other,
NY--Boston ($\hat{\rho}_{AB} \approx 0.60$).

Equation~(\ref{eq:rhoAB_mapping}) implies that the hedge value of
$x_B$ at location $A$ derives from bearing the idiosyncratic component
$\iota_B$ independently of $\iota_A$.
As $\rho_{AB} \to 1$, locations $A$ and $B$ become nearly identical
markets and $x_B$ is a near-substitute for $x_A$; the hedge channel
collapses.
The Round-4 referee pre-registered test~(m) requires $\operatorname{CEV}(E2\_2L \text{ vs } E1\_2L)$
to decrease significantly at $\rho_{AB} = 0.95$.

\subsection{Sensitivity Grid}
\label{sec:calibration:sensitivity}

Table~\ref{tab:sensitivity_grid} lists the pre-registered sensitivity
sweeps queued for execution after the v4 baseline results are available.
The three P1-priority sweeps cover the three parameters most directly
linked to the welfare decomposition channels.

\begin{table}[t]
\centering
\caption{Pre-Registered Sensitivity Sweeps}
\label{tab:sensitivity_grid}
\small
\begin{tabular}{llll}
\toprule
Parameter & Sweep values & Channel & Script \\
\midrule
$\rho_{AB}$          & $\{0,\, 0.25,\, 0.50,\, 0.75,\, 0.95\}$  & Hedge (maintained) & \texttt{sweep\_rhoAB.sh} \\
$p_{\text{work}}$    & $\{0,\, 0.02,\, 0.06,\, 0.12\}$           & Hedge (frequency)  & \texttt{sweep\_prelocate.sh} \\
$\tau_{\text{buy}}$  & $\{0,\, 0.025,\, 0.04,\, 0.06\}$          & Avoided tx-cost    & \texttt{sweep\_txcost.sh} \\
\midrule
\multicolumn{4}{l}{\textit{Phase 2 sweeps (after H2' approval)}} \\
$\gamma$             & $\{3,\, 5,\, 8\}$                         & Risk channel       & --- \\
$\delta_{\text{own}}$& $\{0.02,\, 0.04,\, 0.06\}$               & Rent-saving        & --- \\
$\sigma_{\text{div}}$& $\{0.08,\, 0.10,\, 0.12\}$               & Idio.\ risk share  & --- \\
\bottomrule
\end{tabular}
\end{table}

At $p_{\text{work}} = 0$ and at $\tau_{\text{buy}} = 0$,
the hedge premium~(\ref{eq:hedge_premium}) is exactly zero; any
remaining welfare gap between $E2\_2L$ and $E1\_2L$ in these cells
reflects the continuous-$x$ rent-saving channel only and serves
as a falsification bound.
At $\rho_{AB} = 0.95$, the maintained-hedge channel must also collapse
because $x_B$ and $x_A$ are near-substitutes.
These three boundary conditions define sharp predictions that distinguish
the two decomposed channels (Section~\ref{sec:results:decomp}).

%              Requires amsmath, booktabs, multirow, siunitx (for tabular).

% ─────────────────────────────────────────────────────────────────────────────
\section{Calibration}
\label{sec:calibration}
% ─────────────────────────────────────────────────────────────────────────────

Table~\ref{tab:parameters} summarises the baseline parameter values.
The calibration follows the household-finance lifecycle tradition
(Cocco, Gomes, and Maenhout, 2005; Yao and Zhang, 2005), extended
with three v4-specific objects: the relocation probability $p_{\text{work}}$,
the round-trip transaction-cost pair $(\tau_{\text{sell}}, \tau_{\text{buy}})$,
and the cross-location idiosyncratic correlation $\rho_{AB}$.
We discuss each group in turn.

% ── Table 1: Parameter summary ───────────────────────────────────────────────
\begin{table}[t]
\centering
\caption{Baseline Parameters}
\label{tab:parameters}
\small
\begin{tabular}{llcl}
\toprule
Parameter & Symbol & Value & Source \\
\midrule
\multicolumn{4}{l}{\textit{Preferences and lifecycle}} \\
Risk aversion         & $\gamma$         & 5.0  & CGM (2005) \\
Time discount         & $\beta$          & 0.96 & CGM (2005) \\
Entry age / retire / terminal & --- & 25 / 65 / 80 & Standard \\
\addlinespace
\multicolumn{4}{l}{\textit{Financial assets}} \\
Risk-free gross return     & $R_f$           & 1.02  & CGM (2005) \\
Equity risk premium        & $\mu_S - r_f$   & 4\%   & CGM (2005) \\
Stock return volatility    & $\sigma_S$      & 15.7\% & CGM (2005) \\
\addlinespace
\multicolumn{4}{l}{\textit{Income process (CGM 2005)}} \\
Permanent shock variance   & $\sigma_u^2$    & 0.0106 & CGM Table~I \\
Transitory shock variance  & $\sigma_\varepsilon^2$ & 0.0738 & CGM Table~I \\
Retirement replacement     & $\lambda$       & 0.65  & CGM (2005) \\
\addlinespace
\multicolumn{4}{l}{\textit{Housing return (Cocco 2005; Yao-Zhang 2005)}} \\
Expected log price growth  & $g_h$           & 1.6\%  & Cocco (2005) \\
Total housing return std   & $\sigma_h$      & 11.5\% & Cocco (2005) \\
Aggregate factor std       & $\sigma_{\text{div}}$ & 10.0\% & Case-Shiller / NAREIT \\
Idiosyncratic std (derived) & $\sigma_\iota$ & 5.7\%  & $\sqrt{\sigma_h^2 - \sigma_{\text{div}}^2}$ \\
Rent-to-price ratio        & $\varrho$       & 5.0\%  & Yao and Zhang (2005) \\
Maintenance-to-price ratio & $m$             & 1.0\%  & Cocco (2005) \\
Rent-saving wedge (derived) & $\delta_{\text{own}} = \varrho - m$ & 4.0\% & --- \\
\addlinespace
\multicolumn{4}{l}{\textit{Mobility (PSID-anchored)}} \\
Relocation prob, working age & $p_{\text{work}}$ & 6.0\% & PSID inter-MSA \\
Relocation prob, retired     & $p_{\text{ret}}$  & 2.0\% & PSID retired \\
\addlinespace
\multicolumn{4}{l}{\textit{Transaction costs}} \\
Traditional sell cost  & $\tau_{\text{sell}}$ & 6.0\% & NAR (2023) \\
Traditional buy cost   & $\tau_{\text{buy}}$  & 2.5\% & CFPB / ALTA \\
Token transfer cost    & $\tau_{\text{token}}$& 1.0\% & RealT / Lofty surveys \\
\addlinespace
\multicolumn{4}{l}{\textit{Cross-location correlation}} \\
Idiosyncratic corr., MSA pair & $\rho_{AB}$ & 0.50 & Case-Shiller (2000--2023) \\
\addlinespace
\multicolumn{4}{l}{\textit{Mortgage (baseline off)}} \\
LTV constraint & $\overline{\text{ltv}}$ & 0.0 & Baseline no-mortgage \\
\bottomrule
\end{tabular}
\end{table}

\subsection{Preferences and Income Process}
\label{sec:calibration:prefs}

Preferences and the income process are taken directly from
Cocco, Gomes, and Maenhout (2005, CGM hereafter).
Risk aversion $\gamma = 5$ and the discount factor $\beta = 0.96$ are
their baseline values.
The labor income age profile follows the CGM cubic polynomial fit
to PSID data:
\begin{equation}
    f(\text{age}) \;=\;
    -2.170 + 0.168\,(\text{age}/10)
    - 0.032\,(\text{age}/10)^2
    + 0.002\,(\text{age}/10)^3.
    \label{eq:income_profile}
\end{equation}
Permanent and transitory income shocks have variances $\sigma_u^2 = 0.0106$
and $\sigma_\varepsilon^2 = 0.0738$ (CGM Table~I).
The retirement replacement rate $\lambda = 0.65$ matches the median
Social Security replacement rate at median lifetime earnings.

Financial asset parameters follow the same source:
the equity premium is $4\%$, stock return volatility $15.7\%$,
and the risk-free rate $2\%$ real.

\subsection{Housing Return Decomposition}
\label{sec:calibration:housing}

Single-location housing return volatility $\sigma_h = 11.5\%$ is from
Cocco (2005).
We decompose the single-location log return into a shared aggregate factor
$\eta_{\text{div}}$ and an idiosyncratic component $\iota_\ell$:
\begin{equation}
    \log R_\ell \;=\; \mu_h + \eta_{\text{div}} + \iota_\ell,
    \qquad
    \eta_{\text{div}} \sim \mathcal{N}(0,\sigma_{\text{div}}^2),
    \quad
    \iota_\ell \sim \mathcal{N}(0,\sigma_\iota^2),
    \label{eq:return_decomp}
\end{equation}
with $\sigma_h^2 = \sigma_{\text{div}}^2 + \sigma_\iota^2$.
We set $\sigma_{\text{div}} = 10\%$ so that the aggregate factor accounts
for approximately $75\%$ of total single-location housing variance
($\sigma_{\text{div}}^2 / \sigma_h^2 \approx 0.756$), consistent with
the common factor dominating MSA-level variance in Case-Shiller data
and with Piazzesi and Schneider~(2016).
The idiosyncratic component $\sigma_\iota = \sqrt{0.115^2 - 0.10^2} \approx 5.7\%$
is calibrated residually.

The idiosyncratic components at the two locations follow a bivariate normal
with correlation $\rho_{AB}$ (Section~\ref{sec:calibration:rhoAB}).
The Cholesky factorisation is:
\begin{align}
    \iota_A &= \sigma_\iota\,\xi_A, \label{eq:chol_A} \\
    \iota_B &= \rho_{AB}\,\iota_A + \sqrt{1 - \rho_{AB}^2}\,\sigma_\iota\,\xi_B,
    \label{eq:chol_B}
\end{align}
with $\xi_A, \xi_B \sim \mathcal{N}(0,1)$ independent.

The rent-to-price ratio $\varrho = 5\%$ follows Yao and Zhang~(2005).
Maintenance costs $m = 1\%$ follow Cocco~(2005).
The rent-saving wedge $\delta_{\text{own}} = \varrho - m = 4\%$ governs
the per-period benefit of occupying a fractional token position
($\kappa$ rule, equation~(\ref{eq:kappa}) in Section~\ref{sec:model:kappa}).

\subsection{Mobility Rate}
\label{sec:calibration:mobility}

We anchor the relocation probability to PSID inter-MSA mobility data.
Total annual relocation rates in the PSID average $7$--$12\%$ for
working-age households but include short-distance, within-MSA moves that
do not change the relevant housing market.
Long-distance, inter-MSA moves account for approximately $40$--$50\%$
of total PSID mobility, yielding an inter-MSA rate of $5$--$7\%$ annually
for prime-age workers.
We set $p_{\text{work}} = 6\%$ (PSID midpoint) for working ages
$25$--$64$ and $p_{\text{ret}} = 2\%$ for ages $65$--$80$.

This choice is consistent with Bagliano, Fugazza, and Nicodano~(2014, RFS),
who calibrate career-driven mobility shocks at $5$--$8\%$ annually, and
broadly in the range of Yao and Zhang~(2005) ($4\%$) and Saks and
Wozniak~(2011, JLE) ($2$--$4\%$ inter-state, implying $4$--$8\%$ inter-MSA).

The sensitivity grid (Table~\ref{tab:sensitivity_grid} below)
spans $p_{\text{work}} \in \{0, 0.02, 0.06, 0.12\}$.
At $p_{\text{work}} = 0$, the cross-location hedge channel must vanish by
construction (no relocation, no cost-saving motive); this is a pre-registered
falsification test (Referee Round~4, test~r).

\subsection{Transaction Costs}
\label{sec:calibration:txcost}

\paragraph{Traditional housing ($E1\_2L$).}
We set the seller commission and closing costs at $\tau_{\text{sell}} = 6\%$,
consistent with the NAR~(2023) aggregate for a typical single-family
home sale ($5$--$6\%$ agent commission plus $0.5$--$1\%$ closing costs).
The buyer cost is $\tau_{\text{buy}} = 2.5\%$,
covering origination fees ($0.75\%$), title insurance and escrow
($0.75\%$), appraisal and inspection ($0.25\%$), and prepaid items
($0.25\%$), consistent with CFPB and ALTA surveys.
The implied round-trip cost for an $E1\_2L$ relocation is
$\tau_{\text{sell}} + \tau_{\text{buy}} = 8.5\%$,
within the $5$--$10\%$ range used in the structural housing literature
(Diaz and Luengo-Prado, 2008; Head, Lloyd-Ellis, and Sun, 2014).

\paragraph{Tokenized housing ($E2\_2L$).}
Platform surveys of RealT, Lofty, and DigiShares~(2024) report trading
fees in the range $0.1$--$0.5\%$ plus blockchain gas and settlement
of $0.05$--$0.3\%$ on Layer-2 networks.
We set $\tau_{\text{token}} = 1\%$ to include a regulatory-compliance
margin of $0.5\%$ (reflecting SEC Regulation~D and AML overhead).
The structural distinction from $E1\_2L$ is that $\tau_{\text{token}} \ll \tau_{\text{sell}} + \tau_{\text{buy}}$:
tokens are portable across relocations, so the $8.5\%$ round-trip is avoided
for every unit of cross-location position retained.

\paragraph{Expected hedge premium.}
The per-period value of pre-positioning one unit of $x_B$ while at
location $A$ is
\begin{equation}
    \text{hedge premium per unit}
    \;=\;
    p_{\text{work}} \times \tau_{\text{buy}}
    \;\approx\;
    0.06 \times 0.025
    \;=\; 0.15\%\ \text{per period}.
    \label{eq:hedge_premium}
\end{equation}
This is the incremental saving on buying costs that accrues when the
household has pre-positioned $x_B > 0$ prior to relocation.
Accumulated over a working-life horizon, equation~(\ref{eq:hedge_premium})
implies a lifetime CEV contribution of $1$--$2\%$ from the hedge
channel alone, with the continuous-$x$ rent-saving channel adding a further
$3$--$4\%$ (see Section~\ref{sec:results:decomp}).

\subsection{Cross-Location House-Price Correlation}
\label{sec:calibration:rhoAB}

The parameter $\rho_{AB}$ governs the idiosyncratic correlation between
location $A$ and location $B$ returns.
Using the decomposition~(\ref{eq:return_decomp}),
the observed return correlation between two MSAs is:
\begin{equation}
    \operatorname{Corr}(R_A, R_B)
    \;\approx\;
    \frac{\sigma_{\text{div}}^2}{\sigma_h^2}
    \;+\;
    \frac{\sigma_\iota^2}{\sigma_h^2}\,\rho_{AB}
    \;\approx\; 0.756 + 0.244\,\rho_{AB}.
    \label{eq:rhoAB_mapping}
\end{equation}
We set $\rho_{AB} = 0.50$, which maps to an observed return correlation of
$\approx 0.88$, consistent with pairs of geographically proximate
US metros in the Case-Shiller 20-city panel (2000--2023).
This is the midpoint of the $0.30$--$0.70$ range that covers inter-regional
moves (different metro areas): at one extreme, Chicago--Phoenix
($\hat{\rho}_{AB} \approx 0.35$) and at the other,
NY--Boston ($\hat{\rho}_{AB} \approx 0.60$).

Equation~(\ref{eq:rhoAB_mapping}) implies that the hedge value of
$x_B$ at location $A$ derives from bearing the idiosyncratic component
$\iota_B$ independently of $\iota_A$.
As $\rho_{AB} \to 1$, locations $A$ and $B$ become nearly identical
markets and $x_B$ is a near-substitute for $x_A$; the hedge channel
collapses.
The Round-4 referee pre-registered test~(m) requires $\operatorname{CEV}(E2\_2L \text{ vs } E1\_2L)$
to decrease significantly at $\rho_{AB} = 0.95$.

\subsection{Sensitivity Grid}
\label{sec:calibration:sensitivity}

Table~\ref{tab:sensitivity_grid} lists the pre-registered sensitivity
sweeps queued for execution after the v4 baseline results are available.
The three P1-priority sweeps cover the three parameters most directly
linked to the welfare decomposition channels.

\begin{table}[t]
\centering
\caption{Pre-Registered Sensitivity Sweeps}
\label{tab:sensitivity_grid}
\small
\begin{tabular}{llll}
\toprule
Parameter & Sweep values & Channel & Script \\
\midrule
$\rho_{AB}$          & $\{0,\, 0.25,\, 0.50,\, 0.75,\, 0.95\}$  & Hedge (maintained) & \texttt{sweep\_rhoAB.sh} \\
$p_{\text{work}}$    & $\{0,\, 0.02,\, 0.06,\, 0.12\}$           & Hedge (frequency)  & \texttt{sweep\_prelocate.sh} \\
$\tau_{\text{buy}}$  & $\{0,\, 0.025,\, 0.04,\, 0.06\}$          & Avoided tx-cost    & \texttt{sweep\_txcost.sh} \\
\midrule
\multicolumn{4}{l}{\textit{Phase 2 sweeps (after H2' approval)}} \\
$\gamma$             & $\{3,\, 5,\, 8\}$                         & Risk channel       & --- \\
$\delta_{\text{own}}$& $\{0.02,\, 0.04,\, 0.06\}$               & Rent-saving        & --- \\
$\sigma_{\text{div}}$& $\{0.08,\, 0.10,\, 0.12\}$               & Idio.\ risk share  & --- \\
\bottomrule
\end{tabular}
\end{table}

At $p_{\text{work}} = 0$ and at $\tau_{\text{buy}} = 0$,
the hedge premium~(\ref{eq:hedge_premium}) is exactly zero; any
remaining welfare gap between $E2\_2L$ and $E1\_2L$ in these cells
reflects the continuous-$x$ rent-saving channel only and serves
as a falsification bound.
At $\rho_{AB} = 0.95$, the maintained-hedge channel must also collapse
because $x_B$ and $x_A$ are near-substitutes.
These three boundary conditions define sharp predictions that distinguish
the two decomposed channels (Section~\ref{sec:results:decomp}).

%              Requires amsmath, booktabs, multirow, siunitx (for tabular).

% ─────────────────────────────────────────────────────────────────────────────
\section{Calibration}
\label{sec:calibration}
% ─────────────────────────────────────────────────────────────────────────────

Table~\ref{tab:parameters} summarises the baseline parameter values.
The calibration follows the household-finance lifecycle tradition
(Cocco, Gomes, and Maenhout, 2005; Yao and Zhang, 2005), extended
with three v4-specific objects: the relocation probability $p_{\text{work}}$,
the round-trip transaction-cost pair $(\tau_{\text{sell}}, \tau_{\text{buy}})$,
and the cross-location idiosyncratic correlation $\rho_{AB}$.
We discuss each group in turn.

% ── Table 1: Parameter summary ───────────────────────────────────────────────
\begin{table}[t]
\centering
\caption{Baseline Parameters}
\label{tab:parameters}
\small
\begin{tabular}{llcl}
\toprule
Parameter & Symbol & Value & Source \\
\midrule
\multicolumn{4}{l}{\textit{Preferences and lifecycle}} \\
Risk aversion         & $\gamma$         & 5.0  & CGM (2005) \\
Time discount         & $\beta$          & 0.96 & CGM (2005) \\
Entry age / retire / terminal & --- & 25 / 65 / 80 & Standard \\
\addlinespace
\multicolumn{4}{l}{\textit{Financial assets}} \\
Risk-free gross return     & $R_f$           & 1.02  & CGM (2005) \\
Equity risk premium        & $\mu_S - r_f$   & 4\%   & CGM (2005) \\
Stock return volatility    & $\sigma_S$      & 15.7\% & CGM (2005) \\
\addlinespace
\multicolumn{4}{l}{\textit{Income process (CGM 2005)}} \\
Permanent shock variance   & $\sigma_u^2$    & 0.0106 & CGM Table~I \\
Transitory shock variance  & $\sigma_\varepsilon^2$ & 0.0738 & CGM Table~I \\
Retirement replacement     & $\lambda$       & 0.65  & CGM (2005) \\
\addlinespace
\multicolumn{4}{l}{\textit{Housing return (Cocco 2005; Yao-Zhang 2005)}} \\
Expected log price growth  & $g_h$           & 1.6\%  & Cocco (2005) \\
Total housing return std   & $\sigma_h$      & 11.5\% & Cocco (2005) \\
Aggregate factor std       & $\sigma_{\text{div}}$ & 10.0\% & Case-Shiller / NAREIT \\
Idiosyncratic std (derived) & $\sigma_\iota$ & 5.7\%  & $\sqrt{\sigma_h^2 - \sigma_{\text{div}}^2}$ \\
Rent-to-price ratio        & $\varrho$       & 5.0\%  & Yao and Zhang (2005) \\
Maintenance-to-price ratio & $m$             & 1.0\%  & Cocco (2005) \\
Rent-saving wedge (derived) & $\delta_{\text{own}} = \varrho - m$ & 4.0\% & --- \\
\addlinespace
\multicolumn{4}{l}{\textit{Mobility (PSID-anchored)}} \\
Relocation prob, working age & $p_{\text{work}}$ & 6.0\% & PSID inter-MSA \\
Relocation prob, retired     & $p_{\text{ret}}$  & 2.0\% & PSID retired \\
\addlinespace
\multicolumn{4}{l}{\textit{Transaction costs}} \\
Traditional sell cost  & $\tau_{\text{sell}}$ & 6.0\% & NAR (2023) \\
Traditional buy cost   & $\tau_{\text{buy}}$  & 2.5\% & CFPB / ALTA \\
Token transfer cost    & $\tau_{\text{token}}$& 1.0\% & RealT / Lofty surveys \\
\addlinespace
\multicolumn{4}{l}{\textit{Cross-location correlation}} \\
Idiosyncratic corr., MSA pair & $\rho_{AB}$ & 0.50 & Case-Shiller (2000--2023) \\
\addlinespace
\multicolumn{4}{l}{\textit{Mortgage (baseline off)}} \\
LTV constraint & $\overline{\text{ltv}}$ & 0.0 & Baseline no-mortgage \\
\bottomrule
\end{tabular}
\end{table}

\subsection{Preferences and Income Process}
\label{sec:calibration:prefs}

Preferences and the income process are taken directly from
Cocco, Gomes, and Maenhout (2005, CGM hereafter).
Risk aversion $\gamma = 5$ and the discount factor $\beta = 0.96$ are
their baseline values.
The labor income age profile follows the CGM cubic polynomial fit
to PSID data:
\begin{equation}
    f(\text{age}) \;=\;
    -2.170 + 0.168\,(\text{age}/10)
    - 0.032\,(\text{age}/10)^2
    + 0.002\,(\text{age}/10)^3.
    \label{eq:income_profile}
\end{equation}
Permanent and transitory income shocks have variances $\sigma_u^2 = 0.0106$
and $\sigma_\varepsilon^2 = 0.0738$ (CGM Table~I).
The retirement replacement rate $\lambda = 0.65$ matches the median
Social Security replacement rate at median lifetime earnings.

Financial asset parameters follow the same source:
the equity premium is $4\%$, stock return volatility $15.7\%$,
and the risk-free rate $2\%$ real.

\subsection{Housing Return Decomposition}
\label{sec:calibration:housing}

Single-location housing return volatility $\sigma_h = 11.5\%$ is from
Cocco (2005).
We decompose the single-location log return into a shared aggregate factor
$\eta_{\text{div}}$ and an idiosyncratic component $\iota_\ell$:
\begin{equation}
    \log R_\ell \;=\; \mu_h + \eta_{\text{div}} + \iota_\ell,
    \qquad
    \eta_{\text{div}} \sim \mathcal{N}(0,\sigma_{\text{div}}^2),
    \quad
    \iota_\ell \sim \mathcal{N}(0,\sigma_\iota^2),
    \label{eq:return_decomp}
\end{equation}
with $\sigma_h^2 = \sigma_{\text{div}}^2 + \sigma_\iota^2$.
We set $\sigma_{\text{div}} = 10\%$ so that the aggregate factor accounts
for approximately $75\%$ of total single-location housing variance
($\sigma_{\text{div}}^2 / \sigma_h^2 \approx 0.756$), consistent with
the common factor dominating MSA-level variance in Case-Shiller data
and with Piazzesi and Schneider~(2016).
The idiosyncratic component $\sigma_\iota = \sqrt{0.115^2 - 0.10^2} \approx 5.7\%$
is calibrated residually.

The idiosyncratic components at the two locations follow a bivariate normal
with correlation $\rho_{AB}$ (Section~\ref{sec:calibration:rhoAB}).
The Cholesky factorisation is:
\begin{align}
    \iota_A &= \sigma_\iota\,\xi_A, \label{eq:chol_A} \\
    \iota_B &= \rho_{AB}\,\iota_A + \sqrt{1 - \rho_{AB}^2}\,\sigma_\iota\,\xi_B,
    \label{eq:chol_B}
\end{align}
with $\xi_A, \xi_B \sim \mathcal{N}(0,1)$ independent.

The rent-to-price ratio $\varrho = 5\%$ follows Yao and Zhang~(2005).
Maintenance costs $m = 1\%$ follow Cocco~(2005).
The rent-saving wedge $\delta_{\text{own}} = \varrho - m = 4\%$ governs
the per-period benefit of occupying a fractional token position
($\kappa$ rule, equation~(\ref{eq:kappa}) in Section~\ref{sec:model:kappa}).

\subsection{Mobility Rate}
\label{sec:calibration:mobility}

We anchor the relocation probability to PSID inter-MSA mobility data.
Total annual relocation rates in the PSID average $7$--$12\%$ for
working-age households but include short-distance, within-MSA moves that
do not change the relevant housing market.
Long-distance, inter-MSA moves account for approximately $40$--$50\%$
of total PSID mobility, yielding an inter-MSA rate of $5$--$7\%$ annually
for prime-age workers.
We set $p_{\text{work}} = 6\%$ (PSID midpoint) for working ages
$25$--$64$ and $p_{\text{ret}} = 2\%$ for ages $65$--$80$.

This choice is consistent with Bagliano, Fugazza, and Nicodano~(2014, RFS),
who calibrate career-driven mobility shocks at $5$--$8\%$ annually, and
broadly in the range of Yao and Zhang~(2005) ($4\%$) and Saks and
Wozniak~(2011, JLE) ($2$--$4\%$ inter-state, implying $4$--$8\%$ inter-MSA).

The sensitivity grid (Table~\ref{tab:sensitivity_grid} below)
spans $p_{\text{work}} \in \{0, 0.02, 0.06, 0.12\}$.
At $p_{\text{work}} = 0$, the cross-location hedge channel must vanish by
construction (no relocation, no cost-saving motive); this is a pre-registered
falsification test (Referee Round~4, test~r).

\subsection{Transaction Costs}
\label{sec:calibration:txcost}

\paragraph{Traditional housing ($E1\_2L$).}
We set the seller commission and closing costs at $\tau_{\text{sell}} = 6\%$,
consistent with the NAR~(2023) aggregate for a typical single-family
home sale ($5$--$6\%$ agent commission plus $0.5$--$1\%$ closing costs).
The buyer cost is $\tau_{\text{buy}} = 2.5\%$,
covering origination fees ($0.75\%$), title insurance and escrow
($0.75\%$), appraisal and inspection ($0.25\%$), and prepaid items
($0.25\%$), consistent with CFPB and ALTA surveys.
The implied round-trip cost for an $E1\_2L$ relocation is
$\tau_{\text{sell}} + \tau_{\text{buy}} = 8.5\%$,
within the $5$--$10\%$ range used in the structural housing literature
(Diaz and Luengo-Prado, 2008; Head, Lloyd-Ellis, and Sun, 2014).

\paragraph{Tokenized housing ($E2\_2L$).}
Platform surveys of RealT, Lofty, and DigiShares~(2024) report trading
fees in the range $0.1$--$0.5\%$ plus blockchain gas and settlement
of $0.05$--$0.3\%$ on Layer-2 networks.
We set $\tau_{\text{token}} = 1\%$ to include a regulatory-compliance
margin of $0.5\%$ (reflecting SEC Regulation~D and AML overhead).
The structural distinction from $E1\_2L$ is that $\tau_{\text{token}} \ll \tau_{\text{sell}} + \tau_{\text{buy}}$:
tokens are portable across relocations, so the $8.5\%$ round-trip is avoided
for every unit of cross-location position retained.

\paragraph{Expected hedge premium.}
The per-period value of pre-positioning one unit of $x_B$ while at
location $A$ is
\begin{equation}
    \text{hedge premium per unit}
    \;=\;
    p_{\text{work}} \times \tau_{\text{buy}}
    \;\approx\;
    0.06 \times 0.025
    \;=\; 0.15\%\ \text{per period}.
    \label{eq:hedge_premium}
\end{equation}
This is the incremental saving on buying costs that accrues when the
household has pre-positioned $x_B > 0$ prior to relocation.
Accumulated over a working-life horizon, equation~(\ref{eq:hedge_premium})
implies a lifetime CEV contribution of $1$--$2\%$ from the hedge
channel alone, with the continuous-$x$ rent-saving channel adding a further
$3$--$4\%$ (see Section~\ref{sec:results:decomp}).

\subsection{Cross-Location House-Price Correlation}
\label{sec:calibration:rhoAB}

The parameter $\rho_{AB}$ governs the idiosyncratic correlation between
location $A$ and location $B$ returns.
Using the decomposition~(\ref{eq:return_decomp}),
the observed return correlation between two MSAs is:
\begin{equation}
    \operatorname{Corr}(R_A, R_B)
    \;\approx\;
    \frac{\sigma_{\text{div}}^2}{\sigma_h^2}
    \;+\;
    \frac{\sigma_\iota^2}{\sigma_h^2}\,\rho_{AB}
    \;\approx\; 0.756 + 0.244\,\rho_{AB}.
    \label{eq:rhoAB_mapping}
\end{equation}
We set $\rho_{AB} = 0.50$, which maps to an observed return correlation of
$\approx 0.88$, consistent with pairs of geographically proximate
US metros in the Case-Shiller 20-city panel (2000--2023).
This is the midpoint of the $0.30$--$0.70$ range that covers inter-regional
moves (different metro areas): at one extreme, Chicago--Phoenix
($\hat{\rho}_{AB} \approx 0.35$) and at the other,
NY--Boston ($\hat{\rho}_{AB} \approx 0.60$).

Equation~(\ref{eq:rhoAB_mapping}) implies that the hedge value of
$x_B$ at location $A$ derives from bearing the idiosyncratic component
$\iota_B$ independently of $\iota_A$.
As $\rho_{AB} \to 1$, locations $A$ and $B$ become nearly identical
markets and $x_B$ is a near-substitute for $x_A$; the hedge channel
collapses.
The Round-4 referee pre-registered test~(m) requires $\operatorname{CEV}(E2\_2L \text{ vs } E1\_2L)$
to decrease significantly at $\rho_{AB} = 0.95$.

\subsection{Sensitivity Grid}
\label{sec:calibration:sensitivity}

Table~\ref{tab:sensitivity_grid} lists the pre-registered sensitivity
sweeps queued for execution after the v4 baseline results are available.
The three P1-priority sweeps cover the three parameters most directly
linked to the welfare decomposition channels.

\begin{table}[t]
\centering
\caption{Pre-Registered Sensitivity Sweeps}
\label{tab:sensitivity_grid}
\small
\begin{tabular}{llll}
\toprule
Parameter & Sweep values & Channel & Script \\
\midrule
$\rho_{AB}$          & $\{0,\, 0.25,\, 0.50,\, 0.75,\, 0.95\}$  & Hedge (maintained) & \texttt{sweep\_rhoAB.sh} \\
$p_{\text{work}}$    & $\{0,\, 0.02,\, 0.06,\, 0.12\}$           & Hedge (frequency)  & \texttt{sweep\_prelocate.sh} \\
$\tau_{\text{buy}}$  & $\{0,\, 0.025,\, 0.04,\, 0.06\}$          & Avoided tx-cost    & \texttt{sweep\_txcost.sh} \\
\midrule
\multicolumn{4}{l}{\textit{Phase 2 sweeps (after H2' approval)}} \\
$\gamma$             & $\{3,\, 5,\, 8\}$                         & Risk channel       & --- \\
$\delta_{\text{own}}$& $\{0.02,\, 0.04,\, 0.06\}$               & Rent-saving        & --- \\
$\sigma_{\text{div}}$& $\{0.08,\, 0.10,\, 0.12\}$               & Idio.\ risk share  & --- \\
\bottomrule
\end{tabular}
\end{table}

At $p_{\text{work}} = 0$ and at $\tau_{\text{buy}} = 0$,
the hedge premium~(\ref{eq:hedge_premium}) is exactly zero; any
remaining welfare gap between $E2\_2L$ and $E1\_2L$ in these cells
reflects the continuous-$x$ rent-saving channel only and serves
as a falsification bound.
At $\rho_{AB} = 0.95$, the maintained-hedge channel must also collapse
because $x_B$ and $x_A$ are near-substitutes.
These three boundary conditions define sharp predictions that distinguish
the two decomposed channels (Section~\ref{sec:results:decomp}).

% ─────────────────────────────────────────────────────────────────────────────

\clearpage

% ─────────────────────────────────────────────────────────────────────────────
% s4_results.tex — Results section
% Paper: "Tokenized Housing and Lifecycle Portfolio Choice:
%         A Decoupling of Location from Housing Exposure"
%
% Source: docs/welfare_decomp_v4.md (CEV formula, table shells, channel plan)
%         docs/sensitivity_grid_v4.md (sweep design)
%         paper/outline_v4.md (section 4 structure)
%
% Status: SKELETON — numerical placeholders marked [PLACEHOLDER] throughout.
%         All table shells match the pre-registered spec in welfare_decomp_v4.md.
%         Figures 1 and 2 are placeholder stubs.
%         Fill numerical entries once output/diagnostics/p6_option1_*.json land.
%
% Compilation: included from main.tex via % s4_results.tex — Results section
% Paper: "Tokenized Housing and Lifecycle Portfolio Choice:
%         A Decoupling of Location from Housing Exposure"
%
% Source: docs/welfare_decomp_v4.md (CEV formula, table shells, channel plan)
%         docs/sensitivity_grid_v4.md (sweep design)
%         paper/outline_v4.md (section 4 structure)
%
% Status: SKELETON — numerical placeholders marked [PLACEHOLDER] throughout.
%         All table shells match the pre-registered spec in welfare_decomp_v4.md.
%         Figures 1 and 2 are placeholder stubs.
%         Fill numerical entries once output/diagnostics/p6_option1_*.json land.
%
% Compilation: included from main.tex via % s4_results.tex — Results section
% Paper: "Tokenized Housing and Lifecycle Portfolio Choice:
%         A Decoupling of Location from Housing Exposure"
%
% Source: docs/welfare_decomp_v4.md (CEV formula, table shells, channel plan)
%         docs/sensitivity_grid_v4.md (sweep design)
%         paper/outline_v4.md (section 4 structure)
%
% Status: SKELETON — numerical placeholders marked [PLACEHOLDER] throughout.
%         All table shells match the pre-registered spec in welfare_decomp_v4.md.
%         Figures 1 and 2 are placeholder stubs.
%         Fill numerical entries once output/diagnostics/p6_option1_*.json land.
%
% Compilation: included from main.tex via \input{sections/s4_results}
%              Requires amsmath, booktabs, multirow, siunitx, graphicx, caption.

% ─────────────────────────────────────────────────────────────────────────────
\section{Results}
\label{sec:results}
% ─────────────────────────────────────────────────────────────────────────────

We report the baseline welfare results in Section~\ref{sec:results:baseline},
the channel decomposition in Section~\ref{sec:results:decomp}, the
pre-registered falsification tests in Section~\ref{sec:results:falsification},
and sensitivity analysis in Section~\ref{sec:results:sensitivity}.

All welfare gains are expressed as consumption equivalent variations (CEV) at the
initial state
$\mathbf{s}_0 = (t{=}1,\,w_{\text{mid}},\,z_{\text{mid}},\,\ell{=}A,\,x_{A,0}{=}0,\,x_{B,0}{=}0)$:
\begin{equation}
    \text{CEV}(R_a \text{ vs } R_b) \;=\;
    \left(\frac{V_1^{R_a}(\mathbf{s}_0)}{V_1^{R_b}(\mathbf{s}_0)}\right)^{1/(1-\gamma)}
    - 1,
    \label{eq:cev}
\end{equation}
where $V_1^{R}(\mathbf{s}_0)$ is the time-1 value function under regime $R$
evaluated at $\mathbf{s}_0$.%
\footnote{Both value functions are negative under CRRA with finite wealth,
so the ratio is positive and the formula is well-defined for $\gamma \neq 1$.
The CEV measures the fraction of lifetime consumption in regime $R_b$ that the
household would give up to gain permanent access to regime $R_a$.}
Appendix~\ref{app:numerical} reports convergence diagnostics and full-grid
robustness.

% ─────────────────────────────────────────────────────────────────────────────
\subsection{Baseline Welfare Results}
\label{sec:results:baseline}
% ─────────────────────────────────────────────────────────────────────────────

The headline welfare result is reported in Table~\ref{tab:cev_baseline}.
The first row gives $\text{CEV}(\text{E2\_2L vs E1\_2L})$, the lifetime
welfare value of tokenization relative to traditional binary homeownership.
Rows 2--4 decompose this gain into three channels (detailed in
Section~\ref{sec:results:decomp}).
Row 5 reports $\text{CEV}(\text{E0 vs E1\_2L})$ as a sanity-check benchmark:
the welfare cost of pure renting relative to traditional ownership.

\begin{table}[t]
\centering
\caption{Welfare Value of Tokenization: Baseline and Robustness Panel
         (CEV, percentage points of lifetime consumption)}
\label{tab:cev_baseline}
\small
\begin{tabular}{lcccccc}
\toprule
 & Baseline & $\rho_{AB}=0$ & $\rho_{AB}=0.50$ & $\rho_{AB}=0.95$
 & $p_{\text{reloc}}=0$ & $p_{\text{reloc}}=0.12$ \\
\midrule
$\text{CEV}(\text{E2\_2L vs E1\_2L})$
    & \textsc{[P]} & \textsc{[P]} & \textsc{[P]} & \textsc{[P]}
    & \textsc{[P]} & \textsc{[P]} \\[2pt]
\quad Avoided-tx channel
    & \textsc{[P]} & \textsc{[P]} & \textsc{[P]} & \textsc{[P]}
    & \textsc{[P]} & \textsc{[P]} \\
\quad Maintained-hedge channel
    & \textsc{[P]} & \textsc{[P]} & \textsc{[P]} & \textsc{[P]}
    & \textsc{[P]} & \textsc{[P]} \\
\quad Cross-term $\xi$
    & \textsc{[P]} & \textsc{[P]} & \textsc{[P]} & \textsc{[P]}
    & \textsc{[P]} & \textsc{[P]} \\[4pt]
$\text{CEV}(\text{E0 vs E1\_2L})$ (renter benchmark)
    & \textsc{[P]} & --- & --- & --- & --- & --- \\
\addlinespace
Mean $\bar{x}_B$ at $\ell{=}A$
    & \textsc{[P]} & \textsc{[P]} & \textsc{[P]} & \textsc{[P]}
    & \textsc{[P]} & \textsc{[P]} \\
\bottomrule
\end{tabular}
\vspace{4pt}
\begin{minipage}{\linewidth}
\small\textit{Notes:}
Baseline: $\gamma=5$, $\beta=0.96$, $\rho_{AB}=0.50$,
$p_{\text{reloc,work}}=0.06$, $\tau_{\text{sell}}=0.06$,
$\tau_{\text{buy}}=0.025$, $\tau_{\text{token}}=0.01$.
Grid: $N_W=15$, $N_Z=5$, $N_{x_{\text{prev}}}=3$ (coarse; full-grid
robustness at $N_W=40$, $N_Z=9$, $N_{x_{\text{prev}}}=5$ in
Appendix~\ref{app:numerical}).
\textsc{[P]} = placeholder pending server1 baseline runs.
\end{minipage}
\end{table}

The primary diagnostic for mechanism activation is $\bar{x}_B$ at $\ell{=}A$:
the mean portfolio share of location-$B$ tokens held by households residing
at location $A$.
Under E1\_2L, $\bar{x}_B = 0$ by admissibility (binary tenure forbids
non-occupancy-location ownership).
Under E2\_2L, a positive $\bar{x}_B > 0$ at $\ell{=}A$ indicates that
households voluntarily pre-accumulate location-$B$ tokens before relocation ---
the mechanism that is uniquely enabled by token portability and the 6D state
extension (Option~1).
This is the key discriminator from the v3 Option~3 approximation, under which
$\bar{x}_B = 0$ held at all mobility rates.

Figure~\ref{fig:lifecycle_profiles} illustrates the mechanism at the
lifecycle level: the age profiles of mean $x_A$ and $x_B$ holdings
under E1\_2L and E2\_2L.
The key prediction is that $\bar{x}_B$ rises during working years (ages 25--65),
when $p_{\text{reloc,work}} = 0.06$, and declines post-retirement when the
relocation probability falls to $p_{\text{reloc,ret}} = 0.02$.
This pattern --- confirmed or denied by the server1 policy-array output ---
constitutes the strongest visual evidence for the pre-buying hedge mechanism.

\begin{figure}[t]
\centering
%% PLACEHOLDER — replace with lifecycle profile plots once server1 policy arrays land.
%% See paper/exhibit_memos/fig1_lifecycle_profiles.md for full production spec.
\fbox{%
  \begin{minipage}{0.88\linewidth}
  \vspace{2.5cm}
  \centering
  {\sffamily\small [PLACEHOLDER: Figure~1 — Lifecycle token-holding profiles.
  Panel A (E1\_2L): mean $x_A$ (solid) and $x_B = 0$ (dashed) over ages 25--80.
  Panel B (E2\_2L): mean $x_A$ (solid) and mean $x_B > 0$ (dashed) over ages 25--80.
  Source: per-age mean of \texttt{xA\_policy} and \texttt{xB\_policy} arrays at $\ell{=}A$,
  averaged over feasible $(w, z, x_{\text{prev}})$ states.
  Script: \texttt{scripts/plot\_lifecycle\_profiles.py}.]}
  \vspace{2.5cm}
  \end{minipage}%
}
\caption{Lifecycle mean token-holding profiles under traditional ownership (E1\_2L,
         left panel) and tokenized ownership (E2\_2L, right panel) at location $A$.
         E2\_2L households pre-accumulate location-$B$ tokens during working years
         (positive $\bar{x}_B$), saving the buying cost $\tau_{\text{buy}}$ on
         anticipated future relocation.
         E1\_2L households hold $\bar{x}_B = 0$ at all ages by admissibility.
         \textsc{[P]} = placeholder pending server1 run.}
\label{fig:lifecycle_profiles}
\end{figure}

% ─────────────────────────────────────────────────────────────────────────────
\subsection{Channel Decomposition}
\label{sec:results:decomp}
% ─────────────────────────────────────────────────────────────────────────────

We decompose the total welfare gain using three regime runs.
Let E1\_2L$_{\text{NT}}$ denote the no-transaction-cost counterfactual
($\tau_{\text{sell}} = \tau_{\text{buy}} = 0$, binary tenure otherwise).
The decomposition is:
\begin{align}
    \underbrace{\text{CEV}(\text{E2\_2L vs E1\_2L})}_{\text{total}}
    &\;=\;
    \underbrace{\text{CEV}(\text{E1\_2L}_{\text{NT}} \text{ vs E1\_2L})}_{\text{avoided-tx}}
    \;+\;
    \underbrace{\text{CEV}(\text{E2\_2L vs E1\_2L}_{\text{NT}})}_{\text{maintained-hedge}}
    \;+\; \xi,
    \label{eq:decomp}
\end{align}
where $\xi$ is the cross-term.
This decomposition is pre-registered in \texttt{docs/welfare\_decomp\_v4.md}
and is exact up to the cross-term, which was empirically near zero in v3
($\xi = 0.021$\%) and is expected to remain small in v4.

\begin{table}[t]
\centering
\caption{Channel Decomposition of the Tokenization Welfare Gain
         (CEV, percentage points of lifetime consumption)}
\label{tab:decomp}
\small
\begin{tabular}{lccl}
\toprule
Channel & CEV (\%) & Share of total & Description \\
\midrule
Total: $\text{CEV}(\text{E2\_2L vs E1\_2L})$
    & \textsc{[P]} & 100\% & Headline tokenization gain \\
\addlinespace
\quad (i)~Avoided-tx
    & \textsc{[P]} & \textsc{[P]}\%
    & $\tau_{\text{sell}} + \tau_{\text{buy}}$ burden of E1\_2L moves \\
\qquad of which $\tau_{\text{sell}}$
    & \textsc{[P]} & \textsc{[P]}\%
    & NAR selling commission (6\%) \\
\qquad of which $\tau_{\text{buy}}$
    & \textsc{[P]} & \textsc{[P]}\%
    & Buying closing costs (2.5\%) \\[4pt]
\quad (ii)~Maintained-hedge
    & \textsc{[P]} & \textsc{[P]}\%
    & Continuous-x + pre-buying hedge vs E1\_2L$_\text{NT}$ \\[4pt]
\quad (iii)~Cross-term $\xi$
    & \textsc{[P]} & \textsc{[P]}\%
    & Interaction; expected $\approx 0$ \\
\midrule
v3 Option~3 (no state ext.)
    & $+4.255$\% & (reference) & Previous fire's upper-bound approx \\
\quad of which avoided-tx
    & $+0.816$\% & 19.2\% & From v3 full-grid run \\
\quad of which continuous-x
    & $+3.411$\% & 80.2\% & Liu (2021) territory \\
\quad of which pre-buying hedge
    & $\approx 0$\% & $\approx 0$\% & Absent without state ext. (Option~3 limit) \\
\bottomrule
\end{tabular}
\vspace{4pt}
\begin{minipage}{\linewidth}
\small\textit{Notes:}
Upper panel: v4 Option~1 results (\textsc{[P]} = placeholder).
Lower panel: v3 Option~3 reference numbers (confirmed from server1 run,
commit 186da13).
The maintained-hedge channel in v4 encompasses both the continuous-x
subchannel (present in v3) and the new pre-buying hedge subchannel
enabled by the 6D state; a fourth regime decomposition
($\text{E2\_2L}_{\text{NT}}$) can further separate the two if H1 is confirmed.
\end{minipage}
\end{table}

The maintained-hedge channel in v4 captures two sub-mechanisms:
\begin{enumerate}
    \item \emph{Continuous-$x$ sub-channel} (present in v3): continuous
    fractional ownership of the occupied unit avoids the binary kink at
    $x_\ell = 1$, delivering a Liu~(2021)-style indivisibility relaxation.
    \item \emph{Pre-buying hedge sub-channel} (new in v4): by tracking
    $(x_{A,t-1}, x_{B,t-1})$ in the state, households face a per-period
    buying cost $\tau_{\text{buy}} \cdot \max(\Delta x, 0)$ and can
    pre-accumulate location-$B$ tokens at location $A$, arriving at $B$
    with $x_{B,t-1} > 0$ and avoiding the lump-sum purchase cost.
    The expected per-period premium per unit held is
    $p_{\text{reloc}} \times \tau_{\text{buy}} = 0.06 \times 0.025 = 0.15$\%
    (see Appendix~\ref{app:proof}).
\end{enumerate}

% ─────────────────────────────────────────────────────────────────────────────
\subsection{Falsification Tests}
\label{sec:results:falsification}
% ─────────────────────────────────────────────────────────────────────────────

We conduct three pre-registered falsification tests designed to identify
the pre-buying hedge channel.
These tests were pre-registered in \texttt{docs/welfare\_decomp\_v4.md}
before server1 runs were executed.
Crucially, the analogous tests \emph{failed} under v3 Option~3 (which lacked
the state extension): $\bar{x}_B$ remained zero at all mobility rates,
and the CEV was insensitive to $p_{\text{reloc}}$ and $\rho_{AB}$.

\begin{table}[t]
\centering
\caption{Pre-Registered Falsification Tests
         (Deviations from Baseline)}
\label{tab:falsification}
\small
\begin{tabular}{llcccc}
\toprule
Test & Counterfactual & $\text{CEV}$ & $\bar{x}_B$ at $\ell{=}A$
     & Pass criterion & Result \\
\midrule
(r) Relocation disabled
    & $p_{\text{reloc}} = 0$
    & \textsc{[P]}
    & \textsc{[P]}
    & $\bar{x}_B \to 0$;\; $\Delta\text{CEV} < -0.5$pp
    & \textsc{[P]} \\[4pt]
(m) Perfect correlation
    & $\rho_{AB} = 0.95$
    & \textsc{[P]}
    & \textsc{[P]}
    & $\bar{x}_B$ decreases;\; $\Delta\text{CEV} < 0$
    & \textsc{[P]} \\[4pt]
(q) No buying cost
    & $\tau_{\text{buy}} = 0$
    & \textsc{[P]}
    & \textsc{[P]}
    & Hedge channel $\to 0$;\; $\bar{x}_B \to 0$
    & \textsc{[P]} \\
\bottomrule
\end{tabular}
\vspace{4pt}
\begin{minipage}{\linewidth}
\small\textit{Notes:}
All tests compare to the baseline E2\_2L v4 run.
Pass criterion stated prior to server1 runs.
``Hedge channel'' in test (q) is
$\text{CEV}(\text{E2\_2L}_{\tau=0} \text{ vs E1\_2L}_{\text{NT}}) - \text{CEV}(\text{E2\_2L}_{\text{baseline}} \text{ vs E1\_2L}_{\text{NT}})$.
\textsc{[P]} = placeholder.
\end{minipage}
\end{table}

Test (r) shuts off the relocation shock entirely.
Without future relocation, there is no value to pre-accumulating
location-$B$ tokens while at $A$; the pre-buying hedge channel
should collapse.
If test (r) passes --- i.e., $\bar{x}_B \to 0$ when $p_{\text{reloc}} = 0$
--- it provides strong evidence that the positive $\bar{x}_B$ in the
baseline is driven by the pre-buying motive and not by a model artefact.

Test (m) raises the cross-location correlation to $\rho_{AB} = 0.95$,
nearly eliminating diversification benefit.
When location-$A$ and location-$B$ price processes are nearly identical,
holding $x_B$ at $\ell = A$ provides little hedge against diverging location returns.
The CEV and $\bar{x}_B$ should both decrease relative to the baseline.

Test (q) eliminates the buying cost ($\tau_{\text{buy}} = 0$), which
removes the financial incentive to pre-accumulate $x_B$ (there is no lump-sum
arrival cost to avoid).
The pre-buying hedge channel should collapse to zero.

% ─────────────────────────────────────────────────────────────────────────────
\subsection{Sensitivity Analysis}
\label{sec:results:sensitivity}
% ─────────────────────────────────────────────────────────────────────────────

Figure~\ref{fig:sensitivity_heatmap} depicts the total welfare gain
$\text{CEV}(\text{E2\_2L vs E1\_2L})$ as a function of the two primary
mechanism parameters: the cross-location correlation $\rho_{AB}$ and
the relocation probability $p_{\text{reloc}}$.
The full sweep design follows the pre-registered plan in
\texttt{docs/sensitivity\_grid\_v4.md}; only a subset is shown here.

\begin{figure}[t]
\centering
%% PLACEHOLDER — replace with actual heatmap once sweep results available
\fbox{%
  \begin{minipage}{0.72\linewidth}
  \vspace{2.5cm}
  \centering
  {\sffamily\small [PLACEHOLDER: Figure~2 — CEV sensitivity heatmap.
  X-axis: $\rho_{AB} \in \{0, 0.25, 0.50, 0.75, 0.95\}$.
  Y-axis: $p_{\text{reloc}} \in \{0, 0.02, 0.06, 0.12\}$.
  Heat colour: total CEV(\%), from dark (low) to bright (high).
  Source: sweep outputs from \texttt{scripts/sweep\_rhoAB.sh} and
  \texttt{scripts/sweep\_prelocate.sh} once server1 runs complete.]}
  \vspace{2.5cm}
  \end{minipage}%
}
\caption{Welfare sensitivity to cross-location correlation and relocation rate.
         Each cell reports $\text{CEV}(\text{E2\_2L vs E1\_2L})$ at
         $(\rho_{AB},\, p_{\text{reloc}})$.
         The mechanism prediction is a positive gradient in $p_{\text{reloc}}$
         (higher mobility $\Rightarrow$ higher hedge value) and a negative
         gradient in $\rho_{AB}$ (more correlated locations $\Rightarrow$
         smaller diversification benefit).}
\label{fig:sensitivity_heatmap}
\end{figure}

Table~\ref{tab:sensitivity} reports selected cross-sections of the
sensitivity grid for other key parameters.

\begin{table}[t]
\centering
\caption{Sensitivity of Headline CEV to Key Parameters
         (Total CEV, \%; baseline value on diagonal)}
\label{tab:sensitivity}
\small
\begin{tabular}{lcccccc}
\toprule
Parameter & & \multicolumn{5}{c}{Values} \\
\midrule
Risk aversion $\gamma$
    & & 3 & \textbf{5} & 8 & & \\
\quad $\text{CEV}(\text{E2\_2L vs E1\_2L})$
    & & \textsc{[P]} & \textsc{[P]} & \textsc{[P]} & & \\
\addlinespace
Buying cost $\tau_{\text{buy}}$
    & & 0\% & 1\% & \textbf{2.5\%} & 4\% & 6\% \\
\quad $\text{CEV}(\text{E2\_2L vs E1\_2L})$
    & & \textsc{[P]} & \textsc{[P]} & \textsc{[P]} & \textsc{[P]} & \textsc{[P]} \\
\addlinespace
Relocation rate $p_{\text{reloc,work}}$
    & & 0 & 0.02 & \textbf{0.06} & 0.12 & \\
\quad $\text{CEV}(\text{E2\_2L vs E1\_2L})$
    & & \textsc{[P]} & \textsc{[P]} & \textsc{[P]} & \textsc{[P]} & \\
\quad $\bar{x}_B$ at $\ell{=}A$
    & & \textsc{[P]} & \textsc{[P]} & \textsc{[P]} & \textsc{[P]} & \\
\addlinespace
Cross-location corr.\ $\rho_{AB}$
    & & 0 & 0.25 & \textbf{0.50} & 0.75 & 0.95 \\
\quad $\text{CEV}(\text{E2\_2L vs E1\_2L})$
    & & \textsc{[P]} & \textsc{[P]} & \textsc{[P]} & \textsc{[P]} & \textsc{[P]} \\
\addlinespace
\multicolumn{7}{l}{\textit{Asymmetric robustness (solver v4 extension, fire 22)}} \\
Location-$B$ return premium $\Delta\mu_B$
    & & $-1\%$ & \textbf{0\%} & $+1\%$ & & \\
\quad $\text{CEV}(\text{E2\_2L vs E1\_2L})$
    & & \textsc{[P]} & \textsc{[P]} & \textsc{[P]} & & \\
\quad $\bar{x}_B$ at $\ell{=}A$
    & & \textsc{[P]} & \textsc{[P]} & \textsc{[P]} & & \\[2pt]
Directional mobility ($p_{A\to B}$, $p_{B\to A}$)
    & & (0.06, 0.06) & (0.10, 0.03) & (0.10, 0.10) & & \\
\quad $\text{CEV}(\text{E2\_2L vs E1\_2L})$
    & & \textsc{[P]} & \textsc{[P]} & \textsc{[P]} & & \\
\addlinespace
\multicolumn{7}{l}{\textit{Mortgage activation (LTV constraint)}} \\
LTV ceiling $\theta_{\max}$
    & & \textbf{0} & 0.50 & 0.80 & & \\
\quad $\text{CEV}(\text{E2\_2L vs E1\_2L})$
    & & \textsc{[P]} & \textsc{[P]} & \textsc{[P]} & & \\
\bottomrule
\end{tabular}
\vspace{4pt}
\begin{minipage}{\linewidth}
\small\textit{Notes:}
Bold entries mark the baseline value.
One parameter varies at a time; all others held at baseline.
Sweep scripts: \texttt{scripts/sweep\_rhoAB.sh},
\texttt{scripts/sweep\_prelocate.sh}, \texttt{scripts/sweep\_txcost.sh},
\texttt{scripts/sweep\_asymmetric.sh} (location-$B$ return premium and directional mobility),
\texttt{scripts/sweep\_mortgage.sh} (LTV).
$\Delta\mu_B = \mu_{h,B} - \mu_{h,A}$: the excess Jensen-corrected log-mean return of
location $B$ relative to $A$.
Positive $\Delta\mu_B$ strengthens the pre-buying motive for $x_B$ by adding a return
premium to the tau\_buy-saving incentive.
\textsc{[P]} = placeholder pending server1 sweep runs.
\end{minipage}
\end{table}

Two monotonicity predictions follow from the model mechanism.
First, for $\tau_{\text{buy}}$: as the buying cost rises, the pre-accumulation
motive strengthens and $\bar{x}_B$ should increase, lifting the hedge sub-channel.
Second, for $p_{\text{reloc}}$: higher mobility raises the expected value of
pre-accumulated $x_B$ holdings; at $p_{\text{reloc}} = 0$ the hedge channel
collapses (falsification test r).
If these monotonicity patterns hold in the sweep results, they constitute
additional evidence for the pre-buying hedge mechanism.

% ─────────────────────────────────────────────────────────────────────────────
\subsection{Summary of Mechanism Evidence}
\label{sec:results:summary}
% ─────────────────────────────────────────────────────────────────────────────

The results establish \textsc{[P]} as the headline lifetime welfare gain from
tokenization under the baseline calibration.
The gain decomposes additively into an avoided-transaction-cost channel
(\textsc{[P]}\%) and a maintained-hedge channel (\textsc{[P]}\%), with
a near-zero cross-term.
The maintained-hedge channel is itself divisible into the continuous-$x$
sub-channel present in prior work (Liu, 2021) and the novel pre-buying hedge
sub-channel unlocked by the 6D state extension.

The three falsification tests provide mechanism identification.
Under all three counterfactuals --- zero relocation, high cross-location
correlation, zero buying cost --- the pre-buying hedge channel should
attenuate or collapse.
Confirmation of these patterns will establish that the welfare gain reported
in Table~\ref{tab:cev_baseline} reflects the model's structural mechanism
rather than a calibration artefact.
\textsc{[Fill in pass/fail outcomes from output/diagnostics/p6\_option1\_decomposition.md.]}

%              Requires amsmath, booktabs, multirow, siunitx, graphicx, caption.

% ─────────────────────────────────────────────────────────────────────────────
\section{Results}
\label{sec:results}
% ─────────────────────────────────────────────────────────────────────────────

We report the baseline welfare results in Section~\ref{sec:results:baseline},
the channel decomposition in Section~\ref{sec:results:decomp}, the
pre-registered falsification tests in Section~\ref{sec:results:falsification},
and sensitivity analysis in Section~\ref{sec:results:sensitivity}.

All welfare gains are expressed as consumption equivalent variations (CEV) at the
initial state
$\mathbf{s}_0 = (t{=}1,\,w_{\text{mid}},\,z_{\text{mid}},\,\ell{=}A,\,x_{A,0}{=}0,\,x_{B,0}{=}0)$:
\begin{equation}
    \text{CEV}(R_a \text{ vs } R_b) \;=\;
    \left(\frac{V_1^{R_a}(\mathbf{s}_0)}{V_1^{R_b}(\mathbf{s}_0)}\right)^{1/(1-\gamma)}
    - 1,
    \label{eq:cev}
\end{equation}
where $V_1^{R}(\mathbf{s}_0)$ is the time-1 value function under regime $R$
evaluated at $\mathbf{s}_0$.%
\footnote{Both value functions are negative under CRRA with finite wealth,
so the ratio is positive and the formula is well-defined for $\gamma \neq 1$.
The CEV measures the fraction of lifetime consumption in regime $R_b$ that the
household would give up to gain permanent access to regime $R_a$.}
Appendix~\ref{app:numerical} reports convergence diagnostics and full-grid
robustness.

% ─────────────────────────────────────────────────────────────────────────────
\subsection{Baseline Welfare Results}
\label{sec:results:baseline}
% ─────────────────────────────────────────────────────────────────────────────

The headline welfare result is reported in Table~\ref{tab:cev_baseline}.
The first row gives $\text{CEV}(\text{E2\_2L vs E1\_2L})$, the lifetime
welfare value of tokenization relative to traditional binary homeownership.
Rows 2--4 decompose this gain into three channels (detailed in
Section~\ref{sec:results:decomp}).
Row 5 reports $\text{CEV}(\text{E0 vs E1\_2L})$ as a sanity-check benchmark:
the welfare cost of pure renting relative to traditional ownership.

\begin{table}[t]
\centering
\caption{Welfare Value of Tokenization: Baseline and Robustness Panel
         (CEV, percentage points of lifetime consumption)}
\label{tab:cev_baseline}
\small
\begin{tabular}{lcccccc}
\toprule
 & Baseline & $\rho_{AB}=0$ & $\rho_{AB}=0.50$ & $\rho_{AB}=0.95$
 & $p_{\text{reloc}}=0$ & $p_{\text{reloc}}=0.12$ \\
\midrule
$\text{CEV}(\text{E2\_2L vs E1\_2L})$
    & \textsc{[P]} & \textsc{[P]} & \textsc{[P]} & \textsc{[P]}
    & \textsc{[P]} & \textsc{[P]} \\[2pt]
\quad Avoided-tx channel
    & \textsc{[P]} & \textsc{[P]} & \textsc{[P]} & \textsc{[P]}
    & \textsc{[P]} & \textsc{[P]} \\
\quad Maintained-hedge channel
    & \textsc{[P]} & \textsc{[P]} & \textsc{[P]} & \textsc{[P]}
    & \textsc{[P]} & \textsc{[P]} \\
\quad Cross-term $\xi$
    & \textsc{[P]} & \textsc{[P]} & \textsc{[P]} & \textsc{[P]}
    & \textsc{[P]} & \textsc{[P]} \\[4pt]
$\text{CEV}(\text{E0 vs E1\_2L})$ (renter benchmark)
    & \textsc{[P]} & --- & --- & --- & --- & --- \\
\addlinespace
Mean $\bar{x}_B$ at $\ell{=}A$
    & \textsc{[P]} & \textsc{[P]} & \textsc{[P]} & \textsc{[P]}
    & \textsc{[P]} & \textsc{[P]} \\
\bottomrule
\end{tabular}
\vspace{4pt}
\begin{minipage}{\linewidth}
\small\textit{Notes:}
Baseline: $\gamma=5$, $\beta=0.96$, $\rho_{AB}=0.50$,
$p_{\text{reloc,work}}=0.06$, $\tau_{\text{sell}}=0.06$,
$\tau_{\text{buy}}=0.025$, $\tau_{\text{token}}=0.01$.
Grid: $N_W=15$, $N_Z=5$, $N_{x_{\text{prev}}}=3$ (coarse; full-grid
robustness at $N_W=40$, $N_Z=9$, $N_{x_{\text{prev}}}=5$ in
Appendix~\ref{app:numerical}).
\textsc{[P]} = placeholder pending server1 baseline runs.
\end{minipage}
\end{table}

The primary diagnostic for mechanism activation is $\bar{x}_B$ at $\ell{=}A$:
the mean portfolio share of location-$B$ tokens held by households residing
at location $A$.
Under E1\_2L, $\bar{x}_B = 0$ by admissibility (binary tenure forbids
non-occupancy-location ownership).
Under E2\_2L, a positive $\bar{x}_B > 0$ at $\ell{=}A$ indicates that
households voluntarily pre-accumulate location-$B$ tokens before relocation ---
the mechanism that is uniquely enabled by token portability and the 6D state
extension (Option~1).
This is the key discriminator from the v3 Option~3 approximation, under which
$\bar{x}_B = 0$ held at all mobility rates.

Figure~\ref{fig:lifecycle_profiles} illustrates the mechanism at the
lifecycle level: the age profiles of mean $x_A$ and $x_B$ holdings
under E1\_2L and E2\_2L.
The key prediction is that $\bar{x}_B$ rises during working years (ages 25--65),
when $p_{\text{reloc,work}} = 0.06$, and declines post-retirement when the
relocation probability falls to $p_{\text{reloc,ret}} = 0.02$.
This pattern --- confirmed or denied by the server1 policy-array output ---
constitutes the strongest visual evidence for the pre-buying hedge mechanism.

\begin{figure}[t]
\centering
%% PLACEHOLDER — replace with lifecycle profile plots once server1 policy arrays land.
%% See paper/exhibit_memos/fig1_lifecycle_profiles.md for full production spec.
\fbox{%
  \begin{minipage}{0.88\linewidth}
  \vspace{2.5cm}
  \centering
  {\sffamily\small [PLACEHOLDER: Figure~1 — Lifecycle token-holding profiles.
  Panel A (E1\_2L): mean $x_A$ (solid) and $x_B = 0$ (dashed) over ages 25--80.
  Panel B (E2\_2L): mean $x_A$ (solid) and mean $x_B > 0$ (dashed) over ages 25--80.
  Source: per-age mean of \texttt{xA\_policy} and \texttt{xB\_policy} arrays at $\ell{=}A$,
  averaged over feasible $(w, z, x_{\text{prev}})$ states.
  Script: \texttt{scripts/plot\_lifecycle\_profiles.py}.]}
  \vspace{2.5cm}
  \end{minipage}%
}
\caption{Lifecycle mean token-holding profiles under traditional ownership (E1\_2L,
         left panel) and tokenized ownership (E2\_2L, right panel) at location $A$.
         E2\_2L households pre-accumulate location-$B$ tokens during working years
         (positive $\bar{x}_B$), saving the buying cost $\tau_{\text{buy}}$ on
         anticipated future relocation.
         E1\_2L households hold $\bar{x}_B = 0$ at all ages by admissibility.
         \textsc{[P]} = placeholder pending server1 run.}
\label{fig:lifecycle_profiles}
\end{figure}

% ─────────────────────────────────────────────────────────────────────────────
\subsection{Channel Decomposition}
\label{sec:results:decomp}
% ─────────────────────────────────────────────────────────────────────────────

We decompose the total welfare gain using three regime runs.
Let E1\_2L$_{\text{NT}}$ denote the no-transaction-cost counterfactual
($\tau_{\text{sell}} = \tau_{\text{buy}} = 0$, binary tenure otherwise).
The decomposition is:
\begin{align}
    \underbrace{\text{CEV}(\text{E2\_2L vs E1\_2L})}_{\text{total}}
    &\;=\;
    \underbrace{\text{CEV}(\text{E1\_2L}_{\text{NT}} \text{ vs E1\_2L})}_{\text{avoided-tx}}
    \;+\;
    \underbrace{\text{CEV}(\text{E2\_2L vs E1\_2L}_{\text{NT}})}_{\text{maintained-hedge}}
    \;+\; \xi,
    \label{eq:decomp}
\end{align}
where $\xi$ is the cross-term.
This decomposition is pre-registered in \texttt{docs/welfare\_decomp\_v4.md}
and is exact up to the cross-term, which was empirically near zero in v3
($\xi = 0.021$\%) and is expected to remain small in v4.

\begin{table}[t]
\centering
\caption{Channel Decomposition of the Tokenization Welfare Gain
         (CEV, percentage points of lifetime consumption)}
\label{tab:decomp}
\small
\begin{tabular}{lccl}
\toprule
Channel & CEV (\%) & Share of total & Description \\
\midrule
Total: $\text{CEV}(\text{E2\_2L vs E1\_2L})$
    & \textsc{[P]} & 100\% & Headline tokenization gain \\
\addlinespace
\quad (i)~Avoided-tx
    & \textsc{[P]} & \textsc{[P]}\%
    & $\tau_{\text{sell}} + \tau_{\text{buy}}$ burden of E1\_2L moves \\
\qquad of which $\tau_{\text{sell}}$
    & \textsc{[P]} & \textsc{[P]}\%
    & NAR selling commission (6\%) \\
\qquad of which $\tau_{\text{buy}}$
    & \textsc{[P]} & \textsc{[P]}\%
    & Buying closing costs (2.5\%) \\[4pt]
\quad (ii)~Maintained-hedge
    & \textsc{[P]} & \textsc{[P]}\%
    & Continuous-x + pre-buying hedge vs E1\_2L$_\text{NT}$ \\[4pt]
\quad (iii)~Cross-term $\xi$
    & \textsc{[P]} & \textsc{[P]}\%
    & Interaction; expected $\approx 0$ \\
\midrule
v3 Option~3 (no state ext.)
    & $+4.255$\% & (reference) & Previous fire's upper-bound approx \\
\quad of which avoided-tx
    & $+0.816$\% & 19.2\% & From v3 full-grid run \\
\quad of which continuous-x
    & $+3.411$\% & 80.2\% & Liu (2021) territory \\
\quad of which pre-buying hedge
    & $\approx 0$\% & $\approx 0$\% & Absent without state ext. (Option~3 limit) \\
\bottomrule
\end{tabular}
\vspace{4pt}
\begin{minipage}{\linewidth}
\small\textit{Notes:}
Upper panel: v4 Option~1 results (\textsc{[P]} = placeholder).
Lower panel: v3 Option~3 reference numbers (confirmed from server1 run,
commit 186da13).
The maintained-hedge channel in v4 encompasses both the continuous-x
subchannel (present in v3) and the new pre-buying hedge subchannel
enabled by the 6D state; a fourth regime decomposition
($\text{E2\_2L}_{\text{NT}}$) can further separate the two if H1 is confirmed.
\end{minipage}
\end{table}

The maintained-hedge channel in v4 captures two sub-mechanisms:
\begin{enumerate}
    \item \emph{Continuous-$x$ sub-channel} (present in v3): continuous
    fractional ownership of the occupied unit avoids the binary kink at
    $x_\ell = 1$, delivering a Liu~(2021)-style indivisibility relaxation.
    \item \emph{Pre-buying hedge sub-channel} (new in v4): by tracking
    $(x_{A,t-1}, x_{B,t-1})$ in the state, households face a per-period
    buying cost $\tau_{\text{buy}} \cdot \max(\Delta x, 0)$ and can
    pre-accumulate location-$B$ tokens at location $A$, arriving at $B$
    with $x_{B,t-1} > 0$ and avoiding the lump-sum purchase cost.
    The expected per-period premium per unit held is
    $p_{\text{reloc}} \times \tau_{\text{buy}} = 0.06 \times 0.025 = 0.15$\%
    (see Appendix~\ref{app:proof}).
\end{enumerate}

% ─────────────────────────────────────────────────────────────────────────────
\subsection{Falsification Tests}
\label{sec:results:falsification}
% ─────────────────────────────────────────────────────────────────────────────

We conduct three pre-registered falsification tests designed to identify
the pre-buying hedge channel.
These tests were pre-registered in \texttt{docs/welfare\_decomp\_v4.md}
before server1 runs were executed.
Crucially, the analogous tests \emph{failed} under v3 Option~3 (which lacked
the state extension): $\bar{x}_B$ remained zero at all mobility rates,
and the CEV was insensitive to $p_{\text{reloc}}$ and $\rho_{AB}$.

\begin{table}[t]
\centering
\caption{Pre-Registered Falsification Tests
         (Deviations from Baseline)}
\label{tab:falsification}
\small
\begin{tabular}{llcccc}
\toprule
Test & Counterfactual & $\text{CEV}$ & $\bar{x}_B$ at $\ell{=}A$
     & Pass criterion & Result \\
\midrule
(r) Relocation disabled
    & $p_{\text{reloc}} = 0$
    & \textsc{[P]}
    & \textsc{[P]}
    & $\bar{x}_B \to 0$;\; $\Delta\text{CEV} < -0.5$pp
    & \textsc{[P]} \\[4pt]
(m) Perfect correlation
    & $\rho_{AB} = 0.95$
    & \textsc{[P]}
    & \textsc{[P]}
    & $\bar{x}_B$ decreases;\; $\Delta\text{CEV} < 0$
    & \textsc{[P]} \\[4pt]
(q) No buying cost
    & $\tau_{\text{buy}} = 0$
    & \textsc{[P]}
    & \textsc{[P]}
    & Hedge channel $\to 0$;\; $\bar{x}_B \to 0$
    & \textsc{[P]} \\
\bottomrule
\end{tabular}
\vspace{4pt}
\begin{minipage}{\linewidth}
\small\textit{Notes:}
All tests compare to the baseline E2\_2L v4 run.
Pass criterion stated prior to server1 runs.
``Hedge channel'' in test (q) is
$\text{CEV}(\text{E2\_2L}_{\tau=0} \text{ vs E1\_2L}_{\text{NT}}) - \text{CEV}(\text{E2\_2L}_{\text{baseline}} \text{ vs E1\_2L}_{\text{NT}})$.
\textsc{[P]} = placeholder.
\end{minipage}
\end{table}

Test (r) shuts off the relocation shock entirely.
Without future relocation, there is no value to pre-accumulating
location-$B$ tokens while at $A$; the pre-buying hedge channel
should collapse.
If test (r) passes --- i.e., $\bar{x}_B \to 0$ when $p_{\text{reloc}} = 0$
--- it provides strong evidence that the positive $\bar{x}_B$ in the
baseline is driven by the pre-buying motive and not by a model artefact.

Test (m) raises the cross-location correlation to $\rho_{AB} = 0.95$,
nearly eliminating diversification benefit.
When location-$A$ and location-$B$ price processes are nearly identical,
holding $x_B$ at $\ell = A$ provides little hedge against diverging location returns.
The CEV and $\bar{x}_B$ should both decrease relative to the baseline.

Test (q) eliminates the buying cost ($\tau_{\text{buy}} = 0$), which
removes the financial incentive to pre-accumulate $x_B$ (there is no lump-sum
arrival cost to avoid).
The pre-buying hedge channel should collapse to zero.

% ─────────────────────────────────────────────────────────────────────────────
\subsection{Sensitivity Analysis}
\label{sec:results:sensitivity}
% ─────────────────────────────────────────────────────────────────────────────

Figure~\ref{fig:sensitivity_heatmap} depicts the total welfare gain
$\text{CEV}(\text{E2\_2L vs E1\_2L})$ as a function of the two primary
mechanism parameters: the cross-location correlation $\rho_{AB}$ and
the relocation probability $p_{\text{reloc}}$.
The full sweep design follows the pre-registered plan in
\texttt{docs/sensitivity\_grid\_v4.md}; only a subset is shown here.

\begin{figure}[t]
\centering
%% PLACEHOLDER — replace with actual heatmap once sweep results available
\fbox{%
  \begin{minipage}{0.72\linewidth}
  \vspace{2.5cm}
  \centering
  {\sffamily\small [PLACEHOLDER: Figure~2 — CEV sensitivity heatmap.
  X-axis: $\rho_{AB} \in \{0, 0.25, 0.50, 0.75, 0.95\}$.
  Y-axis: $p_{\text{reloc}} \in \{0, 0.02, 0.06, 0.12\}$.
  Heat colour: total CEV(\%), from dark (low) to bright (high).
  Source: sweep outputs from \texttt{scripts/sweep\_rhoAB.sh} and
  \texttt{scripts/sweep\_prelocate.sh} once server1 runs complete.]}
  \vspace{2.5cm}
  \end{minipage}%
}
\caption{Welfare sensitivity to cross-location correlation and relocation rate.
         Each cell reports $\text{CEV}(\text{E2\_2L vs E1\_2L})$ at
         $(\rho_{AB},\, p_{\text{reloc}})$.
         The mechanism prediction is a positive gradient in $p_{\text{reloc}}$
         (higher mobility $\Rightarrow$ higher hedge value) and a negative
         gradient in $\rho_{AB}$ (more correlated locations $\Rightarrow$
         smaller diversification benefit).}
\label{fig:sensitivity_heatmap}
\end{figure}

Table~\ref{tab:sensitivity} reports selected cross-sections of the
sensitivity grid for other key parameters.

\begin{table}[t]
\centering
\caption{Sensitivity of Headline CEV to Key Parameters
         (Total CEV, \%; baseline value on diagonal)}
\label{tab:sensitivity}
\small
\begin{tabular}{lcccccc}
\toprule
Parameter & & \multicolumn{5}{c}{Values} \\
\midrule
Risk aversion $\gamma$
    & & 3 & \textbf{5} & 8 & & \\
\quad $\text{CEV}(\text{E2\_2L vs E1\_2L})$
    & & \textsc{[P]} & \textsc{[P]} & \textsc{[P]} & & \\
\addlinespace
Buying cost $\tau_{\text{buy}}$
    & & 0\% & 1\% & \textbf{2.5\%} & 4\% & 6\% \\
\quad $\text{CEV}(\text{E2\_2L vs E1\_2L})$
    & & \textsc{[P]} & \textsc{[P]} & \textsc{[P]} & \textsc{[P]} & \textsc{[P]} \\
\addlinespace
Relocation rate $p_{\text{reloc,work}}$
    & & 0 & 0.02 & \textbf{0.06} & 0.12 & \\
\quad $\text{CEV}(\text{E2\_2L vs E1\_2L})$
    & & \textsc{[P]} & \textsc{[P]} & \textsc{[P]} & \textsc{[P]} & \\
\quad $\bar{x}_B$ at $\ell{=}A$
    & & \textsc{[P]} & \textsc{[P]} & \textsc{[P]} & \textsc{[P]} & \\
\addlinespace
Cross-location corr.\ $\rho_{AB}$
    & & 0 & 0.25 & \textbf{0.50} & 0.75 & 0.95 \\
\quad $\text{CEV}(\text{E2\_2L vs E1\_2L})$
    & & \textsc{[P]} & \textsc{[P]} & \textsc{[P]} & \textsc{[P]} & \textsc{[P]} \\
\addlinespace
\multicolumn{7}{l}{\textit{Asymmetric robustness (solver v4 extension, fire 22)}} \\
Location-$B$ return premium $\Delta\mu_B$
    & & $-1\%$ & \textbf{0\%} & $+1\%$ & & \\
\quad $\text{CEV}(\text{E2\_2L vs E1\_2L})$
    & & \textsc{[P]} & \textsc{[P]} & \textsc{[P]} & & \\
\quad $\bar{x}_B$ at $\ell{=}A$
    & & \textsc{[P]} & \textsc{[P]} & \textsc{[P]} & & \\[2pt]
Directional mobility ($p_{A\to B}$, $p_{B\to A}$)
    & & (0.06, 0.06) & (0.10, 0.03) & (0.10, 0.10) & & \\
\quad $\text{CEV}(\text{E2\_2L vs E1\_2L})$
    & & \textsc{[P]} & \textsc{[P]} & \textsc{[P]} & & \\
\addlinespace
\multicolumn{7}{l}{\textit{Mortgage activation (LTV constraint)}} \\
LTV ceiling $\theta_{\max}$
    & & \textbf{0} & 0.50 & 0.80 & & \\
\quad $\text{CEV}(\text{E2\_2L vs E1\_2L})$
    & & \textsc{[P]} & \textsc{[P]} & \textsc{[P]} & & \\
\bottomrule
\end{tabular}
\vspace{4pt}
\begin{minipage}{\linewidth}
\small\textit{Notes:}
Bold entries mark the baseline value.
One parameter varies at a time; all others held at baseline.
Sweep scripts: \texttt{scripts/sweep\_rhoAB.sh},
\texttt{scripts/sweep\_prelocate.sh}, \texttt{scripts/sweep\_txcost.sh},
\texttt{scripts/sweep\_asymmetric.sh} (location-$B$ return premium and directional mobility),
\texttt{scripts/sweep\_mortgage.sh} (LTV).
$\Delta\mu_B = \mu_{h,B} - \mu_{h,A}$: the excess Jensen-corrected log-mean return of
location $B$ relative to $A$.
Positive $\Delta\mu_B$ strengthens the pre-buying motive for $x_B$ by adding a return
premium to the tau\_buy-saving incentive.
\textsc{[P]} = placeholder pending server1 sweep runs.
\end{minipage}
\end{table}

Two monotonicity predictions follow from the model mechanism.
First, for $\tau_{\text{buy}}$: as the buying cost rises, the pre-accumulation
motive strengthens and $\bar{x}_B$ should increase, lifting the hedge sub-channel.
Second, for $p_{\text{reloc}}$: higher mobility raises the expected value of
pre-accumulated $x_B$ holdings; at $p_{\text{reloc}} = 0$ the hedge channel
collapses (falsification test r).
If these monotonicity patterns hold in the sweep results, they constitute
additional evidence for the pre-buying hedge mechanism.

% ─────────────────────────────────────────────────────────────────────────────
\subsection{Summary of Mechanism Evidence}
\label{sec:results:summary}
% ─────────────────────────────────────────────────────────────────────────────

The results establish \textsc{[P]} as the headline lifetime welfare gain from
tokenization under the baseline calibration.
The gain decomposes additively into an avoided-transaction-cost channel
(\textsc{[P]}\%) and a maintained-hedge channel (\textsc{[P]}\%), with
a near-zero cross-term.
The maintained-hedge channel is itself divisible into the continuous-$x$
sub-channel present in prior work (Liu, 2021) and the novel pre-buying hedge
sub-channel unlocked by the 6D state extension.

The three falsification tests provide mechanism identification.
Under all three counterfactuals --- zero relocation, high cross-location
correlation, zero buying cost --- the pre-buying hedge channel should
attenuate or collapse.
Confirmation of these patterns will establish that the welfare gain reported
in Table~\ref{tab:cev_baseline} reflects the model's structural mechanism
rather than a calibration artefact.
\textsc{[Fill in pass/fail outcomes from output/diagnostics/p6\_option1\_decomposition.md.]}

%              Requires amsmath, booktabs, multirow, siunitx, graphicx, caption.

% ─────────────────────────────────────────────────────────────────────────────
\section{Results}
\label{sec:results}
% ─────────────────────────────────────────────────────────────────────────────

We report the baseline welfare results in Section~\ref{sec:results:baseline},
the channel decomposition in Section~\ref{sec:results:decomp}, the
pre-registered falsification tests in Section~\ref{sec:results:falsification},
and sensitivity analysis in Section~\ref{sec:results:sensitivity}.

All welfare gains are expressed as consumption equivalent variations (CEV) at the
initial state
$\mathbf{s}_0 = (t{=}1,\,w_{\text{mid}},\,z_{\text{mid}},\,\ell{=}A,\,x_{A,0}{=}0,\,x_{B,0}{=}0)$:
\begin{equation}
    \text{CEV}(R_a \text{ vs } R_b) \;=\;
    \left(\frac{V_1^{R_a}(\mathbf{s}_0)}{V_1^{R_b}(\mathbf{s}_0)}\right)^{1/(1-\gamma)}
    - 1,
    \label{eq:cev}
\end{equation}
where $V_1^{R}(\mathbf{s}_0)$ is the time-1 value function under regime $R$
evaluated at $\mathbf{s}_0$.%
\footnote{Both value functions are negative under CRRA with finite wealth,
so the ratio is positive and the formula is well-defined for $\gamma \neq 1$.
The CEV measures the fraction of lifetime consumption in regime $R_b$ that the
household would give up to gain permanent access to regime $R_a$.}
Appendix~\ref{app:numerical} reports convergence diagnostics and full-grid
robustness.

% ─────────────────────────────────────────────────────────────────────────────
\subsection{Baseline Welfare Results}
\label{sec:results:baseline}
% ─────────────────────────────────────────────────────────────────────────────

The headline welfare result is reported in Table~\ref{tab:cev_baseline}.
The first row gives $\text{CEV}(\text{E2\_2L vs E1\_2L})$, the lifetime
welfare value of tokenization relative to traditional binary homeownership.
Rows 2--4 decompose this gain into three channels (detailed in
Section~\ref{sec:results:decomp}).
Row 5 reports $\text{CEV}(\text{E0 vs E1\_2L})$ as a sanity-check benchmark:
the welfare cost of pure renting relative to traditional ownership.

\begin{table}[t]
\centering
\caption{Welfare Value of Tokenization: Baseline and Robustness Panel
         (CEV, percentage points of lifetime consumption)}
\label{tab:cev_baseline}
\small
\begin{tabular}{lcccccc}
\toprule
 & Baseline & $\rho_{AB}=0$ & $\rho_{AB}=0.50$ & $\rho_{AB}=0.95$
 & $p_{\text{reloc}}=0$ & $p_{\text{reloc}}=0.12$ \\
\midrule
$\text{CEV}(\text{E2\_2L vs E1\_2L})$
    & \textsc{[P]} & \textsc{[P]} & \textsc{[P]} & \textsc{[P]}
    & \textsc{[P]} & \textsc{[P]} \\[2pt]
\quad Avoided-tx channel
    & \textsc{[P]} & \textsc{[P]} & \textsc{[P]} & \textsc{[P]}
    & \textsc{[P]} & \textsc{[P]} \\
\quad Maintained-hedge channel
    & \textsc{[P]} & \textsc{[P]} & \textsc{[P]} & \textsc{[P]}
    & \textsc{[P]} & \textsc{[P]} \\
\quad Cross-term $\xi$
    & \textsc{[P]} & \textsc{[P]} & \textsc{[P]} & \textsc{[P]}
    & \textsc{[P]} & \textsc{[P]} \\[4pt]
$\text{CEV}(\text{E0 vs E1\_2L})$ (renter benchmark)
    & \textsc{[P]} & --- & --- & --- & --- & --- \\
\addlinespace
Mean $\bar{x}_B$ at $\ell{=}A$
    & \textsc{[P]} & \textsc{[P]} & \textsc{[P]} & \textsc{[P]}
    & \textsc{[P]} & \textsc{[P]} \\
\bottomrule
\end{tabular}
\vspace{4pt}
\begin{minipage}{\linewidth}
\small\textit{Notes:}
Baseline: $\gamma=5$, $\beta=0.96$, $\rho_{AB}=0.50$,
$p_{\text{reloc,work}}=0.06$, $\tau_{\text{sell}}=0.06$,
$\tau_{\text{buy}}=0.025$, $\tau_{\text{token}}=0.01$.
Grid: $N_W=15$, $N_Z=5$, $N_{x_{\text{prev}}}=3$ (coarse; full-grid
robustness at $N_W=40$, $N_Z=9$, $N_{x_{\text{prev}}}=5$ in
Appendix~\ref{app:numerical}).
\textsc{[P]} = placeholder pending server1 baseline runs.
\end{minipage}
\end{table}

The primary diagnostic for mechanism activation is $\bar{x}_B$ at $\ell{=}A$:
the mean portfolio share of location-$B$ tokens held by households residing
at location $A$.
Under E1\_2L, $\bar{x}_B = 0$ by admissibility (binary tenure forbids
non-occupancy-location ownership).
Under E2\_2L, a positive $\bar{x}_B > 0$ at $\ell{=}A$ indicates that
households voluntarily pre-accumulate location-$B$ tokens before relocation ---
the mechanism that is uniquely enabled by token portability and the 6D state
extension (Option~1).
This is the key discriminator from the v3 Option~3 approximation, under which
$\bar{x}_B = 0$ held at all mobility rates.

Figure~\ref{fig:lifecycle_profiles} illustrates the mechanism at the
lifecycle level: the age profiles of mean $x_A$ and $x_B$ holdings
under E1\_2L and E2\_2L.
The key prediction is that $\bar{x}_B$ rises during working years (ages 25--65),
when $p_{\text{reloc,work}} = 0.06$, and declines post-retirement when the
relocation probability falls to $p_{\text{reloc,ret}} = 0.02$.
This pattern --- confirmed or denied by the server1 policy-array output ---
constitutes the strongest visual evidence for the pre-buying hedge mechanism.

\begin{figure}[t]
\centering
%% PLACEHOLDER — replace with lifecycle profile plots once server1 policy arrays land.
%% See paper/exhibit_memos/fig1_lifecycle_profiles.md for full production spec.
\fbox{%
  \begin{minipage}{0.88\linewidth}
  \vspace{2.5cm}
  \centering
  {\sffamily\small [PLACEHOLDER: Figure~1 — Lifecycle token-holding profiles.
  Panel A (E1\_2L): mean $x_A$ (solid) and $x_B = 0$ (dashed) over ages 25--80.
  Panel B (E2\_2L): mean $x_A$ (solid) and mean $x_B > 0$ (dashed) over ages 25--80.
  Source: per-age mean of \texttt{xA\_policy} and \texttt{xB\_policy} arrays at $\ell{=}A$,
  averaged over feasible $(w, z, x_{\text{prev}})$ states.
  Script: \texttt{scripts/plot\_lifecycle\_profiles.py}.]}
  \vspace{2.5cm}
  \end{minipage}%
}
\caption{Lifecycle mean token-holding profiles under traditional ownership (E1\_2L,
         left panel) and tokenized ownership (E2\_2L, right panel) at location $A$.
         E2\_2L households pre-accumulate location-$B$ tokens during working years
         (positive $\bar{x}_B$), saving the buying cost $\tau_{\text{buy}}$ on
         anticipated future relocation.
         E1\_2L households hold $\bar{x}_B = 0$ at all ages by admissibility.
         \textsc{[P]} = placeholder pending server1 run.}
\label{fig:lifecycle_profiles}
\end{figure}

% ─────────────────────────────────────────────────────────────────────────────
\subsection{Channel Decomposition}
\label{sec:results:decomp}
% ─────────────────────────────────────────────────────────────────────────────

We decompose the total welfare gain using three regime runs.
Let E1\_2L$_{\text{NT}}$ denote the no-transaction-cost counterfactual
($\tau_{\text{sell}} = \tau_{\text{buy}} = 0$, binary tenure otherwise).
The decomposition is:
\begin{align}
    \underbrace{\text{CEV}(\text{E2\_2L vs E1\_2L})}_{\text{total}}
    &\;=\;
    \underbrace{\text{CEV}(\text{E1\_2L}_{\text{NT}} \text{ vs E1\_2L})}_{\text{avoided-tx}}
    \;+\;
    \underbrace{\text{CEV}(\text{E2\_2L vs E1\_2L}_{\text{NT}})}_{\text{maintained-hedge}}
    \;+\; \xi,
    \label{eq:decomp}
\end{align}
where $\xi$ is the cross-term.
This decomposition is pre-registered in \texttt{docs/welfare\_decomp\_v4.md}
and is exact up to the cross-term, which was empirically near zero in v3
($\xi = 0.021$\%) and is expected to remain small in v4.

\begin{table}[t]
\centering
\caption{Channel Decomposition of the Tokenization Welfare Gain
         (CEV, percentage points of lifetime consumption)}
\label{tab:decomp}
\small
\begin{tabular}{lccl}
\toprule
Channel & CEV (\%) & Share of total & Description \\
\midrule
Total: $\text{CEV}(\text{E2\_2L vs E1\_2L})$
    & \textsc{[P]} & 100\% & Headline tokenization gain \\
\addlinespace
\quad (i)~Avoided-tx
    & \textsc{[P]} & \textsc{[P]}\%
    & $\tau_{\text{sell}} + \tau_{\text{buy}}$ burden of E1\_2L moves \\
\qquad of which $\tau_{\text{sell}}$
    & \textsc{[P]} & \textsc{[P]}\%
    & NAR selling commission (6\%) \\
\qquad of which $\tau_{\text{buy}}$
    & \textsc{[P]} & \textsc{[P]}\%
    & Buying closing costs (2.5\%) \\[4pt]
\quad (ii)~Maintained-hedge
    & \textsc{[P]} & \textsc{[P]}\%
    & Continuous-x + pre-buying hedge vs E1\_2L$_\text{NT}$ \\[4pt]
\quad (iii)~Cross-term $\xi$
    & \textsc{[P]} & \textsc{[P]}\%
    & Interaction; expected $\approx 0$ \\
\midrule
v3 Option~3 (no state ext.)
    & $+4.255$\% & (reference) & Previous fire's upper-bound approx \\
\quad of which avoided-tx
    & $+0.816$\% & 19.2\% & From v3 full-grid run \\
\quad of which continuous-x
    & $+3.411$\% & 80.2\% & Liu (2021) territory \\
\quad of which pre-buying hedge
    & $\approx 0$\% & $\approx 0$\% & Absent without state ext. (Option~3 limit) \\
\bottomrule
\end{tabular}
\vspace{4pt}
\begin{minipage}{\linewidth}
\small\textit{Notes:}
Upper panel: v4 Option~1 results (\textsc{[P]} = placeholder).
Lower panel: v3 Option~3 reference numbers (confirmed from server1 run,
commit 186da13).
The maintained-hedge channel in v4 encompasses both the continuous-x
subchannel (present in v3) and the new pre-buying hedge subchannel
enabled by the 6D state; a fourth regime decomposition
($\text{E2\_2L}_{\text{NT}}$) can further separate the two if H1 is confirmed.
\end{minipage}
\end{table}

The maintained-hedge channel in v4 captures two sub-mechanisms:
\begin{enumerate}
    \item \emph{Continuous-$x$ sub-channel} (present in v3): continuous
    fractional ownership of the occupied unit avoids the binary kink at
    $x_\ell = 1$, delivering a Liu~(2021)-style indivisibility relaxation.
    \item \emph{Pre-buying hedge sub-channel} (new in v4): by tracking
    $(x_{A,t-1}, x_{B,t-1})$ in the state, households face a per-period
    buying cost $\tau_{\text{buy}} \cdot \max(\Delta x, 0)$ and can
    pre-accumulate location-$B$ tokens at location $A$, arriving at $B$
    with $x_{B,t-1} > 0$ and avoiding the lump-sum purchase cost.
    The expected per-period premium per unit held is
    $p_{\text{reloc}} \times \tau_{\text{buy}} = 0.06 \times 0.025 = 0.15$\%
    (see Appendix~\ref{app:proof}).
\end{enumerate}

% ─────────────────────────────────────────────────────────────────────────────
\subsection{Falsification Tests}
\label{sec:results:falsification}
% ─────────────────────────────────────────────────────────────────────────────

We conduct three pre-registered falsification tests designed to identify
the pre-buying hedge channel.
These tests were pre-registered in \texttt{docs/welfare\_decomp\_v4.md}
before server1 runs were executed.
Crucially, the analogous tests \emph{failed} under v3 Option~3 (which lacked
the state extension): $\bar{x}_B$ remained zero at all mobility rates,
and the CEV was insensitive to $p_{\text{reloc}}$ and $\rho_{AB}$.

\begin{table}[t]
\centering
\caption{Pre-Registered Falsification Tests
         (Deviations from Baseline)}
\label{tab:falsification}
\small
\begin{tabular}{llcccc}
\toprule
Test & Counterfactual & $\text{CEV}$ & $\bar{x}_B$ at $\ell{=}A$
     & Pass criterion & Result \\
\midrule
(r) Relocation disabled
    & $p_{\text{reloc}} = 0$
    & \textsc{[P]}
    & \textsc{[P]}
    & $\bar{x}_B \to 0$;\; $\Delta\text{CEV} < -0.5$pp
    & \textsc{[P]} \\[4pt]
(m) Perfect correlation
    & $\rho_{AB} = 0.95$
    & \textsc{[P]}
    & \textsc{[P]}
    & $\bar{x}_B$ decreases;\; $\Delta\text{CEV} < 0$
    & \textsc{[P]} \\[4pt]
(q) No buying cost
    & $\tau_{\text{buy}} = 0$
    & \textsc{[P]}
    & \textsc{[P]}
    & Hedge channel $\to 0$;\; $\bar{x}_B \to 0$
    & \textsc{[P]} \\
\bottomrule
\end{tabular}
\vspace{4pt}
\begin{minipage}{\linewidth}
\small\textit{Notes:}
All tests compare to the baseline E2\_2L v4 run.
Pass criterion stated prior to server1 runs.
``Hedge channel'' in test (q) is
$\text{CEV}(\text{E2\_2L}_{\tau=0} \text{ vs E1\_2L}_{\text{NT}}) - \text{CEV}(\text{E2\_2L}_{\text{baseline}} \text{ vs E1\_2L}_{\text{NT}})$.
\textsc{[P]} = placeholder.
\end{minipage}
\end{table}

Test (r) shuts off the relocation shock entirely.
Without future relocation, there is no value to pre-accumulating
location-$B$ tokens while at $A$; the pre-buying hedge channel
should collapse.
If test (r) passes --- i.e., $\bar{x}_B \to 0$ when $p_{\text{reloc}} = 0$
--- it provides strong evidence that the positive $\bar{x}_B$ in the
baseline is driven by the pre-buying motive and not by a model artefact.

Test (m) raises the cross-location correlation to $\rho_{AB} = 0.95$,
nearly eliminating diversification benefit.
When location-$A$ and location-$B$ price processes are nearly identical,
holding $x_B$ at $\ell = A$ provides little hedge against diverging location returns.
The CEV and $\bar{x}_B$ should both decrease relative to the baseline.

Test (q) eliminates the buying cost ($\tau_{\text{buy}} = 0$), which
removes the financial incentive to pre-accumulate $x_B$ (there is no lump-sum
arrival cost to avoid).
The pre-buying hedge channel should collapse to zero.

% ─────────────────────────────────────────────────────────────────────────────
\subsection{Sensitivity Analysis}
\label{sec:results:sensitivity}
% ─────────────────────────────────────────────────────────────────────────────

Figure~\ref{fig:sensitivity_heatmap} depicts the total welfare gain
$\text{CEV}(\text{E2\_2L vs E1\_2L})$ as a function of the two primary
mechanism parameters: the cross-location correlation $\rho_{AB}$ and
the relocation probability $p_{\text{reloc}}$.
The full sweep design follows the pre-registered plan in
\texttt{docs/sensitivity\_grid\_v4.md}; only a subset is shown here.

\begin{figure}[t]
\centering
%% PLACEHOLDER — replace with actual heatmap once sweep results available
\fbox{%
  \begin{minipage}{0.72\linewidth}
  \vspace{2.5cm}
  \centering
  {\sffamily\small [PLACEHOLDER: Figure~2 — CEV sensitivity heatmap.
  X-axis: $\rho_{AB} \in \{0, 0.25, 0.50, 0.75, 0.95\}$.
  Y-axis: $p_{\text{reloc}} \in \{0, 0.02, 0.06, 0.12\}$.
  Heat colour: total CEV(\%), from dark (low) to bright (high).
  Source: sweep outputs from \texttt{scripts/sweep\_rhoAB.sh} and
  \texttt{scripts/sweep\_prelocate.sh} once server1 runs complete.]}
  \vspace{2.5cm}
  \end{minipage}%
}
\caption{Welfare sensitivity to cross-location correlation and relocation rate.
         Each cell reports $\text{CEV}(\text{E2\_2L vs E1\_2L})$ at
         $(\rho_{AB},\, p_{\text{reloc}})$.
         The mechanism prediction is a positive gradient in $p_{\text{reloc}}$
         (higher mobility $\Rightarrow$ higher hedge value) and a negative
         gradient in $\rho_{AB}$ (more correlated locations $\Rightarrow$
         smaller diversification benefit).}
\label{fig:sensitivity_heatmap}
\end{figure}

Table~\ref{tab:sensitivity} reports selected cross-sections of the
sensitivity grid for other key parameters.

\begin{table}[t]
\centering
\caption{Sensitivity of Headline CEV to Key Parameters
         (Total CEV, \%; baseline value on diagonal)}
\label{tab:sensitivity}
\small
\begin{tabular}{lcccccc}
\toprule
Parameter & & \multicolumn{5}{c}{Values} \\
\midrule
Risk aversion $\gamma$
    & & 3 & \textbf{5} & 8 & & \\
\quad $\text{CEV}(\text{E2\_2L vs E1\_2L})$
    & & \textsc{[P]} & \textsc{[P]} & \textsc{[P]} & & \\
\addlinespace
Buying cost $\tau_{\text{buy}}$
    & & 0\% & 1\% & \textbf{2.5\%} & 4\% & 6\% \\
\quad $\text{CEV}(\text{E2\_2L vs E1\_2L})$
    & & \textsc{[P]} & \textsc{[P]} & \textsc{[P]} & \textsc{[P]} & \textsc{[P]} \\
\addlinespace
Relocation rate $p_{\text{reloc,work}}$
    & & 0 & 0.02 & \textbf{0.06} & 0.12 & \\
\quad $\text{CEV}(\text{E2\_2L vs E1\_2L})$
    & & \textsc{[P]} & \textsc{[P]} & \textsc{[P]} & \textsc{[P]} & \\
\quad $\bar{x}_B$ at $\ell{=}A$
    & & \textsc{[P]} & \textsc{[P]} & \textsc{[P]} & \textsc{[P]} & \\
\addlinespace
Cross-location corr.\ $\rho_{AB}$
    & & 0 & 0.25 & \textbf{0.50} & 0.75 & 0.95 \\
\quad $\text{CEV}(\text{E2\_2L vs E1\_2L})$
    & & \textsc{[P]} & \textsc{[P]} & \textsc{[P]} & \textsc{[P]} & \textsc{[P]} \\
\addlinespace
\multicolumn{7}{l}{\textit{Asymmetric robustness (solver v4 extension, fire 22)}} \\
Location-$B$ return premium $\Delta\mu_B$
    & & $-1\%$ & \textbf{0\%} & $+1\%$ & & \\
\quad $\text{CEV}(\text{E2\_2L vs E1\_2L})$
    & & \textsc{[P]} & \textsc{[P]} & \textsc{[P]} & & \\
\quad $\bar{x}_B$ at $\ell{=}A$
    & & \textsc{[P]} & \textsc{[P]} & \textsc{[P]} & & \\[2pt]
Directional mobility ($p_{A\to B}$, $p_{B\to A}$)
    & & (0.06, 0.06) & (0.10, 0.03) & (0.10, 0.10) & & \\
\quad $\text{CEV}(\text{E2\_2L vs E1\_2L})$
    & & \textsc{[P]} & \textsc{[P]} & \textsc{[P]} & & \\
\addlinespace
\multicolumn{7}{l}{\textit{Mortgage activation (LTV constraint)}} \\
LTV ceiling $\theta_{\max}$
    & & \textbf{0} & 0.50 & 0.80 & & \\
\quad $\text{CEV}(\text{E2\_2L vs E1\_2L})$
    & & \textsc{[P]} & \textsc{[P]} & \textsc{[P]} & & \\
\bottomrule
\end{tabular}
\vspace{4pt}
\begin{minipage}{\linewidth}
\small\textit{Notes:}
Bold entries mark the baseline value.
One parameter varies at a time; all others held at baseline.
Sweep scripts: \texttt{scripts/sweep\_rhoAB.sh},
\texttt{scripts/sweep\_prelocate.sh}, \texttt{scripts/sweep\_txcost.sh},
\texttt{scripts/sweep\_asymmetric.sh} (location-$B$ return premium and directional mobility),
\texttt{scripts/sweep\_mortgage.sh} (LTV).
$\Delta\mu_B = \mu_{h,B} - \mu_{h,A}$: the excess Jensen-corrected log-mean return of
location $B$ relative to $A$.
Positive $\Delta\mu_B$ strengthens the pre-buying motive for $x_B$ by adding a return
premium to the tau\_buy-saving incentive.
\textsc{[P]} = placeholder pending server1 sweep runs.
\end{minipage}
\end{table}

Two monotonicity predictions follow from the model mechanism.
First, for $\tau_{\text{buy}}$: as the buying cost rises, the pre-accumulation
motive strengthens and $\bar{x}_B$ should increase, lifting the hedge sub-channel.
Second, for $p_{\text{reloc}}$: higher mobility raises the expected value of
pre-accumulated $x_B$ holdings; at $p_{\text{reloc}} = 0$ the hedge channel
collapses (falsification test r).
If these monotonicity patterns hold in the sweep results, they constitute
additional evidence for the pre-buying hedge mechanism.

% ─────────────────────────────────────────────────────────────────────────────
\subsection{Summary of Mechanism Evidence}
\label{sec:results:summary}
% ─────────────────────────────────────────────────────────────────────────────

The results establish \textsc{[P]} as the headline lifetime welfare gain from
tokenization under the baseline calibration.
The gain decomposes additively into an avoided-transaction-cost channel
(\textsc{[P]}\%) and a maintained-hedge channel (\textsc{[P]}\%), with
a near-zero cross-term.
The maintained-hedge channel is itself divisible into the continuous-$x$
sub-channel present in prior work (Liu, 2021) and the novel pre-buying hedge
sub-channel unlocked by the 6D state extension.

The three falsification tests provide mechanism identification.
Under all three counterfactuals --- zero relocation, high cross-location
correlation, zero buying cost --- the pre-buying hedge channel should
attenuate or collapse.
Confirmation of these patterns will establish that the welfare gain reported
in Table~\ref{tab:cev_baseline} reflects the model's structural mechanism
rather than a calibration artefact.
\textsc{[Fill in pass/fail outcomes from output/diagnostics/p6\_option1\_decomposition.md.]}

% ─────────────────────────────────────────────────────────────────────────────

\clearpage

% ─────────────────────────────────────────────────────────────────────────────
% s5_discussion.tex — Discussion
% Paper: "Tokenized Housing and Lifecycle Portfolio Choice:
%         A Decoupling of Location from Housing Exposure"
%
% Source: paper/outline_v4.md §5 (sections 5.1–5.3)
%         docs/welfare_decomp_v4.md §5 (literature comparison table)
%         paper/outline_v4.md Appendix C (hedge-channel sign conditions)
%
% Status: SKELETON — [PLACEHOLDER] marks require server1 numerical results.
%         Structural arguments are complete; fill CEV magnitudes once
%         output/diagnostics/p6_option1_*.json land.
%
% Compilation: % s5_discussion.tex — Discussion
% Paper: "Tokenized Housing and Lifecycle Portfolio Choice:
%         A Decoupling of Location from Housing Exposure"
%
% Source: paper/outline_v4.md §5 (sections 5.1–5.3)
%         docs/welfare_decomp_v4.md §5 (literature comparison table)
%         paper/outline_v4.md Appendix C (hedge-channel sign conditions)
%
% Status: SKELETON — [PLACEHOLDER] marks require server1 numerical results.
%         Structural arguments are complete; fill CEV magnitudes once
%         output/diagnostics/p6_option1_*.json land.
%
% Compilation: % s5_discussion.tex — Discussion
% Paper: "Tokenized Housing and Lifecycle Portfolio Choice:
%         A Decoupling of Location from Housing Exposure"
%
% Source: paper/outline_v4.md §5 (sections 5.1–5.3)
%         docs/welfare_decomp_v4.md §5 (literature comparison table)
%         paper/outline_v4.md Appendix C (hedge-channel sign conditions)
%
% Status: SKELETON — [PLACEHOLDER] marks require server1 numerical results.
%         Structural arguments are complete; fill CEV magnitudes once
%         output/diagnostics/p6_option1_*.json land.
%
% Compilation: \input{sections/s5_discussion} from main.tex
%              Requires amsmath, booktabs, siunitx, hyperref, natbib.

% ─────────────────────────────────────────────────────────────────────────────
\section{Discussion}
\label{sec:discussion}
% ─────────────────────────────────────────────────────────────────────────────

The headline welfare gain of
$\operatorname{CEV}(\text{E2\_2L vs E1\_2L}) = \textsc{[x.xx\%]}$
[\textsc{placeholder}]
arises from two distinct mechanisms, each with a different structural origin
and a different relationship to the existing lifecycle housing-portfolio
literature.
We discuss each mechanism in turn, then address the partial-equilibrium
scope of the welfare estimates and the paper's policy implications.

% ─────────────────────────────────────────────────────────────────────────────
\subsection{Relation to the Continuous-Ownership Literature}
\label{sec:discussion:liu}
% ─────────────────────────────────────────────────────────────────────────────

The continuous-$x$ component of the E2\_2L welfare gain~---~the ability to
hold a fractional share $x \in [0,\infty)$ of the occupied location's
residential value, paying maintenance proportional to share held and receiving
rent savings accordingly~---~shares its economic logic with \citet{Liu2021}.
\citeauthor{Liu2021} analyzes a Minimum Housing Share (MHS) relaxation that
similarly allows a household to hold a continuous fraction of housing rather
than the binary own-or-rent choice of \citet{YaoZhang2005} and \citet{Cocco2005}.
Within a single-location, no-relocation framework, \citeauthor{Liu2021} reports
welfare gains of 5--10\% of lifetime consumption from this indivisibility
relaxation.

Our model confirms the continuous-$x$ channel quantitatively: in the
v3~Option~3 baseline, $\operatorname{CEV}(\text{E2\_2L vs E1\_2L\_NOTX}) =
+3.41\%$ is attributable to the continuous-$x$ mechanism, consistent with
the lower end of \citeauthor{Liu2021}'s range
(the difference reflecting our more conservative grid resolution and the
presence of a relocation shock that occasionally forces households away from
their optimal stationary allocation).

Our contribution \emph{beyond} \citeauthor{Liu2021} rests on two mechanisms
that are absent from the single-location framework:
\begin{enumerate}[(i)]
\item \textbf{Round-trip transaction-cost avoidance.}
      Under traditional binary homeownership (E1\_2L), each relocation forces
      a sell-and-buy round-trip at approximately $\tau_{\text{sell}} +
      \tau_{\text{buy}} = 6\% + 2.5\% = 8.5\%$ of housing value.
      Tokenized holdings avoid the selling cost entirely (tokens are retained
      across moves) and amortize the buying cost over multiple periods via
      gradual pre-accumulation.
      In the v3~Option~3 baseline, this channel contributes
      $+0.82\%$ ($= +0.57\%$ from sell-side avoidance $+ +0.25\%$ from
      buy-side approximation).
      \citeauthor{Liu2021} has no relocation shock and no transaction-cost
      block; this channel is structurally absent from that framework.

\item \textbf{Pre-buying hedge via gradual accumulation.}
      Under the 6D state (v4~Option~1), a household at location~$A$ who
      anticipates possible relocation to~$B$ can pre-accumulate fractional
      $B$-location tokens incrementally.
      Each unit of $x_B$ pre-accumulated while at~$A$ avoids paying
      $\tau_{\text{buy}} = 2.5\%$ on that unit at the moment of relocation
      arrival.
      The per-period expected saving from holding one unit of $x_B$ at~$A$
      is $p_{\text{reloc}} \cdot \tau_{\text{buy}} = 0.06 \times 0.025 =
      0.15\%$; cumulated over the working life, this can deliver a
      lifetime CEV contribution of approximately 0.5--1.5\%.
      [\textsc{Placeholder: fill from server1 p6\_option1\_e2.json.}]
      This mechanism has no analog in \citeauthor{Liu2021} or in any prior
      lifecycle housing-portfolio model known to us.
\end{enumerate}

The structural claim is therefore:
\begin{equation}
    \underbrace{\operatorname{CEV}(\text{E2\_2L vs E1\_2L})}_{\text{Token
    value}} =
    \underbrace{\operatorname{CEV}_{\text{Liu}}}_{\approx 3.4\%}
    + \underbrace{\text{Tx-cost avoidance}}_{\approx 0.8\%}
    + \underbrace{\text{Pre-buying hedge}}_{\approx 0.5\text{--}1.5\%}
    + \text{cross-term}.
    \label{eq:decomp_claim}
\end{equation}
The cross-term is reported directly (Section~\ref{sec:results:decomp});
in the v3 baseline it equals $+0.02\%$, confirming near-additive separability
of channels.

% ─────────────────────────────────────────────────────────────────────────────
\subsection{The Pre-Buying Hedge: Conditions for Activation}
\label{sec:discussion:hedge}
% ─────────────────────────────────────────────────────────────────────────────

The pre-buying hedge mechanism is subtle.
A household at location~$A$ faces a per-unit opportunity cost of holding
$x_B > 0$: the unit of wealth allocated to $x_B$ forgoes the rent saving
that $x_A$ would provide, since only the \emph{occupied-location} token
reduces net housing cost:
\begin{equation}
    \kappa(x_A, x_B \mid \ell{=}A)
    = \rho - x_A \cdot (\rho - m).
    \label{eq:kappa_A}
\end{equation}
Equation~\eqref{eq:kappa_A} implies that each unit of $x_B$ held while
at~$A$ carries a per-period opportunity cost of $\delta_{\text{own}} = \rho - m
= 0.04$ (the rent-saving foregone by not putting that wealth into $x_A$).
The pre-buying benefit of holding $x_B$ is $p_{\text{reloc}} \cdot
\tau_{\text{buy}} = 0.0015$ per period.

Since $0.04 \gg 0.0015$, the naive static calculation suggests $x_B = 0$
is always optimal---consistent with the failed v3~Option~3 mechanism
test (Section~\ref{sec:results:falsification}).\footnote{Under the earlier
over-generous kappa rule $\kappa = \rho - (x_A + x_B)(\rho - m)$, the
$x_B$ opportunity cost was zero, mechanically generating $x_B > 0$
as a spurious artifact (\citealt{research_log_2026_05_01_fix}).
The corrected kappa rule~\eqref{eq:kappa_A} eliminates this artifact.}

Two forces can restore the pre-buying hedge in equilibrium.
First, at low wealth levels, the household cannot afford to bring $x_A$ to
its wealth-constrained ceiling; the marginal unit of wealth invested in $x_B$
rather than $x_A$ sacrifices rent-saving proportional to the \emph{uncovered
fraction} of $x_A$, which may be small when $x_A$ is already near the ceiling
for low-wealth households.
Second, the 6D state tracks $x_{A,\text{prev}}$ and $x_{B,\text{prev}}$
jointly; a household who has already accumulated a large $x_A$ position
(with $x_A \approx x_{A,\text{prev}}$, so $\Delta x_A \approx 0$) pays no
additional $\tau_{\text{buy}}$ to maintain it, making the marginal cost of
diverting a small increment to $x_B$ much lower than the static calculation
suggests.

The pre-registered hedge diagnostic (`mean_xB_at_ellA_xprev0_t1` in the solver
summary) tests whether this pre-buying demand is positive at the initial state
$x_{A,\text{prev}} = x_{B,\text{prev}} = 0$.
If $\bar{x}_B > 0$ (Hypothesis H1), the mechanism activates and channel~(ii)
in equation~\eqref{eq:decomp_claim} is quantitatively non-negligible.
If $\bar{x}_B = 0$ everywhere, the mechanism is dominated by the opportunity
cost and the paper's contribution rests on channels (i) and the Liu-type
continuous-$x$ channel only.
[\textsc{Placeholder: fill \textbf{H1 status} once server1 JSON lands.}]

% ─────────────────────────────────────────────────────────────────────────────
\subsection{Why REITs Cannot Replicate}
\label{sec:discussion:reit}
% ─────────────────────────────────────────────────────────────────────────────

A natural alternative for households seeking housing-market exposure after
relocation is a diversified real estate investment trust (REIT).
We dropped the REIT asset from the v3 model (it was present in the v2 framework)
for a structural reason that this section makes explicit.

REITs aggregate cash flows from a national or regional portfolio of properties.
An investment in a REIT therefore provides exposure to the \emph{common factor}
$R_{\text{div}}$ in housing returns~---~the component shared by all markets~---~but
not to the idiosyncratic location-specific component $\iota_\ell$.
Under our return decomposition
$\log R_\ell = \mu_h + \eta_{\text{div}} + \iota_\ell$
(Section~\ref{sec:model:returns}), a REIT-like diversified asset delivers
$E[\eta_{\text{div}}]$ but $E[\iota_\ell] = 0$ in large diversified portfolios.

Location-$A$ tokens, by contrast, deliver $\iota_A$~---~the idiosyncratic
appreciation specific to the MSA the household has personal ties to.
For a household who lived in location~$A$, built human capital there, and
may return, the idiosyncratic component $\iota_A$ is exactly the exposure
they wish to retain.
No REIT can supply this:
the location-specific fractional token is the only instrument that provides
$\iota_A$-loading for a household currently at~$B$.

This is why the 2026-05-01 pivot from the v2 four-regime (REIT-comparison)
framework to the v3/v4 two-location framework was necessary: the
structurally novel contribution of tokenization is not relative to REIT access
but relative to \emph{direct ownership with forced relocation sale}.

% ─────────────────────────────────────────────────────────────────────────────
\subsection{Relation to Other Lifecycle Housing Models}
\label{sec:discussion:literature}
% ─────────────────────────────────────────────────────────────────────────────

Table~\ref{tab:literature} summarizes the placement of our contribution
relative to the closest precedents.

\begin{table}[t]
\centering
\caption{Relationship to Prior Lifecycle Housing-Portfolio Models}
\label{tab:literature}
\small
\begin{tabular}{@{}lllll@{}}
\toprule
Paper & Relocation & Tx costs & Continuous $x$ & Token portability \\
\midrule
\citet{YaoZhang2005}     & No  & No  & No & No \\
\citet{Cocco2005}         & No  & No  & No & No \\
\citet{KraftMunk2011}     & No  & No  & No & No \\
\citet{KMW2018}           & Yes (reduced-form) & No  & No & No \\
\citet{Liu2021}           & No  & No  & Yes & No \\
\textbf{This paper (E1\_2L)} & Yes & Yes & No  & No \\
\textbf{This paper (E2\_2L)} & Yes & Yes (on $\Delta x$) & Yes & \textbf{Yes} \\
\bottomrule
\end{tabular}
\begin{tablenotes}
\footnotesize
\item \textit{Notes.}
``Relocation'' indicates whether the model includes a stochastic
geographic relocation shock.
``Tx costs'' indicates whether transaction costs (selling and/or
buying commissions) are modelled.
``Continuous $x$'' indicates whether fractional ownership is permitted.
``Token portability'' indicates whether the fractional position can be
retained across a geographic relocation without a forced sale.
\end{tablenotes}
\end{table}

\citet{KMW2018} model relocation explicitly in a lifecycle framework with
habit formation, but transaction costs are absent and the continuous-$x$
channel is not modelled.
Their focus is the housing-as-habit hedge against future consumption
commitments, a mechanism orthogonal to our pre-buying channel.

\citet{SinaiSouleles2005} study the rent-price risk hedge embedded in
homeownership: owning provides insurance against future rent increases for
households with long planned tenures.
Their mechanism is complementary to ours~---~the rent-hedge motive
encourages ownership even when the financial return to housing is below
the equity premium, exactly as in our model~---~but they abstract from
relocation and portability.

\citet{Davidoff2006} studies the income-housing correlation (``double
exposure'' risk) and its interaction with optimal tenure.
This correlation is relevant to our model through the $(\iota_A, \text{income})$
channel; we abstract from it at baseline by maintaining independent income
and housing shocks, and defer an income-correlated location shock extension
to future work (see Section~\ref{sec:discussion:limitations}).

% ─────────────────────────────────────────────────────────────────────────────
\subsection{Robustness of the Mechanism}
\label{sec:discussion:robustness}
% ─────────────────────────────────────────────────────────────────────────────

The pre-registered falsification structure (Section~\ref{sec:results:falsification})
provides the clearest evidence on the robustness of each channel.

\paragraph{Test (r): Relocation disabled ($p_{\text{reloc}} = 0$).}
With no relocation shocks, the motivation to hold $x_B$ while at~$A$
vanishes~---~the household will never visit location~$B$, so no
pre-buying saving accrues.
Under v4~Option~1, the pre-buying hedge channel must therefore collapse
to zero, while the continuous-$x$ channel persists (it does not require
relocation).
Concretely, the expected result is $\bar{x}_B \to 0$ and
$\operatorname{CEV}(E2\_2L \text{ vs } E1\_2L) \to
\operatorname{CEV}_{\text{Liu}}$ (continuous-$x$ only).
[\textsc{Placeholder: report $\Delta\text{CEV}$ and $\bar{x}_B$ from sweep.}]

\paragraph{Test (m): High location correlation ($\rho_{AB} = 0.95$).}
As $\rho_{AB} \to 1$, the location-$B$ token delivers almost the same
return as the location-$A$ token; the diversification value of the
cross-location hedge diminishes.
The pre-buying saving is unchanged in magnitude ($p_{\text{reloc}} \cdot
\tau_{\text{buy}}$ is independent of $\rho_{AB}$), but households who
value the idiosyncratic hedge component will demand less $x_B$ as
$\rho_{AB}$ rises.
We therefore expect the hedge channel to attenuate~---~but not to
collapse completely~---~as $\rho_{AB} \to 0.95$.
[\textsc{Placeholder: report $\bar{x}_B(\rho_{AB} = 0.95)$ from sweep.}]

\paragraph{Test (q): No buying cost ($\tau_{\text{buy}} = 0$).}
With $\tau_{\text{buy}} = 0$, pre-buying $x_B$ saves nothing; the
pre-buying motive is eliminated.
The expected outcome is $\bar{x}_B = 0$ (identical to v3~Option~3 under
zero buying cost) and a CEV change equal only to the sell-side avoidance
and continuous-$x$ channels.
This isolates the pre-buying contribution to the total welfare gain.
[\textsc{Placeholder: report from sweep\_txcost.sh run.}]

\paragraph{Sensitivity to $p_{\text{reloc}}$ and $\rho_{AB}$ jointly.}
Figure~\ref{fig:heatmap_cev} (placeholder) reports the sensitivity heatmap
over $(\rho_{AB}, p_{\text{reloc}}) \in \{0, 0.25, 0.50, 0.75, 0.95\} \times
\{0, 0.02, 0.06, 0.12\}$.
The model predicts a weakly increasing CEV in $p_{\text{reloc}}$ (more
mobility amplifies pre-buying value) and a weakly decreasing CEV in
$\rho_{AB}$ (higher correlation reduces the idiosyncratic hedge value).
If the empirical pattern deviates from these monotone predictions,
the mechanism claim requires reassessment.
[\textsc{Placeholder: fill from scripts/sweep\_rhoAB.sh and
scripts/sweep\_prelocate.sh.}]

% ─────────────────────────────────────────────────────────────────────────────
\subsection{Partial-Equilibrium Scope and Limitations}
\label{sec:discussion:limitations}
% ─────────────────────────────────────────────────────────────────────────────

All welfare estimates in this paper are partial-equilibrium in three respects.

\paragraph{Exogenous house prices.}
We take the location-specific house-price processes $(R_A, R_B)$ as
exogenous, calibrated to observed MSA-level moments.
Widespread adoption of token portability could affect housing demand in
origin and destination locations, potentially altering $R_A$ and $R_B$.
A general-equilibrium model with endogenous prices is left to future work.
Our CEV estimates should be interpreted as the household-level welfare value
of \emph{access} to tokenized instruments, conditional on other households'
behavior remaining unchanged.

\paragraph{Symmetric and stationary location shocks.}
The relocation shock is i.i.d.\ Bernoulli with constant probability
$p_{\text{reloc}}(t)$.
In practice, relocation rates are correlated with employment shocks,
income shocks, and housing-market conditions (\citealt{PSID}).
An extension with income-correlated location shocks (e.g.,
$\text{corr}(\varepsilon_t, \iota_{A,t}) > 0$) would add a second
channel to the pre-buying demand~---~holding $x_B$ at~$A$ also hedges
against a negative income draw at~$A$ that triggers relocation to~$B$.
This extension is queued as a robustness check but is not included in
the baseline results.

\paragraph{Two-location, symmetric-returns simplification.}
We model two symmetric locations with exchangeable house-price processes
(same $\mu_h$, $\sigma_h$; differing only by idiosyncratic component).
Empirically, households relocate across heterogeneous MSAs with different
mean appreciation rates and rent-to-price ratios.
Asymmetric robustness (different $\mu_A \neq \mu_B$, different
$p_{A \to B} \neq p_{B \to A}$) is queued in Phase~2 sensitivity work
(Table~\ref{tab:sensitivity}, rows 6--7).

\paragraph{No mortgage.}
The baseline results set $\text{LTV\_MAX} = 0$.
Mortgages reduce the effective down payment on housing, raising the
household's ability to reach the optimal $x_A$ at low wealth levels.
This can attenuate the pre-buying demand for $x_B$ (see
Section~\ref{sec:discussion:hedge}) if households who were
wealth-constrained in $x_A$ can now fully fund their desired $x_A$
via borrowing.
Mortgage activation (LTV $\in \{0.5, 0.8\}$) is queued as a Phase~2
sensitivity check.
In the v2 baseline (\citealt{research_log_2026_05_01_mortgage}),
mortgage reduced $\operatorname{CEV}(\text{E2\_2L vs E1\_2L})$ by
approximately 37\% under no-relocation; the direction under the
2-location model may differ.

\paragraph{Platform and regulatory risk.}
We abstract from platform failure, regulatory uncertainty, and
information asymmetries between token issuers and holders.
These factors introduce a risk premium on tokenized positions that
is not captured in the $\tau_{\text{token}} = 1\%$ transfer cost
alone.
A fuller treatment of platform risk is deferred to companion work.

% ─────────────────────────────────────────────────────────────────────────────
\subsection{Policy Implications}
\label{sec:discussion:policy}
% ─────────────────────────────────────────────────────────────────────────────

Our findings have three direct policy implications for household finance
and housing-market regulation.

\paragraph{Mobility lock-in as a welfare cost.}
The $+0.82\%$ transaction-cost-avoidance channel in the v3~Option~3 baseline
quantifies the welfare cost of \emph{mobility lock-in}: households who own
traditional homes are effectively penalized for relocating, reducing their
geographic flexibility.
Token portability eliminates this distortion.
Policy interventions that reduce $\tau_{\text{sell}}$ or $\tau_{\text{buy}}$
(e.g., real-estate transfer tax reform) would deliver similar gains without
requiring tokenization, but would not restore the idiosyncratic hedge or the
pre-buying channel.

\paragraph{Token design and the buying-cost schedule.}
If $\tau_{\text{buy}}$ is too high, the pre-buying hedge becomes economically
negligible ($p_{\text{reloc}} \cdot \tau_{\text{buy}}$ too small relative to
the $\delta_{\text{own}}$ opportunity cost).
If $\tau_{\text{buy}}$ is too low (approaching $\tau_{\text{token}} = 1\%$),
the pre-buying saving is small and the mechanism collapses.
The sensitivity analysis over $\tau_{\text{buy}} \in \{0, 1\%, 2.5\%, 4\%, 6\%\}$
identifies the buying-cost levels at which the pre-buying channel is most
economically significant.
This informs platform pricing: the welfare case for token portability relative
to direct ownership is maximized at intermediate buying cost levels.

\paragraph{Household heterogeneity.}
The pre-buying mechanism is most valuable for households who are
(i)~frequently mobile ($p_{\text{reloc}}$ high),
(ii)~have location-specific ties to prior MSAs ($\iota_A$ relevant to their
     welfare), and
(iii)~are wealth-constrained in their $x_A$ position (forcing $\Delta x_A = 0$
      and making $x_B$ investment relatively cheap).
From a regulatory standpoint, this suggests that token portability delivers
disproportionate welfare gains to mobile younger households with limited
wealth~---~a population underserved by existing housing finance instruments.
 from main.tex
%              Requires amsmath, booktabs, siunitx, hyperref, natbib.

% ─────────────────────────────────────────────────────────────────────────────
\section{Discussion}
\label{sec:discussion}
% ─────────────────────────────────────────────────────────────────────────────

The headline welfare gain of
$\operatorname{CEV}(\text{E2\_2L vs E1\_2L}) = \textsc{[x.xx\%]}$
[\textsc{placeholder}]
arises from two distinct mechanisms, each with a different structural origin
and a different relationship to the existing lifecycle housing-portfolio
literature.
We discuss each mechanism in turn, then address the partial-equilibrium
scope of the welfare estimates and the paper's policy implications.

% ─────────────────────────────────────────────────────────────────────────────
\subsection{Relation to the Continuous-Ownership Literature}
\label{sec:discussion:liu}
% ─────────────────────────────────────────────────────────────────────────────

The continuous-$x$ component of the E2\_2L welfare gain~---~the ability to
hold a fractional share $x \in [0,\infty)$ of the occupied location's
residential value, paying maintenance proportional to share held and receiving
rent savings accordingly~---~shares its economic logic with \citet{Liu2021}.
\citeauthor{Liu2021} analyzes a Minimum Housing Share (MHS) relaxation that
similarly allows a household to hold a continuous fraction of housing rather
than the binary own-or-rent choice of \citet{YaoZhang2005} and \citet{Cocco2005}.
Within a single-location, no-relocation framework, \citeauthor{Liu2021} reports
welfare gains of 5--10\% of lifetime consumption from this indivisibility
relaxation.

Our model confirms the continuous-$x$ channel quantitatively: in the
v3~Option~3 baseline, $\operatorname{CEV}(\text{E2\_2L vs E1\_2L\_NOTX}) =
+3.41\%$ is attributable to the continuous-$x$ mechanism, consistent with
the lower end of \citeauthor{Liu2021}'s range
(the difference reflecting our more conservative grid resolution and the
presence of a relocation shock that occasionally forces households away from
their optimal stationary allocation).

Our contribution \emph{beyond} \citeauthor{Liu2021} rests on two mechanisms
that are absent from the single-location framework:
\begin{enumerate}[(i)]
\item \textbf{Round-trip transaction-cost avoidance.}
      Under traditional binary homeownership (E1\_2L), each relocation forces
      a sell-and-buy round-trip at approximately $\tau_{\text{sell}} +
      \tau_{\text{buy}} = 6\% + 2.5\% = 8.5\%$ of housing value.
      Tokenized holdings avoid the selling cost entirely (tokens are retained
      across moves) and amortize the buying cost over multiple periods via
      gradual pre-accumulation.
      In the v3~Option~3 baseline, this channel contributes
      $+0.82\%$ ($= +0.57\%$ from sell-side avoidance $+ +0.25\%$ from
      buy-side approximation).
      \citeauthor{Liu2021} has no relocation shock and no transaction-cost
      block; this channel is structurally absent from that framework.

\item \textbf{Pre-buying hedge via gradual accumulation.}
      Under the 6D state (v4~Option~1), a household at location~$A$ who
      anticipates possible relocation to~$B$ can pre-accumulate fractional
      $B$-location tokens incrementally.
      Each unit of $x_B$ pre-accumulated while at~$A$ avoids paying
      $\tau_{\text{buy}} = 2.5\%$ on that unit at the moment of relocation
      arrival.
      The per-period expected saving from holding one unit of $x_B$ at~$A$
      is $p_{\text{reloc}} \cdot \tau_{\text{buy}} = 0.06 \times 0.025 =
      0.15\%$; cumulated over the working life, this can deliver a
      lifetime CEV contribution of approximately 0.5--1.5\%.
      [\textsc{Placeholder: fill from server1 p6\_option1\_e2.json.}]
      This mechanism has no analog in \citeauthor{Liu2021} or in any prior
      lifecycle housing-portfolio model known to us.
\end{enumerate}

The structural claim is therefore:
\begin{equation}
    \underbrace{\operatorname{CEV}(\text{E2\_2L vs E1\_2L})}_{\text{Token
    value}} =
    \underbrace{\operatorname{CEV}_{\text{Liu}}}_{\approx 3.4\%}
    + \underbrace{\text{Tx-cost avoidance}}_{\approx 0.8\%}
    + \underbrace{\text{Pre-buying hedge}}_{\approx 0.5\text{--}1.5\%}
    + \text{cross-term}.
    \label{eq:decomp_claim}
\end{equation}
The cross-term is reported directly (Section~\ref{sec:results:decomp});
in the v3 baseline it equals $+0.02\%$, confirming near-additive separability
of channels.

% ─────────────────────────────────────────────────────────────────────────────
\subsection{The Pre-Buying Hedge: Conditions for Activation}
\label{sec:discussion:hedge}
% ─────────────────────────────────────────────────────────────────────────────

The pre-buying hedge mechanism is subtle.
A household at location~$A$ faces a per-unit opportunity cost of holding
$x_B > 0$: the unit of wealth allocated to $x_B$ forgoes the rent saving
that $x_A$ would provide, since only the \emph{occupied-location} token
reduces net housing cost:
\begin{equation}
    \kappa(x_A, x_B \mid \ell{=}A)
    = \rho - x_A \cdot (\rho - m).
    \label{eq:kappa_A}
\end{equation}
Equation~\eqref{eq:kappa_A} implies that each unit of $x_B$ held while
at~$A$ carries a per-period opportunity cost of $\delta_{\text{own}} = \rho - m
= 0.04$ (the rent-saving foregone by not putting that wealth into $x_A$).
The pre-buying benefit of holding $x_B$ is $p_{\text{reloc}} \cdot
\tau_{\text{buy}} = 0.0015$ per period.

Since $0.04 \gg 0.0015$, the naive static calculation suggests $x_B = 0$
is always optimal---consistent with the failed v3~Option~3 mechanism
test (Section~\ref{sec:results:falsification}).\footnote{Under the earlier
over-generous kappa rule $\kappa = \rho - (x_A + x_B)(\rho - m)$, the
$x_B$ opportunity cost was zero, mechanically generating $x_B > 0$
as a spurious artifact (\citealt{research_log_2026_05_01_fix}).
The corrected kappa rule~\eqref{eq:kappa_A} eliminates this artifact.}

Two forces can restore the pre-buying hedge in equilibrium.
First, at low wealth levels, the household cannot afford to bring $x_A$ to
its wealth-constrained ceiling; the marginal unit of wealth invested in $x_B$
rather than $x_A$ sacrifices rent-saving proportional to the \emph{uncovered
fraction} of $x_A$, which may be small when $x_A$ is already near the ceiling
for low-wealth households.
Second, the 6D state tracks $x_{A,\text{prev}}$ and $x_{B,\text{prev}}$
jointly; a household who has already accumulated a large $x_A$ position
(with $x_A \approx x_{A,\text{prev}}$, so $\Delta x_A \approx 0$) pays no
additional $\tau_{\text{buy}}$ to maintain it, making the marginal cost of
diverting a small increment to $x_B$ much lower than the static calculation
suggests.

The pre-registered hedge diagnostic (`mean_xB_at_ellA_xprev0_t1` in the solver
summary) tests whether this pre-buying demand is positive at the initial state
$x_{A,\text{prev}} = x_{B,\text{prev}} = 0$.
If $\bar{x}_B > 0$ (Hypothesis H1), the mechanism activates and channel~(ii)
in equation~\eqref{eq:decomp_claim} is quantitatively non-negligible.
If $\bar{x}_B = 0$ everywhere, the mechanism is dominated by the opportunity
cost and the paper's contribution rests on channels (i) and the Liu-type
continuous-$x$ channel only.
[\textsc{Placeholder: fill \textbf{H1 status} once server1 JSON lands.}]

% ─────────────────────────────────────────────────────────────────────────────
\subsection{Why REITs Cannot Replicate}
\label{sec:discussion:reit}
% ─────────────────────────────────────────────────────────────────────────────

A natural alternative for households seeking housing-market exposure after
relocation is a diversified real estate investment trust (REIT).
We dropped the REIT asset from the v3 model (it was present in the v2 framework)
for a structural reason that this section makes explicit.

REITs aggregate cash flows from a national or regional portfolio of properties.
An investment in a REIT therefore provides exposure to the \emph{common factor}
$R_{\text{div}}$ in housing returns~---~the component shared by all markets~---~but
not to the idiosyncratic location-specific component $\iota_\ell$.
Under our return decomposition
$\log R_\ell = \mu_h + \eta_{\text{div}} + \iota_\ell$
(Section~\ref{sec:model:returns}), a REIT-like diversified asset delivers
$E[\eta_{\text{div}}]$ but $E[\iota_\ell] = 0$ in large diversified portfolios.

Location-$A$ tokens, by contrast, deliver $\iota_A$~---~the idiosyncratic
appreciation specific to the MSA the household has personal ties to.
For a household who lived in location~$A$, built human capital there, and
may return, the idiosyncratic component $\iota_A$ is exactly the exposure
they wish to retain.
No REIT can supply this:
the location-specific fractional token is the only instrument that provides
$\iota_A$-loading for a household currently at~$B$.

This is why the 2026-05-01 pivot from the v2 four-regime (REIT-comparison)
framework to the v3/v4 two-location framework was necessary: the
structurally novel contribution of tokenization is not relative to REIT access
but relative to \emph{direct ownership with forced relocation sale}.

% ─────────────────────────────────────────────────────────────────────────────
\subsection{Relation to Other Lifecycle Housing Models}
\label{sec:discussion:literature}
% ─────────────────────────────────────────────────────────────────────────────

Table~\ref{tab:literature} summarizes the placement of our contribution
relative to the closest precedents.

\begin{table}[t]
\centering
\caption{Relationship to Prior Lifecycle Housing-Portfolio Models}
\label{tab:literature}
\small
\begin{tabular}{@{}lllll@{}}
\toprule
Paper & Relocation & Tx costs & Continuous $x$ & Token portability \\
\midrule
\citet{YaoZhang2005}     & No  & No  & No & No \\
\citet{Cocco2005}         & No  & No  & No & No \\
\citet{KraftMunk2011}     & No  & No  & No & No \\
\citet{KMW2018}           & Yes (reduced-form) & No  & No & No \\
\citet{Liu2021}           & No  & No  & Yes & No \\
\textbf{This paper (E1\_2L)} & Yes & Yes & No  & No \\
\textbf{This paper (E2\_2L)} & Yes & Yes (on $\Delta x$) & Yes & \textbf{Yes} \\
\bottomrule
\end{tabular}
\begin{tablenotes}
\footnotesize
\item \textit{Notes.}
``Relocation'' indicates whether the model includes a stochastic
geographic relocation shock.
``Tx costs'' indicates whether transaction costs (selling and/or
buying commissions) are modelled.
``Continuous $x$'' indicates whether fractional ownership is permitted.
``Token portability'' indicates whether the fractional position can be
retained across a geographic relocation without a forced sale.
\end{tablenotes}
\end{table}

\citet{KMW2018} model relocation explicitly in a lifecycle framework with
habit formation, but transaction costs are absent and the continuous-$x$
channel is not modelled.
Their focus is the housing-as-habit hedge against future consumption
commitments, a mechanism orthogonal to our pre-buying channel.

\citet{SinaiSouleles2005} study the rent-price risk hedge embedded in
homeownership: owning provides insurance against future rent increases for
households with long planned tenures.
Their mechanism is complementary to ours~---~the rent-hedge motive
encourages ownership even when the financial return to housing is below
the equity premium, exactly as in our model~---~but they abstract from
relocation and portability.

\citet{Davidoff2006} studies the income-housing correlation (``double
exposure'' risk) and its interaction with optimal tenure.
This correlation is relevant to our model through the $(\iota_A, \text{income})$
channel; we abstract from it at baseline by maintaining independent income
and housing shocks, and defer an income-correlated location shock extension
to future work (see Section~\ref{sec:discussion:limitations}).

% ─────────────────────────────────────────────────────────────────────────────
\subsection{Robustness of the Mechanism}
\label{sec:discussion:robustness}
% ─────────────────────────────────────────────────────────────────────────────

The pre-registered falsification structure (Section~\ref{sec:results:falsification})
provides the clearest evidence on the robustness of each channel.

\paragraph{Test (r): Relocation disabled ($p_{\text{reloc}} = 0$).}
With no relocation shocks, the motivation to hold $x_B$ while at~$A$
vanishes~---~the household will never visit location~$B$, so no
pre-buying saving accrues.
Under v4~Option~1, the pre-buying hedge channel must therefore collapse
to zero, while the continuous-$x$ channel persists (it does not require
relocation).
Concretely, the expected result is $\bar{x}_B \to 0$ and
$\operatorname{CEV}(E2\_2L \text{ vs } E1\_2L) \to
\operatorname{CEV}_{\text{Liu}}$ (continuous-$x$ only).
[\textsc{Placeholder: report $\Delta\text{CEV}$ and $\bar{x}_B$ from sweep.}]

\paragraph{Test (m): High location correlation ($\rho_{AB} = 0.95$).}
As $\rho_{AB} \to 1$, the location-$B$ token delivers almost the same
return as the location-$A$ token; the diversification value of the
cross-location hedge diminishes.
The pre-buying saving is unchanged in magnitude ($p_{\text{reloc}} \cdot
\tau_{\text{buy}}$ is independent of $\rho_{AB}$), but households who
value the idiosyncratic hedge component will demand less $x_B$ as
$\rho_{AB}$ rises.
We therefore expect the hedge channel to attenuate~---~but not to
collapse completely~---~as $\rho_{AB} \to 0.95$.
[\textsc{Placeholder: report $\bar{x}_B(\rho_{AB} = 0.95)$ from sweep.}]

\paragraph{Test (q): No buying cost ($\tau_{\text{buy}} = 0$).}
With $\tau_{\text{buy}} = 0$, pre-buying $x_B$ saves nothing; the
pre-buying motive is eliminated.
The expected outcome is $\bar{x}_B = 0$ (identical to v3~Option~3 under
zero buying cost) and a CEV change equal only to the sell-side avoidance
and continuous-$x$ channels.
This isolates the pre-buying contribution to the total welfare gain.
[\textsc{Placeholder: report from sweep\_txcost.sh run.}]

\paragraph{Sensitivity to $p_{\text{reloc}}$ and $\rho_{AB}$ jointly.}
Figure~\ref{fig:heatmap_cev} (placeholder) reports the sensitivity heatmap
over $(\rho_{AB}, p_{\text{reloc}}) \in \{0, 0.25, 0.50, 0.75, 0.95\} \times
\{0, 0.02, 0.06, 0.12\}$.
The model predicts a weakly increasing CEV in $p_{\text{reloc}}$ (more
mobility amplifies pre-buying value) and a weakly decreasing CEV in
$\rho_{AB}$ (higher correlation reduces the idiosyncratic hedge value).
If the empirical pattern deviates from these monotone predictions,
the mechanism claim requires reassessment.
[\textsc{Placeholder: fill from scripts/sweep\_rhoAB.sh and
scripts/sweep\_prelocate.sh.}]

% ─────────────────────────────────────────────────────────────────────────────
\subsection{Partial-Equilibrium Scope and Limitations}
\label{sec:discussion:limitations}
% ─────────────────────────────────────────────────────────────────────────────

All welfare estimates in this paper are partial-equilibrium in three respects.

\paragraph{Exogenous house prices.}
We take the location-specific house-price processes $(R_A, R_B)$ as
exogenous, calibrated to observed MSA-level moments.
Widespread adoption of token portability could affect housing demand in
origin and destination locations, potentially altering $R_A$ and $R_B$.
A general-equilibrium model with endogenous prices is left to future work.
Our CEV estimates should be interpreted as the household-level welfare value
of \emph{access} to tokenized instruments, conditional on other households'
behavior remaining unchanged.

\paragraph{Symmetric and stationary location shocks.}
The relocation shock is i.i.d.\ Bernoulli with constant probability
$p_{\text{reloc}}(t)$.
In practice, relocation rates are correlated with employment shocks,
income shocks, and housing-market conditions (\citealt{PSID}).
An extension with income-correlated location shocks (e.g.,
$\text{corr}(\varepsilon_t, \iota_{A,t}) > 0$) would add a second
channel to the pre-buying demand~---~holding $x_B$ at~$A$ also hedges
against a negative income draw at~$A$ that triggers relocation to~$B$.
This extension is queued as a robustness check but is not included in
the baseline results.

\paragraph{Two-location, symmetric-returns simplification.}
We model two symmetric locations with exchangeable house-price processes
(same $\mu_h$, $\sigma_h$; differing only by idiosyncratic component).
Empirically, households relocate across heterogeneous MSAs with different
mean appreciation rates and rent-to-price ratios.
Asymmetric robustness (different $\mu_A \neq \mu_B$, different
$p_{A \to B} \neq p_{B \to A}$) is queued in Phase~2 sensitivity work
(Table~\ref{tab:sensitivity}, rows 6--7).

\paragraph{No mortgage.}
The baseline results set $\text{LTV\_MAX} = 0$.
Mortgages reduce the effective down payment on housing, raising the
household's ability to reach the optimal $x_A$ at low wealth levels.
This can attenuate the pre-buying demand for $x_B$ (see
Section~\ref{sec:discussion:hedge}) if households who were
wealth-constrained in $x_A$ can now fully fund their desired $x_A$
via borrowing.
Mortgage activation (LTV $\in \{0.5, 0.8\}$) is queued as a Phase~2
sensitivity check.
In the v2 baseline (\citealt{research_log_2026_05_01_mortgage}),
mortgage reduced $\operatorname{CEV}(\text{E2\_2L vs E1\_2L})$ by
approximately 37\% under no-relocation; the direction under the
2-location model may differ.

\paragraph{Platform and regulatory risk.}
We abstract from platform failure, regulatory uncertainty, and
information asymmetries between token issuers and holders.
These factors introduce a risk premium on tokenized positions that
is not captured in the $\tau_{\text{token}} = 1\%$ transfer cost
alone.
A fuller treatment of platform risk is deferred to companion work.

% ─────────────────────────────────────────────────────────────────────────────
\subsection{Policy Implications}
\label{sec:discussion:policy}
% ─────────────────────────────────────────────────────────────────────────────

Our findings have three direct policy implications for household finance
and housing-market regulation.

\paragraph{Mobility lock-in as a welfare cost.}
The $+0.82\%$ transaction-cost-avoidance channel in the v3~Option~3 baseline
quantifies the welfare cost of \emph{mobility lock-in}: households who own
traditional homes are effectively penalized for relocating, reducing their
geographic flexibility.
Token portability eliminates this distortion.
Policy interventions that reduce $\tau_{\text{sell}}$ or $\tau_{\text{buy}}$
(e.g., real-estate transfer tax reform) would deliver similar gains without
requiring tokenization, but would not restore the idiosyncratic hedge or the
pre-buying channel.

\paragraph{Token design and the buying-cost schedule.}
If $\tau_{\text{buy}}$ is too high, the pre-buying hedge becomes economically
negligible ($p_{\text{reloc}} \cdot \tau_{\text{buy}}$ too small relative to
the $\delta_{\text{own}}$ opportunity cost).
If $\tau_{\text{buy}}$ is too low (approaching $\tau_{\text{token}} = 1\%$),
the pre-buying saving is small and the mechanism collapses.
The sensitivity analysis over $\tau_{\text{buy}} \in \{0, 1\%, 2.5\%, 4\%, 6\%\}$
identifies the buying-cost levels at which the pre-buying channel is most
economically significant.
This informs platform pricing: the welfare case for token portability relative
to direct ownership is maximized at intermediate buying cost levels.

\paragraph{Household heterogeneity.}
The pre-buying mechanism is most valuable for households who are
(i)~frequently mobile ($p_{\text{reloc}}$ high),
(ii)~have location-specific ties to prior MSAs ($\iota_A$ relevant to their
     welfare), and
(iii)~are wealth-constrained in their $x_A$ position (forcing $\Delta x_A = 0$
      and making $x_B$ investment relatively cheap).
From a regulatory standpoint, this suggests that token portability delivers
disproportionate welfare gains to mobile younger households with limited
wealth~---~a population underserved by existing housing finance instruments.
 from main.tex
%              Requires amsmath, booktabs, siunitx, hyperref, natbib.

% ─────────────────────────────────────────────────────────────────────────────
\section{Discussion}
\label{sec:discussion}
% ─────────────────────────────────────────────────────────────────────────────

The headline welfare gain of
$\operatorname{CEV}(\text{E2\_2L vs E1\_2L}) = \textsc{[x.xx\%]}$
[\textsc{placeholder}]
arises from two distinct mechanisms, each with a different structural origin
and a different relationship to the existing lifecycle housing-portfolio
literature.
We discuss each mechanism in turn, then address the partial-equilibrium
scope of the welfare estimates and the paper's policy implications.

% ─────────────────────────────────────────────────────────────────────────────
\subsection{Relation to the Continuous-Ownership Literature}
\label{sec:discussion:liu}
% ─────────────────────────────────────────────────────────────────────────────

The continuous-$x$ component of the E2\_2L welfare gain~---~the ability to
hold a fractional share $x \in [0,\infty)$ of the occupied location's
residential value, paying maintenance proportional to share held and receiving
rent savings accordingly~---~shares its economic logic with \citet{Liu2021}.
\citeauthor{Liu2021} analyzes a Minimum Housing Share (MHS) relaxation that
similarly allows a household to hold a continuous fraction of housing rather
than the binary own-or-rent choice of \citet{YaoZhang2005} and \citet{Cocco2005}.
Within a single-location, no-relocation framework, \citeauthor{Liu2021} reports
welfare gains of 5--10\% of lifetime consumption from this indivisibility
relaxation.

Our model confirms the continuous-$x$ channel quantitatively: in the
v3~Option~3 baseline, $\operatorname{CEV}(\text{E2\_2L vs E1\_2L\_NOTX}) =
+3.41\%$ is attributable to the continuous-$x$ mechanism, consistent with
the lower end of \citeauthor{Liu2021}'s range
(the difference reflecting our more conservative grid resolution and the
presence of a relocation shock that occasionally forces households away from
their optimal stationary allocation).

Our contribution \emph{beyond} \citeauthor{Liu2021} rests on two mechanisms
that are absent from the single-location framework:
\begin{enumerate}[(i)]
\item \textbf{Round-trip transaction-cost avoidance.}
      Under traditional binary homeownership (E1\_2L), each relocation forces
      a sell-and-buy round-trip at approximately $\tau_{\text{sell}} +
      \tau_{\text{buy}} = 6\% + 2.5\% = 8.5\%$ of housing value.
      Tokenized holdings avoid the selling cost entirely (tokens are retained
      across moves) and amortize the buying cost over multiple periods via
      gradual pre-accumulation.
      In the v3~Option~3 baseline, this channel contributes
      $+0.82\%$ ($= +0.57\%$ from sell-side avoidance $+ +0.25\%$ from
      buy-side approximation).
      \citeauthor{Liu2021} has no relocation shock and no transaction-cost
      block; this channel is structurally absent from that framework.

\item \textbf{Pre-buying hedge via gradual accumulation.}
      Under the 6D state (v4~Option~1), a household at location~$A$ who
      anticipates possible relocation to~$B$ can pre-accumulate fractional
      $B$-location tokens incrementally.
      Each unit of $x_B$ pre-accumulated while at~$A$ avoids paying
      $\tau_{\text{buy}} = 2.5\%$ on that unit at the moment of relocation
      arrival.
      The per-period expected saving from holding one unit of $x_B$ at~$A$
      is $p_{\text{reloc}} \cdot \tau_{\text{buy}} = 0.06 \times 0.025 =
      0.15\%$; cumulated over the working life, this can deliver a
      lifetime CEV contribution of approximately 0.5--1.5\%.
      [\textsc{Placeholder: fill from server1 p6\_option1\_e2.json.}]
      This mechanism has no analog in \citeauthor{Liu2021} or in any prior
      lifecycle housing-portfolio model known to us.
\end{enumerate}

The structural claim is therefore:
\begin{equation}
    \underbrace{\operatorname{CEV}(\text{E2\_2L vs E1\_2L})}_{\text{Token
    value}} =
    \underbrace{\operatorname{CEV}_{\text{Liu}}}_{\approx 3.4\%}
    + \underbrace{\text{Tx-cost avoidance}}_{\approx 0.8\%}
    + \underbrace{\text{Pre-buying hedge}}_{\approx 0.5\text{--}1.5\%}
    + \text{cross-term}.
    \label{eq:decomp_claim}
\end{equation}
The cross-term is reported directly (Section~\ref{sec:results:decomp});
in the v3 baseline it equals $+0.02\%$, confirming near-additive separability
of channels.

% ─────────────────────────────────────────────────────────────────────────────
\subsection{The Pre-Buying Hedge: Conditions for Activation}
\label{sec:discussion:hedge}
% ─────────────────────────────────────────────────────────────────────────────

The pre-buying hedge mechanism is subtle.
A household at location~$A$ faces a per-unit opportunity cost of holding
$x_B > 0$: the unit of wealth allocated to $x_B$ forgoes the rent saving
that $x_A$ would provide, since only the \emph{occupied-location} token
reduces net housing cost:
\begin{equation}
    \kappa(x_A, x_B \mid \ell{=}A)
    = \rho - x_A \cdot (\rho - m).
    \label{eq:kappa_A}
\end{equation}
Equation~\eqref{eq:kappa_A} implies that each unit of $x_B$ held while
at~$A$ carries a per-period opportunity cost of $\delta_{\text{own}} = \rho - m
= 0.04$ (the rent-saving foregone by not putting that wealth into $x_A$).
The pre-buying benefit of holding $x_B$ is $p_{\text{reloc}} \cdot
\tau_{\text{buy}} = 0.0015$ per period.

Since $0.04 \gg 0.0015$, the naive static calculation suggests $x_B = 0$
is always optimal---consistent with the failed v3~Option~3 mechanism
test (Section~\ref{sec:results:falsification}).\footnote{Under the earlier
over-generous kappa rule $\kappa = \rho - (x_A + x_B)(\rho - m)$, the
$x_B$ opportunity cost was zero, mechanically generating $x_B > 0$
as a spurious artifact (\citealt{research_log_2026_05_01_fix}).
The corrected kappa rule~\eqref{eq:kappa_A} eliminates this artifact.}

Two forces can restore the pre-buying hedge in equilibrium.
First, at low wealth levels, the household cannot afford to bring $x_A$ to
its wealth-constrained ceiling; the marginal unit of wealth invested in $x_B$
rather than $x_A$ sacrifices rent-saving proportional to the \emph{uncovered
fraction} of $x_A$, which may be small when $x_A$ is already near the ceiling
for low-wealth households.
Second, the 6D state tracks $x_{A,\text{prev}}$ and $x_{B,\text{prev}}$
jointly; a household who has already accumulated a large $x_A$ position
(with $x_A \approx x_{A,\text{prev}}$, so $\Delta x_A \approx 0$) pays no
additional $\tau_{\text{buy}}$ to maintain it, making the marginal cost of
diverting a small increment to $x_B$ much lower than the static calculation
suggests.

The pre-registered hedge diagnostic (`mean_xB_at_ellA_xprev0_t1` in the solver
summary) tests whether this pre-buying demand is positive at the initial state
$x_{A,\text{prev}} = x_{B,\text{prev}} = 0$.
If $\bar{x}_B > 0$ (Hypothesis H1), the mechanism activates and channel~(ii)
in equation~\eqref{eq:decomp_claim} is quantitatively non-negligible.
If $\bar{x}_B = 0$ everywhere, the mechanism is dominated by the opportunity
cost and the paper's contribution rests on channels (i) and the Liu-type
continuous-$x$ channel only.
[\textsc{Placeholder: fill \textbf{H1 status} once server1 JSON lands.}]

% ─────────────────────────────────────────────────────────────────────────────
\subsection{Why REITs Cannot Replicate}
\label{sec:discussion:reit}
% ─────────────────────────────────────────────────────────────────────────────

A natural alternative for households seeking housing-market exposure after
relocation is a diversified real estate investment trust (REIT).
We dropped the REIT asset from the v3 model (it was present in the v2 framework)
for a structural reason that this section makes explicit.

REITs aggregate cash flows from a national or regional portfolio of properties.
An investment in a REIT therefore provides exposure to the \emph{common factor}
$R_{\text{div}}$ in housing returns~---~the component shared by all markets~---~but
not to the idiosyncratic location-specific component $\iota_\ell$.
Under our return decomposition
$\log R_\ell = \mu_h + \eta_{\text{div}} + \iota_\ell$
(Section~\ref{sec:model:returns}), a REIT-like diversified asset delivers
$E[\eta_{\text{div}}]$ but $E[\iota_\ell] = 0$ in large diversified portfolios.

Location-$A$ tokens, by contrast, deliver $\iota_A$~---~the idiosyncratic
appreciation specific to the MSA the household has personal ties to.
For a household who lived in location~$A$, built human capital there, and
may return, the idiosyncratic component $\iota_A$ is exactly the exposure
they wish to retain.
No REIT can supply this:
the location-specific fractional token is the only instrument that provides
$\iota_A$-loading for a household currently at~$B$.

This is why the 2026-05-01 pivot from the v2 four-regime (REIT-comparison)
framework to the v3/v4 two-location framework was necessary: the
structurally novel contribution of tokenization is not relative to REIT access
but relative to \emph{direct ownership with forced relocation sale}.

% ─────────────────────────────────────────────────────────────────────────────
\subsection{Relation to Other Lifecycle Housing Models}
\label{sec:discussion:literature}
% ─────────────────────────────────────────────────────────────────────────────

Table~\ref{tab:literature} summarizes the placement of our contribution
relative to the closest precedents.

\begin{table}[t]
\centering
\caption{Relationship to Prior Lifecycle Housing-Portfolio Models}
\label{tab:literature}
\small
\begin{tabular}{@{}lllll@{}}
\toprule
Paper & Relocation & Tx costs & Continuous $x$ & Token portability \\
\midrule
\citet{YaoZhang2005}     & No  & No  & No & No \\
\citet{Cocco2005}         & No  & No  & No & No \\
\citet{KraftMunk2011}     & No  & No  & No & No \\
\citet{KMW2018}           & Yes (reduced-form) & No  & No & No \\
\citet{Liu2021}           & No  & No  & Yes & No \\
\textbf{This paper (E1\_2L)} & Yes & Yes & No  & No \\
\textbf{This paper (E2\_2L)} & Yes & Yes (on $\Delta x$) & Yes & \textbf{Yes} \\
\bottomrule
\end{tabular}
\begin{tablenotes}
\footnotesize
\item \textit{Notes.}
``Relocation'' indicates whether the model includes a stochastic
geographic relocation shock.
``Tx costs'' indicates whether transaction costs (selling and/or
buying commissions) are modelled.
``Continuous $x$'' indicates whether fractional ownership is permitted.
``Token portability'' indicates whether the fractional position can be
retained across a geographic relocation without a forced sale.
\end{tablenotes}
\end{table}

\citet{KMW2018} model relocation explicitly in a lifecycle framework with
habit formation, but transaction costs are absent and the continuous-$x$
channel is not modelled.
Their focus is the housing-as-habit hedge against future consumption
commitments, a mechanism orthogonal to our pre-buying channel.

\citet{SinaiSouleles2005} study the rent-price risk hedge embedded in
homeownership: owning provides insurance against future rent increases for
households with long planned tenures.
Their mechanism is complementary to ours~---~the rent-hedge motive
encourages ownership even when the financial return to housing is below
the equity premium, exactly as in our model~---~but they abstract from
relocation and portability.

\citet{Davidoff2006} studies the income-housing correlation (``double
exposure'' risk) and its interaction with optimal tenure.
This correlation is relevant to our model through the $(\iota_A, \text{income})$
channel; we abstract from it at baseline by maintaining independent income
and housing shocks, and defer an income-correlated location shock extension
to future work (see Section~\ref{sec:discussion:limitations}).

% ─────────────────────────────────────────────────────────────────────────────
\subsection{Robustness of the Mechanism}
\label{sec:discussion:robustness}
% ─────────────────────────────────────────────────────────────────────────────

The pre-registered falsification structure (Section~\ref{sec:results:falsification})
provides the clearest evidence on the robustness of each channel.

\paragraph{Test (r): Relocation disabled ($p_{\text{reloc}} = 0$).}
With no relocation shocks, the motivation to hold $x_B$ while at~$A$
vanishes~---~the household will never visit location~$B$, so no
pre-buying saving accrues.
Under v4~Option~1, the pre-buying hedge channel must therefore collapse
to zero, while the continuous-$x$ channel persists (it does not require
relocation).
Concretely, the expected result is $\bar{x}_B \to 0$ and
$\operatorname{CEV}(E2\_2L \text{ vs } E1\_2L) \to
\operatorname{CEV}_{\text{Liu}}$ (continuous-$x$ only).
[\textsc{Placeholder: report $\Delta\text{CEV}$ and $\bar{x}_B$ from sweep.}]

\paragraph{Test (m): High location correlation ($\rho_{AB} = 0.95$).}
As $\rho_{AB} \to 1$, the location-$B$ token delivers almost the same
return as the location-$A$ token; the diversification value of the
cross-location hedge diminishes.
The pre-buying saving is unchanged in magnitude ($p_{\text{reloc}} \cdot
\tau_{\text{buy}}$ is independent of $\rho_{AB}$), but households who
value the idiosyncratic hedge component will demand less $x_B$ as
$\rho_{AB}$ rises.
We therefore expect the hedge channel to attenuate~---~but not to
collapse completely~---~as $\rho_{AB} \to 0.95$.
[\textsc{Placeholder: report $\bar{x}_B(\rho_{AB} = 0.95)$ from sweep.}]

\paragraph{Test (q): No buying cost ($\tau_{\text{buy}} = 0$).}
With $\tau_{\text{buy}} = 0$, pre-buying $x_B$ saves nothing; the
pre-buying motive is eliminated.
The expected outcome is $\bar{x}_B = 0$ (identical to v3~Option~3 under
zero buying cost) and a CEV change equal only to the sell-side avoidance
and continuous-$x$ channels.
This isolates the pre-buying contribution to the total welfare gain.
[\textsc{Placeholder: report from sweep\_txcost.sh run.}]

\paragraph{Sensitivity to $p_{\text{reloc}}$ and $\rho_{AB}$ jointly.}
Figure~\ref{fig:heatmap_cev} (placeholder) reports the sensitivity heatmap
over $(\rho_{AB}, p_{\text{reloc}}) \in \{0, 0.25, 0.50, 0.75, 0.95\} \times
\{0, 0.02, 0.06, 0.12\}$.
The model predicts a weakly increasing CEV in $p_{\text{reloc}}$ (more
mobility amplifies pre-buying value) and a weakly decreasing CEV in
$\rho_{AB}$ (higher correlation reduces the idiosyncratic hedge value).
If the empirical pattern deviates from these monotone predictions,
the mechanism claim requires reassessment.
[\textsc{Placeholder: fill from scripts/sweep\_rhoAB.sh and
scripts/sweep\_prelocate.sh.}]

% ─────────────────────────────────────────────────────────────────────────────
\subsection{Partial-Equilibrium Scope and Limitations}
\label{sec:discussion:limitations}
% ─────────────────────────────────────────────────────────────────────────────

All welfare estimates in this paper are partial-equilibrium in three respects.

\paragraph{Exogenous house prices.}
We take the location-specific house-price processes $(R_A, R_B)$ as
exogenous, calibrated to observed MSA-level moments.
Widespread adoption of token portability could affect housing demand in
origin and destination locations, potentially altering $R_A$ and $R_B$.
A general-equilibrium model with endogenous prices is left to future work.
Our CEV estimates should be interpreted as the household-level welfare value
of \emph{access} to tokenized instruments, conditional on other households'
behavior remaining unchanged.

\paragraph{Symmetric and stationary location shocks.}
The relocation shock is i.i.d.\ Bernoulli with constant probability
$p_{\text{reloc}}(t)$.
In practice, relocation rates are correlated with employment shocks,
income shocks, and housing-market conditions (\citealt{PSID}).
An extension with income-correlated location shocks (e.g.,
$\text{corr}(\varepsilon_t, \iota_{A,t}) > 0$) would add a second
channel to the pre-buying demand~---~holding $x_B$ at~$A$ also hedges
against a negative income draw at~$A$ that triggers relocation to~$B$.
This extension is queued as a robustness check but is not included in
the baseline results.

\paragraph{Two-location, symmetric-returns simplification.}
We model two symmetric locations with exchangeable house-price processes
(same $\mu_h$, $\sigma_h$; differing only by idiosyncratic component).
Empirically, households relocate across heterogeneous MSAs with different
mean appreciation rates and rent-to-price ratios.
Asymmetric robustness (different $\mu_A \neq \mu_B$, different
$p_{A \to B} \neq p_{B \to A}$) is queued in Phase~2 sensitivity work
(Table~\ref{tab:sensitivity}, rows 6--7).

\paragraph{No mortgage.}
The baseline results set $\text{LTV\_MAX} = 0$.
Mortgages reduce the effective down payment on housing, raising the
household's ability to reach the optimal $x_A$ at low wealth levels.
This can attenuate the pre-buying demand for $x_B$ (see
Section~\ref{sec:discussion:hedge}) if households who were
wealth-constrained in $x_A$ can now fully fund their desired $x_A$
via borrowing.
Mortgage activation (LTV $\in \{0.5, 0.8\}$) is queued as a Phase~2
sensitivity check.
In the v2 baseline (\citealt{research_log_2026_05_01_mortgage}),
mortgage reduced $\operatorname{CEV}(\text{E2\_2L vs E1\_2L})$ by
approximately 37\% under no-relocation; the direction under the
2-location model may differ.

\paragraph{Platform and regulatory risk.}
We abstract from platform failure, regulatory uncertainty, and
information asymmetries between token issuers and holders.
These factors introduce a risk premium on tokenized positions that
is not captured in the $\tau_{\text{token}} = 1\%$ transfer cost
alone.
A fuller treatment of platform risk is deferred to companion work.

% ─────────────────────────────────────────────────────────────────────────────
\subsection{Policy Implications}
\label{sec:discussion:policy}
% ─────────────────────────────────────────────────────────────────────────────

Our findings have three direct policy implications for household finance
and housing-market regulation.

\paragraph{Mobility lock-in as a welfare cost.}
The $+0.82\%$ transaction-cost-avoidance channel in the v3~Option~3 baseline
quantifies the welfare cost of \emph{mobility lock-in}: households who own
traditional homes are effectively penalized for relocating, reducing their
geographic flexibility.
Token portability eliminates this distortion.
Policy interventions that reduce $\tau_{\text{sell}}$ or $\tau_{\text{buy}}$
(e.g., real-estate transfer tax reform) would deliver similar gains without
requiring tokenization, but would not restore the idiosyncratic hedge or the
pre-buying channel.

\paragraph{Token design and the buying-cost schedule.}
If $\tau_{\text{buy}}$ is too high, the pre-buying hedge becomes economically
negligible ($p_{\text{reloc}} \cdot \tau_{\text{buy}}$ too small relative to
the $\delta_{\text{own}}$ opportunity cost).
If $\tau_{\text{buy}}$ is too low (approaching $\tau_{\text{token}} = 1\%$),
the pre-buying saving is small and the mechanism collapses.
The sensitivity analysis over $\tau_{\text{buy}} \in \{0, 1\%, 2.5\%, 4\%, 6\%\}$
identifies the buying-cost levels at which the pre-buying channel is most
economically significant.
This informs platform pricing: the welfare case for token portability relative
to direct ownership is maximized at intermediate buying cost levels.

\paragraph{Household heterogeneity.}
The pre-buying mechanism is most valuable for households who are
(i)~frequently mobile ($p_{\text{reloc}}$ high),
(ii)~have location-specific ties to prior MSAs ($\iota_A$ relevant to their
     welfare), and
(iii)~are wealth-constrained in their $x_A$ position (forcing $\Delta x_A = 0$
      and making $x_B$ investment relatively cheap).
From a regulatory standpoint, this suggests that token portability delivers
disproportionate welfare gains to mobile younger households with limited
wealth~---~a population underserved by existing housing finance instruments.

% ─────────────────────────────────────────────────────────────────────────────

\clearpage

% ─────────────────────────────────────────────────────────────────────────────
% s6_conclusion.tex — Conclusion
% Paper: "Tokenized Housing and Lifecycle Portfolio Choice:
%         A Decoupling of Location from Housing Exposure"
%
% Source: paper/outline_v4.md §6 (3–4 paragraph structure)
%         docs/welfare_decomp_v4.md §4 (channel definitions)
%         paper/sections/s5_discussion.tex (Liu comparison, policy)
%
% Status: skeleton complete; [PLACEHOLDER] marks require server1 numerical results.
%         Fill in: CEV headline, channel magnitudes, mean_xB evidence.
%
% Compilation: % s6_conclusion.tex — Conclusion
% Paper: "Tokenized Housing and Lifecycle Portfolio Choice:
%         A Decoupling of Location from Housing Exposure"
%
% Source: paper/outline_v4.md §6 (3–4 paragraph structure)
%         docs/welfare_decomp_v4.md §4 (channel definitions)
%         paper/sections/s5_discussion.tex (Liu comparison, policy)
%
% Status: skeleton complete; [PLACEHOLDER] marks require server1 numerical results.
%         Fill in: CEV headline, channel magnitudes, mean_xB evidence.
%
% Compilation: % s6_conclusion.tex — Conclusion
% Paper: "Tokenized Housing and Lifecycle Portfolio Choice:
%         A Decoupling of Location from Housing Exposure"
%
% Source: paper/outline_v4.md §6 (3–4 paragraph structure)
%         docs/welfare_decomp_v4.md §4 (channel definitions)
%         paper/sections/s5_discussion.tex (Liu comparison, policy)
%
% Status: skeleton complete; [PLACEHOLDER] marks require server1 numerical results.
%         Fill in: CEV headline, channel magnitudes, mean_xB evidence.
%
% Compilation: \input{sections/s6_conclusion} from main.tex

% ─────────────────────────────────────────────────────────────────────────────
\section{Conclusion}
\label{sec:conclusion}
% ─────────────────────────────────────────────────────────────────────────────

This paper quantifies the lifetime welfare value of a structural capability
that distinguishes residential housing tokens from all prior instruments in
the household's choice set: \emph{the decoupling of geographic location from
housing-asset exposure}.
In a two-location lifecycle model with stochastic relocation shocks, we
compare three regimes --- rent-only (E0), traditional binary homeownership
(E1\textsubscript{2L}), and continuous fractional token ownership portable
across relocations (E2\textsubscript{2L}) --- and solve for the welfare
difference using a 6-dimensional backward-induction algorithm that tracks
prior-period token holdings as explicit state variables.
Our headline result is $\text{CEV}(\text{E2}_\text{2L}\ \text{vs}\
\text{E1}_\text{2L}) = \textsc{[P]}\%$ at baseline calibration
($\gamma = 5$, $p_{\text{reloc}} = 6\%$, $\tau_{\text{buy}} = 2.5\%$,
$\rho_{AB} = 0.50$).

\paragraph{What the welfare gain decomposes into.}
Following the three-channel decomposition of Section~\ref{sec:results:decomp}
and the formal claim in equation~\eqref{eq:decomp_claim}, the total gain
separates into two structurally distinct components.
First, an \emph{avoided-transaction-cost channel} of \textsc{[P]}\% arises
because tokenized households escape the $\tau_{\text{sell}} + \tau_{\text{buy}}
\approx 8.5\%$ round-trip cost that E1\textsubscript{2L} households pay at
each relocation event.
This channel is present in any portable financial instrument that avoids
forced sale at moving and delivers the same CEV gain regardless of
location-return correlations.
Second, a \emph{pre-buying hedge channel} of \textsc{[P]}\% emerges
specifically from the 6-dimensional state extension: a household at location~A
can incrementally accumulate location-B tokens at cost $\tau_{\text{buy}}$
per unit per period, reducing the lump-sum buying cost on future relocation
to B.
This channel is \emph{uniquely activated by the state extension} (Option~1):
it is absent in the Option~3 approximation and in any model that does not
track prior-period holdings.
The falsification tests of Section~\ref{sec:results:falsification} confirm the
mechanism: the hedge channel collapses when $p_{\text{reloc}} = 0$ (no
relocation motive) and attenuates monotonically as $\rho_{AB} \to 1$
(cross-location exposure becomes redundant).
The cross-term $\xi = \textsc{[P]}\%$ (equation~\eqref{eq:decomp}) confirms
near-additive separability of the two channels.

\paragraph{How this contributes beyond existing work.}
The closest benchmark is \citet{Liu2021}, who shows that relaxing binary
homeownership to continuous fractional ownership delivers 5--10\% welfare
gains under a single-location model without relocation.
Our avoided-transaction-cost channel (\textsc{[P]}\%) is partially Liu-adjacent:
it arises because tokens reduce the lump-sum trade in of E1 at relocation.
However, Liu's setup has no relocation, and hence this channel is outside his
framework.
The pre-buying hedge channel (\textsc{[P]}\%) is structurally novel: it
requires (i)~two locations, (ii)~stochastic relocation, (iii)~positive
$\tau_{\text{buy}}$, and (iv)~state-tracked prior holdings --- none of which
appear in \citet{Liu2021}, \citet{YaoZhang2005}, \citet{Cocco2005}, or
\citet{KMW2018}.
Table~\ref{tab:literature} in Section~\ref{sec:discussion:literature}
formalizes this comparison.
Our results also add a quantitative claim to the qualitative portability
arguments of \citet{SinaiSouleles2005}: token portability is worth \textsc{[P]}\%
in lifetime CEV, with the hedge component growing in mobility rate and
shrinking in cross-location correlation.

\paragraph{Policy implications and token design.}
The welfare case for token portability is strongest for mobile, wealth-constrained
younger households (Section~\ref{sec:discussion:policy}).
This has a direct implication for regulatory design: mandatory secondary-market
liquidity and low $\tau_{\text{token}}$ are important for the avoided-tx
channel, while a moderate $\tau_{\text{buy}}$ (roughly 1--4\%) maximises the
pre-buying hedge (the benefit-per-unit is $p_{\text{reloc}} \times
\tau_{\text{buy}} \approx 0.15\%$/period at baseline, attaining maximum
welfare contribution at intermediate values of $\tau_{\text{buy}}$).
Policy interventions that reduce $\tau_{\text{sell}}$ or $\tau_{\text{buy}}$
via transfer-tax reform deliver similar avoided-tx gains without requiring
tokenization, but cannot restore the idiosyncratic location-specific hedge.
The two channels therefore have different policy substitutes, which sharpens
the design case for token portability specifically.

\paragraph{Limitations and future work.}
This analysis is partial equilibrium: house prices are exogenous, and we
do not model the general-equilibrium price impact of widespread tokenization
adoption.
The model abstracts from heterogeneous preferences, household composition
changes (marriage, children), and tax treatment (property taxes, capital
gains deferral) that may interact with relocation incentives.
On the numerical side, the coarse x-prev grid ($N_{x,\text{prev}} = 3$) is
a pragmatic first step; finer grids would sharpen the pre-buying hedge
estimates at the cost of substantially more compute.
The sensitivity panel (Table~\ref{tab:sensitivity}) documents that the
headline CEV is robust across plausible parameter ranges, but robustness to
structural assumptions --- homogeneous agents, symmetric locations,
no mortgage interaction with token holdings --- remains to be established.
Three natural extensions are: (i)~a general-equilibrium version that endogenizes
house prices to study the price-impact of tokenization adoption at scale;
(ii)~heterogeneous-agent lifecycle models with idiosyncratic labor-market shocks
correlated with local housing returns (building on \citealt{BFN2014});
and (iii)~an explicit treatment of tax sheltering and depreciation schedules
that would affect the effective $\tau_{\text{buy}}$ and $\tau_{\text{token}}$
in practice.
The companion paper explores the case where the token portfolio spans multiple
distinct locations simultaneously, examining the optimal diversification
across a richer location menu.
 from main.tex

% ─────────────────────────────────────────────────────────────────────────────
\section{Conclusion}
\label{sec:conclusion}
% ─────────────────────────────────────────────────────────────────────────────

This paper quantifies the lifetime welfare value of a structural capability
that distinguishes residential housing tokens from all prior instruments in
the household's choice set: \emph{the decoupling of geographic location from
housing-asset exposure}.
In a two-location lifecycle model with stochastic relocation shocks, we
compare three regimes --- rent-only (E0), traditional binary homeownership
(E1\textsubscript{2L}), and continuous fractional token ownership portable
across relocations (E2\textsubscript{2L}) --- and solve for the welfare
difference using a 6-dimensional backward-induction algorithm that tracks
prior-period token holdings as explicit state variables.
Our headline result is $\text{CEV}(\text{E2}_\text{2L}\ \text{vs}\
\text{E1}_\text{2L}) = \textsc{[P]}\%$ at baseline calibration
($\gamma = 5$, $p_{\text{reloc}} = 6\%$, $\tau_{\text{buy}} = 2.5\%$,
$\rho_{AB} = 0.50$).

\paragraph{What the welfare gain decomposes into.}
Following the three-channel decomposition of Section~\ref{sec:results:decomp}
and the formal claim in equation~\eqref{eq:decomp_claim}, the total gain
separates into two structurally distinct components.
First, an \emph{avoided-transaction-cost channel} of \textsc{[P]}\% arises
because tokenized households escape the $\tau_{\text{sell}} + \tau_{\text{buy}}
\approx 8.5\%$ round-trip cost that E1\textsubscript{2L} households pay at
each relocation event.
This channel is present in any portable financial instrument that avoids
forced sale at moving and delivers the same CEV gain regardless of
location-return correlations.
Second, a \emph{pre-buying hedge channel} of \textsc{[P]}\% emerges
specifically from the 6-dimensional state extension: a household at location~A
can incrementally accumulate location-B tokens at cost $\tau_{\text{buy}}$
per unit per period, reducing the lump-sum buying cost on future relocation
to B.
This channel is \emph{uniquely activated by the state extension} (Option~1):
it is absent in the Option~3 approximation and in any model that does not
track prior-period holdings.
The falsification tests of Section~\ref{sec:results:falsification} confirm the
mechanism: the hedge channel collapses when $p_{\text{reloc}} = 0$ (no
relocation motive) and attenuates monotonically as $\rho_{AB} \to 1$
(cross-location exposure becomes redundant).
The cross-term $\xi = \textsc{[P]}\%$ (equation~\eqref{eq:decomp}) confirms
near-additive separability of the two channels.

\paragraph{How this contributes beyond existing work.}
The closest benchmark is \citet{Liu2021}, who shows that relaxing binary
homeownership to continuous fractional ownership delivers 5--10\% welfare
gains under a single-location model without relocation.
Our avoided-transaction-cost channel (\textsc{[P]}\%) is partially Liu-adjacent:
it arises because tokens reduce the lump-sum trade in of E1 at relocation.
However, Liu's setup has no relocation, and hence this channel is outside his
framework.
The pre-buying hedge channel (\textsc{[P]}\%) is structurally novel: it
requires (i)~two locations, (ii)~stochastic relocation, (iii)~positive
$\tau_{\text{buy}}$, and (iv)~state-tracked prior holdings --- none of which
appear in \citet{Liu2021}, \citet{YaoZhang2005}, \citet{Cocco2005}, or
\citet{KMW2018}.
Table~\ref{tab:literature} in Section~\ref{sec:discussion:literature}
formalizes this comparison.
Our results also add a quantitative claim to the qualitative portability
arguments of \citet{SinaiSouleles2005}: token portability is worth \textsc{[P]}\%
in lifetime CEV, with the hedge component growing in mobility rate and
shrinking in cross-location correlation.

\paragraph{Policy implications and token design.}
The welfare case for token portability is strongest for mobile, wealth-constrained
younger households (Section~\ref{sec:discussion:policy}).
This has a direct implication for regulatory design: mandatory secondary-market
liquidity and low $\tau_{\text{token}}$ are important for the avoided-tx
channel, while a moderate $\tau_{\text{buy}}$ (roughly 1--4\%) maximises the
pre-buying hedge (the benefit-per-unit is $p_{\text{reloc}} \times
\tau_{\text{buy}} \approx 0.15\%$/period at baseline, attaining maximum
welfare contribution at intermediate values of $\tau_{\text{buy}}$).
Policy interventions that reduce $\tau_{\text{sell}}$ or $\tau_{\text{buy}}$
via transfer-tax reform deliver similar avoided-tx gains without requiring
tokenization, but cannot restore the idiosyncratic location-specific hedge.
The two channels therefore have different policy substitutes, which sharpens
the design case for token portability specifically.

\paragraph{Limitations and future work.}
This analysis is partial equilibrium: house prices are exogenous, and we
do not model the general-equilibrium price impact of widespread tokenization
adoption.
The model abstracts from heterogeneous preferences, household composition
changes (marriage, children), and tax treatment (property taxes, capital
gains deferral) that may interact with relocation incentives.
On the numerical side, the coarse x-prev grid ($N_{x,\text{prev}} = 3$) is
a pragmatic first step; finer grids would sharpen the pre-buying hedge
estimates at the cost of substantially more compute.
The sensitivity panel (Table~\ref{tab:sensitivity}) documents that the
headline CEV is robust across plausible parameter ranges, but robustness to
structural assumptions --- homogeneous agents, symmetric locations,
no mortgage interaction with token holdings --- remains to be established.
Three natural extensions are: (i)~a general-equilibrium version that endogenizes
house prices to study the price-impact of tokenization adoption at scale;
(ii)~heterogeneous-agent lifecycle models with idiosyncratic labor-market shocks
correlated with local housing returns (building on \citealt{BFN2014});
and (iii)~an explicit treatment of tax sheltering and depreciation schedules
that would affect the effective $\tau_{\text{buy}}$ and $\tau_{\text{token}}$
in practice.
The companion paper explores the case where the token portfolio spans multiple
distinct locations simultaneously, examining the optimal diversification
across a richer location menu.
 from main.tex

% ─────────────────────────────────────────────────────────────────────────────
\section{Conclusion}
\label{sec:conclusion}
% ─────────────────────────────────────────────────────────────────────────────

This paper quantifies the lifetime welfare value of a structural capability
that distinguishes residential housing tokens from all prior instruments in
the household's choice set: \emph{the decoupling of geographic location from
housing-asset exposure}.
In a two-location lifecycle model with stochastic relocation shocks, we
compare three regimes --- rent-only (E0), traditional binary homeownership
(E1\textsubscript{2L}), and continuous fractional token ownership portable
across relocations (E2\textsubscript{2L}) --- and solve for the welfare
difference using a 6-dimensional backward-induction algorithm that tracks
prior-period token holdings as explicit state variables.
Our headline result is $\text{CEV}(\text{E2}_\text{2L}\ \text{vs}\
\text{E1}_\text{2L}) = \textsc{[P]}\%$ at baseline calibration
($\gamma = 5$, $p_{\text{reloc}} = 6\%$, $\tau_{\text{buy}} = 2.5\%$,
$\rho_{AB} = 0.50$).

\paragraph{What the welfare gain decomposes into.}
Following the three-channel decomposition of Section~\ref{sec:results:decomp}
and the formal claim in equation~\eqref{eq:decomp_claim}, the total gain
separates into two structurally distinct components.
First, an \emph{avoided-transaction-cost channel} of \textsc{[P]}\% arises
because tokenized households escape the $\tau_{\text{sell}} + \tau_{\text{buy}}
\approx 8.5\%$ round-trip cost that E1\textsubscript{2L} households pay at
each relocation event.
This channel is present in any portable financial instrument that avoids
forced sale at moving and delivers the same CEV gain regardless of
location-return correlations.
Second, a \emph{pre-buying hedge channel} of \textsc{[P]}\% emerges
specifically from the 6-dimensional state extension: a household at location~A
can incrementally accumulate location-B tokens at cost $\tau_{\text{buy}}$
per unit per period, reducing the lump-sum buying cost on future relocation
to B.
This channel is \emph{uniquely activated by the state extension} (Option~1):
it is absent in the Option~3 approximation and in any model that does not
track prior-period holdings.
The falsification tests of Section~\ref{sec:results:falsification} confirm the
mechanism: the hedge channel collapses when $p_{\text{reloc}} = 0$ (no
relocation motive) and attenuates monotonically as $\rho_{AB} \to 1$
(cross-location exposure becomes redundant).
The cross-term $\xi = \textsc{[P]}\%$ (equation~\eqref{eq:decomp}) confirms
near-additive separability of the two channels.

\paragraph{How this contributes beyond existing work.}
The closest benchmark is \citet{Liu2021}, who shows that relaxing binary
homeownership to continuous fractional ownership delivers 5--10\% welfare
gains under a single-location model without relocation.
Our avoided-transaction-cost channel (\textsc{[P]}\%) is partially Liu-adjacent:
it arises because tokens reduce the lump-sum trade in of E1 at relocation.
However, Liu's setup has no relocation, and hence this channel is outside his
framework.
The pre-buying hedge channel (\textsc{[P]}\%) is structurally novel: it
requires (i)~two locations, (ii)~stochastic relocation, (iii)~positive
$\tau_{\text{buy}}$, and (iv)~state-tracked prior holdings --- none of which
appear in \citet{Liu2021}, \citet{YaoZhang2005}, \citet{Cocco2005}, or
\citet{KMW2018}.
Table~\ref{tab:literature} in Section~\ref{sec:discussion:literature}
formalizes this comparison.
Our results also add a quantitative claim to the qualitative portability
arguments of \citet{SinaiSouleles2005}: token portability is worth \textsc{[P]}\%
in lifetime CEV, with the hedge component growing in mobility rate and
shrinking in cross-location correlation.

\paragraph{Policy implications and token design.}
The welfare case for token portability is strongest for mobile, wealth-constrained
younger households (Section~\ref{sec:discussion:policy}).
This has a direct implication for regulatory design: mandatory secondary-market
liquidity and low $\tau_{\text{token}}$ are important for the avoided-tx
channel, while a moderate $\tau_{\text{buy}}$ (roughly 1--4\%) maximises the
pre-buying hedge (the benefit-per-unit is $p_{\text{reloc}} \times
\tau_{\text{buy}} \approx 0.15\%$/period at baseline, attaining maximum
welfare contribution at intermediate values of $\tau_{\text{buy}}$).
Policy interventions that reduce $\tau_{\text{sell}}$ or $\tau_{\text{buy}}$
via transfer-tax reform deliver similar avoided-tx gains without requiring
tokenization, but cannot restore the idiosyncratic location-specific hedge.
The two channels therefore have different policy substitutes, which sharpens
the design case for token portability specifically.

\paragraph{Limitations and future work.}
This analysis is partial equilibrium: house prices are exogenous, and we
do not model the general-equilibrium price impact of widespread tokenization
adoption.
The model abstracts from heterogeneous preferences, household composition
changes (marriage, children), and tax treatment (property taxes, capital
gains deferral) that may interact with relocation incentives.
On the numerical side, the coarse x-prev grid ($N_{x,\text{prev}} = 3$) is
a pragmatic first step; finer grids would sharpen the pre-buying hedge
estimates at the cost of substantially more compute.
The sensitivity panel (Table~\ref{tab:sensitivity}) documents that the
headline CEV is robust across plausible parameter ranges, but robustness to
structural assumptions --- homogeneous agents, symmetric locations,
no mortgage interaction with token holdings --- remains to be established.
Three natural extensions are: (i)~a general-equilibrium version that endogenizes
house prices to study the price-impact of tokenization adoption at scale;
(ii)~heterogeneous-agent lifecycle models with idiosyncratic labor-market shocks
correlated with local housing returns (building on \citealt{BFN2014});
and (iii)~an explicit treatment of tax sheltering and depreciation schedules
that would affect the effective $\tau_{\text{buy}}$ and $\tau_{\text{token}}$
in practice.
The companion paper explores the case where the token portfolio spans multiple
distinct locations simultaneously, examining the optimal diversification
across a richer location menu.

% ─────────────────────────────────────────────────────────────────────────────

\clearpage

% ─── References ──────────────────────────────────────────────────────────────
\bibliography{references}

% ─────────────────────────────────────────────────────────────────────────────
\clearpage
\begin{appendices}
% ─────────────────────────────────────────────────────────────────────────────

\section{Numerical Implementation Details}
\label{app:numerical}

\subsection{Value function iteration}

The Bellman equation is solved by backward induction over periods $t = T, T-1,
\ldots, 1$.
Each period, we iterate over the 6-dimensional state grid $(w, z, \ell,
x_{A,\text{prev}}, x_{B,\text{prev}})$ with $T = 56$ periods (ages 25--80).
The continuation value is approximated by 7-dimensional Gauss-Hermite
quadrature with $n = 3$ nodes per dimension ($3^7 = 2{,}187$ quadrature points),
integrating jointly over the stock return, aggregate housing factor, two
idiosyncratic location shocks, house-price normalisation, permanent income shock,
and transitory income shock.

\subsection{Interpolation}

Continuation values are computed by bilinear interpolation over the $(w, z)$
grid dimensions.
The $(x_{A,\text{prev}}, x_{B,\text{prev}})$ dimensions are indexed exactly:
by restricting the choice of $x_{\text{new}}$ to the $x_{\text{prev}}$ grid,
the next-period state is always an exact grid point, eliminating the need to
interpolate over the $x_{\text{prev}}$ dimensions.

\subsection{Coarse grid and computation time}

Baseline grids: $N_W = 15$, $N_Z = 5$, $N_{x,\text{prev}} = 3$,
$N_{\ell} = 2$, giving $15 \times 5 \times 2 \times 3 \times 3 = 1{,}350$
states per period.
Control search: $N_{\text{asset}} = 7$ points for $b$ and $s$ candidates;
$N_{x,\text{prev}} = 3$ choices for each of $x_{A,\text{new}}$ and
$x_{B,\text{new}}$ (9 combinations in E2\textsubscript{2L}).
Estimated wall time per regime at baseline: \ph{X} hours on [server1 spec].
Full-grid runs use $N_W = 40$, $N_Z = 9$, $N_{x,\text{prev}} = 5$ and are
used for robustness checks.

\section{Convergence Diagnostics}
\label{app:convergence}

\ph{To be populated from output/diagnostics/p6\_option1\_smoke.md after server1 run.}

Table~\ref{tab:convergence} reports Euler-residual statistics and
value-function sup-norm convergence at the coarse and full grids.

\begin{table}[h]
\centering
\caption{Convergence diagnostics \ph{(placeholder — fill from server1 output)}}
\label{tab:convergence}
\begin{tabular}{lcc}
\toprule
Metric & Coarse grid & Full grid \\
\midrule
Feasible states at $t=1$ (\%) & \ph{X} & \ph{X} \\
Euler residual p50             & \ph{X} & \ph{X} \\
Euler residual p95             & \ph{X} & \ph{X} \\
NaN / Inf health check         & PASS   & \ph{TBD} \\
Terminal-identity check        & PASS   & \ph{TBD} \\
\bottomrule
\end{tabular}
\end{table}

\section{Hedge-Channel Sign Conditions}
\label{app:sign}

\begin{proposition}[Necessary condition for a positive pre-buying hedge channel]
\label{prop:sign}
In regime \EtwosubTwoL, a household at location~$A$ holds $x_{B,\text{new}} > 0$
in equilibrium if and only if the expected present value of future buying-cost
savings exceeds the current opportunity cost of holding a non-occupied token:
\begin{equation}
  \underbrace{p_{\text{reloc}} \cdot \taubuy}_{\text{expected saving per unit}}
  >
  \underbrace{\delta_{\text{own}} - \mathbb{E}[R_B - R_f]}_{\text{opportunity cost of }x_B \text{ over bond}},
  \label{eq:sign_condition}
\end{equation}
where $\delta_{\text{own}} = \rho - m$ is the rent-saving forgone by holding a
non-occupied unit, and $\mathbb{E}[R_B - R_f]$ is the expected excess return
on location-B housing.
\end{proposition}

\begin{remark}
At baseline calibration: $p_{\text{reloc}} \cdot \taubuy = 0.06 \times 0.025 =
0.0015$ per period per unit; $\delta_{\text{own}} = 0.04$;
$\mathbb{E}[R_B - R_f] = \ph{X}$.
The sign of the hedge channel is therefore an empirical question, resolved by
the server1 baseline runs.
If the condition holds, $x_{B,\text{new}} > 0$ at ell=A (H1 confirmed);
if it fails, the channel is inactive (fall back to Path~D).
\end{remark}

\end{appendices}

% ─────────────────────────────────────────────────────────────────────────────
\end{document}
% ─────────────────────────────────────────────────────────────────────────────
